\documentclass[UTF8,12pt,a4paper,fontset=fandol]{ctexbook}

% 支持从项目根目录或当前笔记目录编译。
\makeatletter
\def\input@path{{../}{../../}}
\makeatother
\usepackage{researchnotes}

\title{泛函分析笔记}
\author{}
\date{}

\begin{document}

\maketitle
\tableofcontents

\chapter{绪论：从线性代数到泛函分析}
\label{chap:fa-introduction}

\section{研究对象与基本问题}

\keyterm{泛函分析（functional analysis）}研究带有拓扑、范数或内积结构的线性空间，以及这些空间上的函数与算子。它将有限维线性代数中的向量、矩阵和线性方程推广到函数空间，使微分方程、积分方程、函数逼近与变分问题能够在共同的框架中讨论。

全书以三个问题组织内容：首先，如何选择容纳未知函数的空间，并保证极限仍留在该空间中；其次，怎样判断一个方程是否有解、解是否唯一，以及解对数据是否连续依赖；最后，当有限维的紧性与对角化方法失效时，如何借助弱收敛、紧算子和谱理论恢复可用的结构。

学习这些问题时，可以把函数暂时看成有无穷多个坐标的向量：数列的坐标用自然数编号，函数的取值用自变量编号。但这种类比只解释线性运算，不能替代分析。有限维中逐坐标收敛就是向量收敛；无限维中，每个坐标都变小，全部坐标合起来的误差却可能不变。因而本书始终同时追问两个问题：对象之间有什么代数关系，以及所用极限究竟在哪一种意义下成立。

\section{从一个积分方程开始}
考虑未知连续函数 $u:[0,1]\to\mathbb R$ 满足
\begin{equation}
 u(t)=t+\frac12\int_0^t u(s)\,\mathrm ds.
 \label{eq:fa-intro-integral}
\end{equation}
如果已经知道 $u$ 连续，右端可微，就能求导得到 $u'=1+u/2$、$u(0)=0$，从而猜到 $u(t)=2(e^{t/2}-1)$。不过，更一般的积分方程未必能这样求导，微分后的方程也未必有显式解。泛函分析提供的办法是直接研究整个函数所组成的空间。

从 $u_0(t)=0$ 出发，把已知函数代入右端，依次得到
\[
 u_1(t)=t,\qquad u_2(t)=t+\frac{t^2}{4},\qquad
 u_3(t)=t+\frac{t^2}{4}+\frac{t^3}{24}.
\]
记右端变换为 $F$，用 $\|u\|_\infty=\max_{0\leq t\leq1}|u(t)|$ 测量函数大小，则
\[
 \|F(u)-F(v)\|_\infty
 \leq\frac12\|u-v\|_\infty.
\]
也就是说，每次代入至少把两条输入曲线的最大差缩小一半。因此相邻迭代的误差形成可求和的几何级数，$(u_n)$ 是 Cauchy 序列。连续函数空间在这一距离下完备，保证极限仍为连续函数；$F$ 连续，保证极限满足原方程；同一个收缩估计还保证两个解只能相同。这正是第\ref{chap:fa-metric}章压缩映射原理的具体模型。

这个例子中，空间回答“解住在哪里”，范数回答“怎样测量误差”，算子回答“方程怎样变换函数”，完备性回答“近似过程会不会跑出空间”。后面处理最佳逼近时，用投影代替迭代；处理能量极小化时，用弱收敛寻找极限；处理积分方程时，用谱判断哪些参数会破坏可解性。方法虽然不同，选择空间、建立估计和传递极限的次序保持不变。

\section{课程层次与先修知识}

本书面向学过数学分析与高等代数的读者，主体定位为本科高年级课程。先修知识包括极限、连续性、函数列、级数、线性空间、线性映射、内积与有限维特征值理论。涉及可测函数空间时，需要补充测度与 Lebesgue 积分；涉及一般谱理论时，需要复分析中的解析函数与 Liouville 定理。

\begin{description}
  \item[入门基础] 第\ref{chap:fa-metric}章至第\ref{chap:fa-baire}章建立度量空间、Banach 空间、Hilbert 空间及基本定理，是第一次学习的重点。
  \item[核心理论] 第\ref{chap:fa-weak}章至第\ref{chap:fa-compact-spectral}章讨论弱拓扑、弱紧性、算子谱与紧算子，构成完整基础课程的后半部分。
  \item[应用专题] 第\ref{chap:fa-variational}、\ref{chap:fa-fixed-point}章将抽象理论用于变分问题、弱解和非线性方程，可按课程需要选读。
  \item[进阶专题] 第\ref{chap:fa-unbounded}章至第\ref{chap:fa-semigroups}章介绍无界算子、一般谱定理和算子半群，作为进一步学习的入口。
\end{description}

附录\ref{app:fa-measure}补充测度与积分，附录\ref{app:fa-topology}补充一般拓扑与集合论，附录\ref{app:fa-complex}补充复分析工具，附录\ref{app:fa-solutions}提供部分习题解答。尚未学习实分析的读者，可先用连续函数空间与数列空间理解主体概念，再结合附录进入可测函数空间。

\section{全书章节总览}

以下顺序兼顾知识依赖与学习动机。首次学习以基础与核心部分为主，应用部分择要阅读，进阶部分用于了解后续方向。

\begin{longtable}{L{0.08\textwidth}L{0.32\textwidth}L{0.48\textwidth}}
\toprule
章节 & 主题 & 核心问题 \\
\midrule
\endhead
\ref{chap:fa-introduction} & 从线性代数到泛函分析 & 从有限维到无限维，哪些方法仍然有效？ \\
\ref{chap:fa-metric} & 度量空间、完备性与紧性 & 极限何时存在，序列何时能够提取收敛子列？ \\
\ref{chap:fa-banach} & 赋范线性空间与 Banach 空间 & 怎样选择函数空间，怎样识别无限维现象？ \\
\ref{chap:fa-operators} & 有界线性算子与对偶空间 & 如何用范数控制线性映射，如何用标量观测向量？ \\
\ref{chap:fa-hilbert} & Hilbert 空间与正交展开 & 投影、最佳逼近和正交坐标如何推广？ \\
\ref{chap:fa-hahn-banach} & Hahn--Banach 定理与凸集分离 & 局部定义的泛函如何延拓，如何分离点与凸集？ \\
\ref{chap:fa-baire} & Baire 纲定理与算子基本定理 & 完备性怎样产生一致估计和逆映射的连续性？ \\
\ref{chap:fa-weak} & 弱拓扑、弱星拓扑与收敛 & 放宽收敛要求以后，哪些线性结构仍被保留？ \\
\ref{chap:fa-reflexive} & 弱紧性与自反空间 & 何时能从有界序列中提取弱收敛子列？ \\
\ref{chap:fa-spectrum} & 有界算子的谱理论 & 无限维中的不可逆性为何不限于特征值？ \\
\ref{chap:fa-compact} & 紧算子与 Fredholm 理论 & 哪些无限维方程仍具有有限维式的可解性结构？ \\
\ref{chap:fa-compact-spectral} & 紧自伴算子的谱分解 & 哪些算子能按正交特征向量展开？ \\
\ref{chap:fa-variational} & 变分方法与方程的弱解 & 如何由能量、弱紧性与强制性证明解的存在？ \\
\ref{chap:fa-fixed-point} & 不动点方法与非线性方程 & 非线性方程何时可由压缩性或紧性求解？ \\
\ref{chap:fa-unbounded} & 无界算子与定义域 & 微分算子的定义域为何是算子的一部分？ \\
\ref{chap:fa-general-spectral} & 一般谱定理与函数演算 & 无法用离散特征向量表示时，如何进行谱分解？ \\
\ref{chap:fa-semigroups} & 算子半群与演化方程 & 如何把矩阵指数推广为无限维时间演化？ \\
\bottomrule
\end{longtable}

\section{学习路线与贯穿全书的例子}

基本路线依次是空间结构、连续线性映射、基本存在性定理、弱收敛与紧性、算子谱及应用。Hilbert 空间的几何方法可以先于一般对偶理论学习；Hahn--Banach 定理与 Baire 纲定理各有独立的证明来源，不能将二者混为一条推导链。

全书反复使用连续函数空间、数列空间、可积函数空间，以及移位算子、对角算子、积分算子和微分算子。第一次出现时给出定义，此后在新的结构下重新考察同一对象。这样能够具体比较：同一线性映射换用范数后可能失去连续性；同一有界序列在弱拓扑下可能收敛；同一微分表达式采用不同边界条件会给出不同算子。

每章习题按概念辨析、证明训练、计算与应用组织。解题时先写清空间、标量域、范数或拓扑，再检查所用定理的假设。反例与正面结论同样重要，尤其要追踪完备性、闭性、有限维性、紧性和自反性分别用在何处。

第一次阅读时，可先以 $C([0,1])$、$\ell^1$、$\ell^2$ 和 $c_0$ 为四个主要模型，把涉及一般测度空间的段落留待第二遍阅读。弱紧性一章先掌握 Hilbert 空间中的子列构造，再读一般拓扑证明；谱理论一章先理解对角算子与乘法算子，再进入依赖复分析的谱半径证明。应用章只调用已经说明的结论；进阶章中引用的谱定理和生成定理均明确标出证明范围。

\section{记号与约定}

标量域记为 $\mathbb K$，取 $\mathbb R$ 或 $\mathbb C$。一般赋范线性空间记为 $X,Y$，Hilbert 空间记为 $H$，零向量记为 $0$。$\|x\|_X$ 表示 $X$ 中向量 $x$ 的范数；上下文明确时省略下标。$B_X$ 表示闭单位球 $\{x\in X:\|x\|_X\leq 1\}$。对集合 $M\subset X$，$\overline M$ 默认表示范数闭包，$\operatorname{span}M$ 表示有限线性组合的集合。

连续线性映射 $T:X\to Y$ 的全体记为 $\mathcal B(X,Y)$，并令 $\mathcal B(X)=\mathcal B(X,X)$；其核与值域分别记为 $\ker T$、$\operatorname{Ran}T$。连续线性泛函构成的对偶空间记为 $X^*$，双对偶记为 $X^{**}$。恒等算子记为 $I$。

复内积 $\langle x,y\rangle$ 约定对第一个变量线性、对第二个变量共轭线性。Banach 空间算子的对偶算子记为 $T':Y^*\to X^*$，Hilbert 空间算子的伴随记为 $T^*$，以免混淆两种构造。涉及弱星拓扑时使用 $\sigma(X^*,X)$，涉及弱拓扑时使用 $\sigma(X,X^*)$；完整定义见第\ref{chap:fa-weak}章。

\section{参考资源与本书的组织方式}

核心主题的覆盖范围参考 MIT 的泛函分析入门课程大纲\cite{mit-syllabus}，相关公开讲义可由课程阅读页\cite{mit-notes}获取。本书采用自己的章节顺序，将弱紧性与自反空间单列，并另设应用和进阶部分。章节总览用于组织知识，不等同于任何一门课程的考试范围。

\chapter{度量空间、完备性与紧性}
\label{chap:fa-metric}

函数列可以逐点收敛，也可以一致收敛；同一列函数在不同意义下可能具有不同的极限性质。度量把“接近”变成可以估计的数值，完备性保证 Cauchy 逼近有极限，紧性则保证从任意序列中选出收敛子列。

\section{度量与基本拓扑概念}
\begin{definition}
集合 $M$ 上的\keyterm{度量（metric）}是映射 $d:M\times M\to[0,\infty)$，满足对所有 $x,y,z\in M$，
\[
 d(x,y)=0\iff x=y,\qquad d(x,y)=d(y,x),\qquad
 d(x,z)\leq d(x,y)+d(y,z).
\]
二元组 $(M,d)$ 称为度量空间。开球记为 $B(x,r)=\{y:d(x,y)<r\}$，其中 $r>0$。
\end{definition}
集合 $U$ 开，是指每个 $x\in U$ 都有某个开球 $B(x,r)\subset U$；集合 $F$ 闭，是指 $M\setminus F$ 开。集合 $A$ 的闭包 $\overline A$ 由所有满足 $B(x,r)\cap A\ne\varnothing$（对每个 $r>0$）的点组成。若 $\overline A=M$，称 $A$ 稠密；存在可数稠密子集的空间称为\keyterm{可分空间（separable space）}。

度量空间中，$x\in\overline A$ 当且仅当存在 $a_n\in A$ 使 $d(a_n,x)\to0$：一个方向由三角不等式得到，另一个方向在每个 $B(x,1/n)$ 中选择一点即可。因此，闭集恰是包含其所有收敛序列极限的集合。

这里的开、闭总是相对于所在空间而言。比如 $[0,1)$ 在 $\mathbb R$ 中既不开也不闭，在子空间 $[0,2]$ 中却是开集，因为它等于 $(-1,1)\cap[0,2]$。闭性也不表示集合“没有空隙”：整数集在 $\mathbb R$ 中闭，但并不是区间。它只表示从集合内取点所得到的、在环境空间中存在的极限不会离开集合。开集与闭集不是互为否定的两个类别，整个空间和空集既开又闭。

\begin{proposition}[距离函数的连续性]\label{prop:fa-distance-lipschitz}
对度量空间中的非空集合 $A$，令 $d(x,A)=\inf_{a\in A}d(x,a)$，则
\[
 |d(x,A)-d(y,A)|\leq d(x,y),\qquad
 d(x,A)=0\iff x\in\overline A.
\]
\end{proposition}
\begin{proof}
对任意 $a\in A$，三角不等式给出 $d(x,a)\leq d(x,y)+d(y,a)$。对 $a$ 取下确界，再交换 $x,y$，便得第一式。距离为零表示每个正半径的球都与 $A$ 相交，这恰是闭包的定义。下确界不一定达到；例如 $0$ 到 $(0,1)$ 的距离为零，但集合中没有距离为零的点。
\end{proof}

\begin{example}
在 $C([0,1])$ 上，$d_\infty(f,g)=\max_{t\in[0,1]}|f(t)-g(t)|$ 给出一致收敛。另一个度量 $d_1(f,g)=\int_0^1|f-g|\,\mathrm dt$ 给出不同的收敛：$f_n(t)=t^n$ 满足 $d_1(f_n,0)=1/(n+1)\to0$，但 $d_\infty(f_n,0)=1$。
\end{example}

\section{连续性与完备性}
映射 $F:(M,d)\to(N,\rho)$ 在 $x$ 连续，是指对每个 $\varepsilon>0$ 存在 $\delta>0$，使 $d(x,y)<\delta$ 蕴含 $\rho(Fx,Fy)<\varepsilon$。若 $\delta$ 可对所有 $x$ 统一选择，则称一致连续。用半径 $1/n$ 选择反例可证明：度量空间中的连续性等价于保持序列极限。

\begin{definition}
若对每个 $\varepsilon>0$，存在 $N$ 使 $m,n\geq N$ 时 $d(x_m,x_n)<\varepsilon$，则称 $(x_n)$ 为 Cauchy 序列。每个 Cauchy 序列都在空间内收敛时，称空间\keyterm{完备（complete）}。
\end{definition}
收敛序列必为 Cauchy 序列，因为 $d(x_m,x_n)\leq d(x_m,x)+d(x_n,x)$。逆命题需要空间本身的性质：有理数列可以在实数中趋于 $\sqrt2$，却在 $\mathbb Q$ 中没有极限。

这两个定义的区别在于，收敛先指定一个候选极限 $x$，Cauchy 条件只比较已经算出的两项。在迭代计算中，通常先得到的是后者。证明一个空间完备的常用过程是：先在更简单的环境中找出候选极限，再证明它属于原空间，最后证明原序列按指定的度量趋于它。这三个步骤缺少任何一个，都不能完成完备性证明。

\begin{proposition}\label{prop:fa-closed-complete}
完备度量空间的子集在诱导度量下完备，当且仅当它闭。任意度量空间中的完备子集都是闭集。
\end{proposition}
\begin{proof}
闭子集中的 Cauchy 序列先在完备环境中收敛，再由闭性知极限仍在子集中。反之，若子集 $A$ 完备且 $x\in\overline A$，取 $a_n\in A$ 收敛到 $x$。此列为 Cauchy 序列，故在 $A$ 中有极限；度量极限唯一，因而 $x\in A$。
\end{proof}

\begin{example}[一致极限与完备性]
$C([0,1])$ 在 $d_\infty$ 下完备。
\end{example}
\begin{proof}
若 $(f_n)$ 为一致距离下的 Cauchy 序列，则对每个 $t$，$(f_n(t))$ 在 $\mathbb K$ 中收敛，记极限为 $f(t)$。在 $|f_n(t)-f_m(t)|<\varepsilon$ 中令 $m\to\infty$，得到对所有 $t$ 同时成立的 $|f_n(t)-f(t)|\leq\varepsilon$，故收敛一致。给定 $t_0$，将 $|f(t)-f(t_0)|$ 分解为两个一致逼近误差与 $|f_n(t)-f_n(t_0)|$，即知 $f$ 连续。
\end{proof}

最后一步可具体写成
\[
 |f(t)-f(t_0)|\leq |f(t)-f_n(t)|
 +|f_n(t)-f_n(t_0)|+|f_n(t_0)-f(t_0)|.
\]
先选一个固定的 $n$，使首尾两项各小于 $\varepsilon/3$；再用这个 $f_n$ 的连续性选择 $\delta$，使中间一项小于 $\varepsilon/3$。先选 $n$、后选 $\delta$ 的次序很关键：逐点收敛允许 $n$ 随 $t$ 改变，就不能直接作上述统一控制。

\section{完备化与连续延拓}
\begin{theorem}[完备化]
每个度量空间 $M$ 都可等距嵌入一个完备空间 $\widehat M$，且其像稠密。这样的完备化在保持原空间嵌入的等距同构意义下唯一。
\end{theorem}
\begin{proof}
在 $M$ 的 Cauchy 序列之间规定 $(x_n)\sim(y_n)$ 当且仅当 $d(x_n,y_n)\to0$，并令
\[
 \widehat d([x_n],[y_n])=\lim_n d(x_n,y_n).
\]
三角不等式保证右侧存在、与代表元无关，并使其成为等价类集合上的度量。常值序列给出等距嵌入 $\iota$；由于 $\iota(x_n)\to[x_n]$，其像稠密。

若 $(\xi_n)$ 为 $\widehat M$ 中的 Cauchy 序列，由稠密性取 $z_n\in M$ 使 $\widehat d(\iota(z_n),\xi_n)<1/n$。于是 $(z_n)$ 为 Cauchy 序列，而 $\xi_n\to[z_n]$，所以 $\widehat M$ 完备。在另一完备化中，把 $[x_n]$ 映为嵌入序列的极限；这个映射良定义、保距且满射，给出唯一性。
\end{proof}

若 $A$ 在 $M$ 中稠密，$F:A\to N$ 一致连续且 $N$ 完备，则令 $\overline F(x)=\lim_nF(a_n)$，其中 $a_n\to x$，可得唯一的一致连续延拓。存在性来自一致连续映射将 Cauchy 序列映为 Cauchy 序列；两种逼近序列交错排列即可验证极限无关性。仅有连续性通常不足，例如 $x\mapsto1/x$ 不能从 $(0,1)$ 连续延拓到 $[0,1]$。

\section{紧性、列紧性与全有界性}
\begin{definition}
度量空间称为\keyterm{紧的（compact）}，若每个开覆盖都有有限子覆盖；称为\keyterm{列紧的（sequentially compact）}，若每个序列都有收敛子列；称为\keyterm{全有界的（totally bounded）}，若对每个 $\varepsilon>0$，有限多个半径为 $\varepsilon$ 的开球就能覆盖整个空间。
\end{definition}
\begin{theorem}\label{thm:fa-metric-compactness}
对度量空间，紧、列紧、完备且全有界三者等价。
\end{theorem}
\begin{proof}
紧空间必全有界，因为全体半径 $\varepsilon$ 的球构成开覆盖。若一个序列没有收敛子列，每个点都有仅包含有限多个序列指标的邻域；否则可逐次选点得到收敛子列。这些邻域的有限子覆盖不可能覆盖所有指标，故紧空间列紧。

列紧空间完备，因为 Cauchy 序列一旦有收敛子列，整个序列就收敛。若不全有界，可递归构造点列，使任意两点距离至少为某个 $\varepsilon>0$，与列紧性矛盾。

完备且全有界时，对半径 $2^{-k}$ 的有限覆盖逐次保留含无穷多项的球，抽得 Cauchy 子列，故列紧。最后，列紧空间的任意开覆盖具有 Lebesgue 数：存在 $\delta>0$，每个半径 $\delta$ 的球都包含于某个覆盖成员。否则选择 $x_n$ 使 $B(x_n,1/n)$ 不被任何成员包含，取 $x_n$ 的收敛子列，在极限点附近的一个覆盖成员内即得矛盾。再取有限个半径 $\delta/2$ 的球覆盖空间，得到有限子覆盖。
\end{proof}
完备与有界并不保证紧。例如无穷集合上的离散度量 $d(x,y)=1$（$x\ne y$）完备且有界，但半径 $1/2$ 的球都是单点，故不全有界。

\begin{example}[三种条件分别控制什么]
在通常距离下，$(0,1)$ 全有界但不完备，$\mathbb R$ 完备但不全有界，$[0,1]$ 同时满足两者，因而紧。
\end{example}
\begin{solve}
任意小的正数 $\varepsilon$ 都允许把 $(0,1)$ 分成有限个长度小于 $\varepsilon$ 的区间，所以它全有界；但 $1/n$ 的极限 $0$ 不在其中。实数空间完备，而有限个固定半径的球只能覆盖一个有界区域，不能覆盖全部实数。闭区间既保留所有 Cauchy 极限，又能被有限小区间覆盖。完备性防止极限落到空间外，全有界性防止序列在固定精度下出现无穷多个彼此分离的位置。
\end{solve}

紧性的一个直接用途是证明最值存在。若非空紧集 $K$ 上的实连续函数 $g$ 无界，可取 $x_n\in K$ 使 $|g(x_n)|>n$，收敛子列与连续性立刻造成矛盾。于是 $m=\inf_Kg$ 为有限数；取 $g(x_n)\to m$，再提取收敛子列 $x_{n_k}\to x_*$，便有 $g(x_*)=m$。后面证明能量极小点存在时，仍采用“取极小化序列、提取子列、传递极限”的结构，只是所用紧性和连续性会改变。

\section{压缩映射原理}
\begin{theorem}[Banach 不动点定理]\label{thm:fa-contraction}
设 $(M,d)$ 非空且完备，$F:M\to M$ 满足 $d(Fx,Fy)\leq qd(x,y)$，其中 $0\leq q<1$。则 $F$ 有唯一不动点 $x_*$。任取 $x_0\in M$，迭代 $x_{n+1}=F(x_n)$ 收敛到 $x_*$，且
\[
 d(x_n,x_*)\leq\frac{q^n}{1-q}d(x_1,x_0),\qquad
 d(x_n,x_*)\leq\frac{q}{1-q}d(x_n,x_{n-1})\quad(n\geq1).
\]
\end{theorem}
\begin{proof}
归纳得 $d(x_{n+1},x_n)\leq q^nd(x_1,x_0)$。对相邻距离求和可知 $(x_n)$ 为 Cauchy 序列，故有极限 $x_*$。压缩条件蕴含连续性，因此 $F(x_*)=x_*$。若 $y_*$ 也是不动点，则 $d(x_*,y_*)\leq qd(x_*,y_*)$，只能有 $x_*=y_*$。将上述尾部几何级数求和即得两项误差估计。
\end{proof}

\section{例题与习题}
\begin{example}
方程 $x=\cos x$ 在 $[0,1]$ 中有唯一解，迭代 $x_{n+1}=\cos x_n$ 对任意 $x_0\in[0,1]$ 收敛。
\end{example}
\begin{solve}
$F([0,1])\subset[0,1]$，且由中值定理 $|F(x)-F(y)|\leq(\sin1)|x-y|$。区间完备且 $\sin1<1$，故可用定理\ref{thm:fa-contraction}。若希望误差不超过 $\varepsilon$，可在 $\frac{\sin1}{1-\sin1}|x_n-x_{n-1}|\leq\varepsilon$ 时停止。
\end{solve}

使用压缩映射原理时，应分别验证集合非空、完备、映射保持集合以及统一收缩常数小于一。比如 $F(x)=x/2$ 在 $(0,1)$ 上是压缩自映射，却没有不动点，因为唯一可能的不动点 $0$ 被排除；在 $[1,2]$ 上同一个公式虽然满足收缩估计，却不保持集合。再如 $F(x)=x+1$ 在 $\mathbb R$ 上满足 Lipschitz 常数一而没有不动点，$F(x)=x$ 则有无穷多个不动点。这些例子分别指出各项条件的作用。
\begin{enumerate}
 \item\label{exer:fa-cauchy} 证明 Cauchy 序列有界，并证明它若有收敛子列则整个序列收敛。
 \item 比较 $(0,1)$、$[0,1]$、$\mathbb R$ 的完备性、全有界性和紧性。
 \item 证明紧度量空间上的连续映射一致连续。
 \item 给出 $q=1$ 时不动点不存在或不唯一的例子。
\end{enumerate}

\chapter{赋范线性空间与 Banach 空间}
\label{chap:fa-banach}

在向量空间中，极限还必须与加法、数乘相容。范数同时描述向量的大小和向量间的距离，是建立无限维线性分析的基本结构。

\section{范数、半范数与等价范数}
\begin{definition}
$\mathbb K$ 上线性空间 $X$ 的\keyterm{范数（norm）}是映射 $\|\cdot\|:X\to[0,\infty)$，满足
\[
 \|x\|=0\iff x=0,\qquad \|\alpha x\|=|\alpha|\|x\|,\qquad
 \|x+y\|\leq\|x\|+\|y\|.
\]
不要求第一项的正定性时称为半范数（seminorm）。范数诱导度量 $d(x,y)=\|x-y\|$；此度量完备时，$X$ 称为\keyterm{Banach 空间（Banach space）}。
\end{definition}
反三角不等式 $|\|x\|-\|y\||\leq\|x-y\|$ 表明范数连续。若 $x_n\to x$、$y_n\to y$、$\alpha_n\to\alpha$，则三角不等式以及 $(x_n)$ 的有界性给出 $x_n+y_n\to x+y$ 和 $\alpha_nx_n\to\alpha x$。

同一空间的两个范数 $\|\cdot\|_a,\|\cdot\|_b$ 称为等价，若存在 $c,C>0$ 使
\[
 c\|x\|_a\leq\|x\|_b\leq C\|x\|_a\quad(x\in X).
\]
两侧估计表明等价范数有相同的收敛序列、Cauchy 序列和开集，因而完备性也相同。反过来，若两个范数给出相同的拓扑，则两个恒等映射在零点连续，利用线性缩放就得到上述估计。

\section{数列空间与连续函数空间}
有限维中常见的三个范数是 $\|x\|_1=\sum_{j=1}^d|x_j|$、$\|x\|_2=(\sum_{j=1}^d|x_j|^2)^{1/2}$ 和 $\|x\|_\infty=\max_j|x_j|$。例如 $x=(3,4)$ 的三个范数分别是 $7,5,4$。它们衡量大小的方式不同，但满足 $\|x\|_\infty\leq\|x\|_2\leq\|x\|_1\leq d\|x\|_\infty$，因此给出相同的收敛。把有限求和改成无限求和时，首先要筛选使求和有意义的数列，这就产生下面的空间。

对 $1\leq p<\infty$ 定义
\[
 \ell^p=\left\{x=(x_j)_{j\geq1}:\sum_{j=1}^\infty|x_j|^p<\infty\right\},
 \qquad \|x\|_p=\left(\sum_{j=1}^\infty|x_j|^p\right)^{1/p}.
\]
另令 $\ell^\infty=\{x:\sup_j|x_j|<\infty\}$，$\|x\|_\infty=\sup_j|x_j|$；$c_0$ 为其中趋于零的数列，$c_{00}$ 为仅有限个坐标非零的数列。$\ell^p$ 的三角不等式是 Minkowski 不等式，见附录\ref{app:fa-measure}；计数测度将积分形式化为求和形式。

\begin{example}[一个空间中的点本身也可以是数列]
数列 $x=(1,1/2,1/3,\ldots)$ 属于 $\ell^2$ 与 $c_0$，但不属于 $\ell^1$；常数列 $\boldsymbol 1=(1,1,\ldots)$ 属于 $\ell^\infty$，却不属于 $c_0$。一般地，对 $a>0$，$(j^{-a})\in\ell^p$ 当且仅当 $ap>1$。
\end{example}
\begin{solve}
这些判断分别归结为级数 $\sum j^{-2}$ 收敛、调和级数发散，以及一般正项级数 $\sum j^{-ap}$ 的判别。注意区分两个指标：$x_j$ 是一个向量的第 $j$ 个坐标，而 $x^{(n)}$ 是空间中第 $n$ 个向量。$x^{(n)}\to x$ 是整条数列作为向量趋近另一条数列，不等于仅检查每个固定坐标。
\end{solve}

\begin{proposition}\label{prop:fa-sequence-complete}
$\ell^p$（$1\leq p\leq\infty$）和配备上确界范数的 $c_0$ 都是 Banach 空间。若 $K$ 为非空紧度量空间，则 $C(K)$ 在一致范数 $\|f\|_\infty=\max_{t\in K}|f(t)|$ 下也是 Banach 空间。
\end{proposition}
\begin{proof}
先设 $p<\infty$，$(x^{(n)})$ 为 $\ell^p$ 中的 Cauchy 序列。由 $|x_j^{(n)}-x_j^{(m)}|\leq\|x^{(n)}-x^{(m)}\|_p$，每个坐标收敛，记极限为 $x_j$。若 $n,m\geq N$ 时范数差小于 $\varepsilon$，则先对有限个坐标求和并令 $m\to\infty$，再令坐标数趋于无穷，得到
\[
 \sum_j|x_j^{(n)}-x_j|^p\leq\varepsilon^p\quad(n\geq N).
\]
因此 $x-x^{(N)}\in\ell^p$，从而 $x\in\ell^p$ 且 $x^{(n)}\to x$。$p=\infty$ 时，逐坐标取极限后保留同一上界即可。$c_0$ 在 $\ell^\infty$ 中闭：先用一个趋零数列一致逼近极限数列，再控制其尾部。$C(K)$ 的证明与上一章 $C([0,1])$ 的一致极限证明相同。
\end{proof}
截断前 $N$ 个坐标表明，$c_{00}$ 在每个 $\ell^p$（$p<\infty$）和 $c_0$ 中稠密。用实部、虚部为有理数的有限支撑数列逼近，可见这些空间可分；$\ell^\infty$ 不可分，因为所有零一数列构成不可数集，任意两项的距离为一。

这里截断误差可以直接计算。令 $P_Nx=(x_1,\ldots,x_N,0,\ldots)$，则 $p<\infty$ 时 $\|x-P_Nx\|_p^p=\sum_{j>N}|x_j|^p\to0$；在 $c_0$ 中误差为 $\sup_{j>N}|x_j|\to0$。但对 $\boldsymbol1\in\ell^\infty$，截断误差始终等于一。为了补全不可分性的论证，若存在可数稠密集，每个零一数列附近半径 $1/3$ 的开球都应含有稠密集的一点；这些球两两不交，这将需要不可数多个不同的稠密集元素，矛盾。

\section{可测函数空间与积分范数}
给定测度空间 $(\Omega,\mathcal A,\mu)$，$L^p(\mu)$ 的元素是几乎处处相等的可测函数的等价类。对 $1\leq p<\infty$ 令
\[
 \|f\|_p=\left(\int_\Omega|f|^p\,\mathrm d\mu\right)^{1/p},
 \qquad
 \|f\|_\infty=\inf\{C\geq0:|f|\leq C\ \text{几乎处处}\}.
\]
只保留范数有限的等价类。取等价类是必要的：在原函数集合上，积分范数为零并不意味着处处为零。

比如在 Lebesgue 测度下，恒零函数与只在 $t=1/2$ 处取一的函数代表同一个 $L^p(0,1)$ 元素。因此对一个抽象的 $L^p$ 元素说“它在 $1/2$ 处的值”，通常没有定义；但它的积分和范数都有定义。$\|f\|_\infty$ 中的本质上确界（essential supremum）也忽略零测集上的异常值，与对每一点取上确界不同。

对 $f(t)=t^{-a}$、$a>0$，有 $f\in L^p(0,1)$ 当且仅当 $ap<1$，因为 $\int_0^1t^{-ap}\,\mathrm dt$ 的收敛条件如此。这个条件与数列 $(j^{-a})$ 的条件方向相反：积分例子检查零点附近的奇性，数列例子检查无穷远的尾部。在有限测度空间上，若 $1\leq p<q\leq\infty$，Hölder 不等式给出
\[
 \|f\|_p\leq\mu(\Omega)^{1/p-1/q}\|f\|_q
 \quad(0<\mu(\Omega)<\infty),
\]
故 $L^q\subset L^p$；数列空间则是 $\ell^p\subset\ell^q$。使用包含关系前，必须先说明测度空间，不能只记一个箭头方向。

\begin{theorem}\label{thm:fa-lp-complete}
对任意测度空间及 $1\leq p\leq\infty$，$L^p(\mu)$ 是 Banach 空间。
\end{theorem}
\begin{proof}
设 $p<\infty$。从 Cauchy 序列 $(f_n)$ 选取子列，使 $\|f_{n_{k+1}}-f_{n_k}\|_p\leq2^{-k}$。令 $g_k=|f_{n_{k+1}}-f_{n_k}|$。Minkowski 不等式和单调收敛定理给出 $\|\sum_kg_k\|_p\leq\sum_k2^{-k}<\infty$，故该和几乎处处有限。于是函数级数绝对收敛，定义
\[
 f=f_{n_1}+\sum_{k=1}^{\infty}(f_{n_{k+1}}-f_{n_k})
 \quad\text{几乎处处}.
\]
尾项的 $L^p$ 范数不超过相应几何级数尾和，故子列在范数下趋于 $f\in L^p$，原 Cauchy 序列也趋于 $f$。当 $p=\infty$ 时，去掉控制各对差的可数个零测集后，序列在剩余集合上一致 Cauchy，逐点极限便给出 $L^\infty$ 极限。
\end{proof}
范数收敛不要求原序列逐点收敛。例如在 $[0,1)$ 上依次列出各级二进区间的示性函数，其 $L^p$ 范数在 $p<\infty$ 时趋于零，但每个点都被无穷多个区间包含，也被无穷多个区间排除。

\begin{example}[为什么要从连续函数走向可积函数]\label{ex:fa-incomplete-integral}
对 $n\geq3$，令 $f_n$ 在 $[0,1/2]$ 上为零，在 $[1/2,1/2+1/n]$ 上线性地从零升到一，在剩余区间上为一。则 $(f_n)$ 在积分范数下为 Cauchy 序列，但不在 $C([0,1])$ 中收敛。
\end{example}
\begin{solve}
令 $g=\mathbf1_{(1/2,1]}$。两者只在宽度 $1/n$ 的过渡区间有差别，三角形面积给出 $\|f_n-g\|_1=1/(2n)$，所以 $\|f_n-f_m\|_1\leq1/(2n)+1/(2m)$。若存在连续 $f$ 使 $\|f_n-f\|_1\to0$，则 $\|f-g\|_1=0$。连续性迫使 $f$ 在 $(0,1/2)$ 上恒为零、在 $(1/2,1)$ 上恒为一：否则某点附近有正长度区间上的差不为零，积分不可能为零。这与 $f$ 在 $1/2$ 连续矛盾。新空间 $L^1$ 正好容纳了这个极限，而一致范数下的连续函数空间本来就是完备的。
\end{solve}

\section{子空间、商空间与级数判别}
Banach 空间的线性子空间完备当且仅当它闭，这是命题\ref{prop:fa-closed-complete}的直接应用。若 $M\subset X$ 为闭线性子空间，在商空间 $X/M$ 上定义
\begin{equation}
 \|x+M\|_{X/M}=\inf_{m\in M}\|x-m\|_X.
 \label{eq:fa-quotient-norm}
\end{equation}
右侧也记为 $\operatorname{dist}(x,M)$，即点到集合的距离。更换代表元不会改变右侧；齐次性和三角不等式由取近似极小元得到。范数为零等价于 $x\in\overline M=M$，所以闭性保证正定性。

商空间把相差一个 $M$ 中元素的向量视为相同。其零元素是整个集合 $M$，而不是只含原点的集合。例如 $X=\mathbb R^2$、$M=\mathbb R\times\{0\}$ 时，$(a,b)+M$ 是高度为 $b$ 的水平直线；在欧氏范数下，$\|(a,b)+M\|=|b|$，故商空间等距对应于纵坐标轴。函数空间中，对常数函数构成的子空间取商，就把“相差一个常数”的函数视为同一对象。微分算子无法区分这些函数，因而商空间自然适合消除这类不唯一性。

商范数中的下确界未必有最小元。证明三角不等式时只需对任意 $\varepsilon>0$，分别选 $m_1,m_2\in M$，使 $\|x-m_1\|<\|x+M\|+\varepsilon$、$\|y-m_2\|<\|y+M\|+\varepsilon$。用 $m_1+m_2$ 作为 $x+y$ 的一个候选代表，再令 $\varepsilon\downarrow0$ 即可；不应事先假定存在最近点。

\begin{proposition}\label{prop:fa-series-complete}
赋范空间完备当且仅当每个满足 $\sum_n\|x_n\|<\infty$ 的级数 $\sum_nx_n$ 都收敛。Banach 空间对闭线性子空间取商后仍为 Banach 空间。
\end{proposition}
\begin{proof}
完备空间中，绝对可和级数的部分和为 Cauchy 序列。反之，对任意 Cauchy 序列 $(y_n)$ 选子列使 $\|y_{n_{k+1}}-y_{n_k}\|\leq2^{-k}$；差分级数收敛，所以子列及原序列均收敛。

设 $\sum_n\|\xi_n\|_{X/M}<\infty$。对每个商类 $\xi_n$ 选择代表元 $x_n$ 使 $\|x_n\|\leq\|\xi_n\|_{X/M}+2^{-n}$。则 $\sum_nx_n$ 在 $X$ 中收敛，连续商映射把其和映为 $\sum_n\xi_n$ 的极限。由刚证的级数判别，$X/M$ 完备。
\end{proof}
赋范空间的度量完备化可通过代表序列逐项相加、数乘，赋予唯一延续原运算的 Banach 空间结构；范数定义为原序列范数的极限。

级数判别中的“绝对”不能省略，也不表示所有收敛级数都绝对收敛。在 $\ell^2$ 中，$\sum_{n\geq1}e_n/n$ 收敛到 $(1/n)$，因为尾项范数平方为 $\sum_{n>N}1/n^2\to0$；但各项范数之和为 $\sum1/n=\infty$。完备性保证绝对可和是收敛的充分条件，并未把它变成必要条件。

\section{有限维定理与 Riesz 引理}
\begin{theorem}\label{thm:fa-finite-dimension}
有限维线性空间上的所有范数等价；有限维赋范空间完备，其闭有界集紧；任意赋范空间的有限维线性子空间闭。
\end{theorem}
\begin{proof}
取基 $e_1,\ldots,e_d$，写 $x=\sum_ja_je_j$。三角不等式给出 $\|x\|\leq C\|a\|_2$，故 $a\mapsto\|\sum_ja_je_j\|$ 在欧氏单位球面上连续且正，存在正最小值 $c$。由齐次性得 $c\|a\|_2\leq\|x\|$，从而任意范数等价于坐标欧氏范数。完备与紧性由欧氏空间转移；有限维子空间完备，因而在环境空间中闭。
\end{proof}
\begin{lemma}[Riesz 引理]\label{lem:fa-riesz}
若 $M$ 是赋范空间 $X$ 的真闭线性子空间，$0<\alpha<1$，则存在单位向量 $x$ 满足 $\inf_{m\in M}\|x-m\|>\alpha$。
\end{lemma}
\begin{proof}
取 $y\notin M$，令 $d=\inf_{m\in M}\|y-m\|>0$，选 $m_0\in M$ 使 $\|y-m_0\|<d/\alpha$。令 $x=(y-m_0)/\|y-m_0\|$，则对任意 $m\in M$，$\|x-m\|\geq d/\|y-m_0\|>\alpha$。
\end{proof}
\begin{corollary}\label{cor:fa-unit-ball}
赋范空间的闭单位球紧当且仅当空间有限维。
\end{corollary}
\begin{proof}
无限维时，对逐步生成的有限维子空间反复使用引理\ref{lem:fa-riesz}，得到单位向量 $x_n$，任意两项距离大于 $1/2$，因而闭单位球不列紧。有限维方向由定理\ref{thm:fa-finite-dimension}得到。
\end{proof}

Riesz 引理在没有内积的空间中起到寻找“新方向”的作用。在 Hilbert 空间里可以选择与已有有限维空间正交的单位向量，其距离恰为一；一般赋范空间没有角度，却仍能找到与已有空间保持固定距离的单位向量。由此可见，无限维单位球不紧并非某个特殊范数造成的偶然现象，而是任意范数都会遇到的结构性障碍。

\section{例题与习题}
\begin{example}[不同范数的作用]
$C([0,1])$ 上的 $\|f\|_1=\int_0^1|f|$ 与 $\|f\|_\infty$ 不等价。
\end{example}
\begin{solve}
虽有 $\|f\|_1\leq\|f\|_\infty$，但取 $f_n(t)=\max(1-nt,0)$，则 $\|f_n\|_\infty=1$、$\|f_n\|_1=1/(2n)$，反向统一估计不可能成立。这说明在无限维空间中，改变范数可能改变收敛与连续性。
\end{solve}
\begin{enumerate}
 \item 证明 $c_{00}$ 在 $\ell^\infty$ 中的闭包为 $c_0$。
 \item\label{exer:fa-lp-inclusion} 证明 $\ell^p\subset\ell^q$ 且 $\|x\|_q\leq\|x\|_p$，其中 $1\leq p<q\leq\infty$。
 \item 在 $\mathbb R^2$ 的欧氏范数下，计算对 $x$ 轴取商后的范数。
 \item 对例\ref{ex:fa-incomplete-integral}中的函数列，计算 $\|f_n-g\|_p$（$1\leq p<\infty$），说明为何同一反例适用于每个有限积分范数。
\end{enumerate}

\chapter{有界线性算子与对偶空间}
\label{chap:fa-operators}

有限维线性映射由矩阵决定，而且自动连续。无限维空间中，线性本身不控制极限；需要一个与向量无关的估计，才能把线性代数与分析连接起来。

\section{连续性与有界性}
\begin{theorem}\label{thm:fa-bounded-continuous}
赋范空间间的线性映射 $T:X\to Y$ 连续，当且仅当它在零点连续，当且仅当存在 $C\geq0$ 使 $\|Tx\|_Y\leq C\|x\|_X$ 对所有 $x\in X$ 成立。
\end{theorem}
\begin{proof}
若 $T$ 在零点连续，取 $\delta>0$ 使 $\|u\|<\delta$ 时 $\|Tu\|<1$。对 $x\ne0$ 令 $u=\delta x/(2\|x\|)$，得到 $\|Tx\|\leq2\|x\|/\delta$。反之，统一估计给出 $\|Tx-Ty\|\leq C\|x-y\|$，即处处连续。
\end{proof}
满足上述估计的线性映射称为\keyterm{有界线性算子（bounded linear operator）}，其范数为
\begin{equation}
 \|T\|=\sup_{\|x\|_X\leq1}\|Tx\|_Y.
 \label{eq:fa-operator-norm}
\end{equation}
它等于估计中最小的常数。这里“有界”指有界集的像有界；非零线性映射在整个空间上的像不可能有界，因为 $T(nx)=nTx$。

求算子范数通常分两步：先对任意输入证明 $\|Tx\|\leq C\|x\|$，得到上界；再构造单位向量或一列单位向量，使输出范数达到或趋近 $C$，得到下界。仅证明一个估计并不能断言其中的常数最优。对非零定义域，还有 $\|T\|=\sup_{x\ne0}\|Tx\|/\|x\|$，它表示线性变换对误差的最大放大倍数。

\begin{example}[积分算子的精确范数]\label{ex:fa-volterra-norm}
在 $C([0,1])$ 的一致范数下，$Vf(t)=\int_0^t f(s)\,\mathrm ds$ 是范数为一的有界线性算子。
\end{example}
\begin{solve}
积分的线性性保证 $V$ 线性，且 $|Vf(t)|\leq t\|f\|_\infty\leq\|f\|_\infty$。因此 $\|V\|\leq1$。取 $f(t)=1$，则 $Vf(t)=t$，输入与输出的一致范数都为一，所以 $\|V\|=1$。这同时表明积分操作在这一范数下不会放大扰动：$\|V(f+h)-Vf\|_\infty\leq\|h\|_\infty$。
\end{solve}

\begin{example}[上确界未必达到]
在 $\ell^2$ 上令 $Dx=(a_jx_j)$，其中 $a_j=1-1/(j+1)$。则 $\|D\|=1$，但不存在单位向量 $x$ 使 $\|Dx\|=1$。
\end{example}
\begin{solve}
一般有界对角算子满足 $\|D\|=\sup_j|a_j|$：上界逐坐标估计，下界用 $e_j$ 测试。本例的上确界为一，但每个 $a_j<1$。任意非零 $x$ 至少有一个非零坐标，故 $\sum a_j^2|x_j|^2<\sum|x_j|^2$；另一方面 $\|De_j\|=a_j\to1$。因此范数定义必须使用上确界，而不能无条件改成最大值。
\end{solve}

\section{算子空间与稠密延拓}
$\mathcal B(X,Y)$ 在逐点运算及算子范数下是赋范空间。若 $S\in\mathcal B(Y,Z)$，则 $\|ST\|\leq\|S\|\|T\|$。
\begin{proposition}\label{prop:fa-operator-complete}
若 $Y$ 为 Banach 空间，则 $\mathcal B(X,Y)$ 完备，无须假定 $X$ 完备。
\end{proposition}
\begin{proof}
若 $(T_n)$ 在算子范数下 Cauchy，则对每个 $x$，$(T_nx)$ 在 $Y$ 中 Cauchy。定义 $Tx=\lim_nT_nx$，逐点取极限可知 $T$ 线性。由 $\sup_n\|T_n\|<\infty$ 得 $T$ 有界；在 $\|(T_n-T_m)x\|\leq\varepsilon\|x\|$ 中令 $m\to\infty$，得到 $\|T_n-T\|\leq\varepsilon$。
\end{proof}
若 $D\subset X$ 是稠密线性子空间，$Y$ 完备，且 $T:D\to Y$ 有界，则
\[
 \widetilde T x=\lim_nTx_n\quad(x_n\in D,\ x_n\to x)
\]
给出唯一有界线性延拓，且 $\|\widetilde T\|=\|T\|$。极限的存在与无关性由 $\|Tx_n-Ty_n\|\leq\|T\|\|x_n-y_n\|$ 保证，线性和范数估计在极限下保持。

\section{连续对偶与双对偶}
\begin{definition}
线性映射 $f:X\to\mathbb K$ 称为线性泛函（linear functional）。所有连续线性泛函组成\keyterm{对偶空间（dual space）}$X^*=\mathcal B(X,\mathbb K)$；所有线性泛函组成的代数对偶通常更大。
\end{definition}
由于标量域完备，命题\ref{prop:fa-operator-complete}说明 $X^*$ 总是 Banach 空间。对每个 $x\in X$，规定
\[
 (Jx)(f)=f(x)\quad(f\in X^*).
\]
则 $Jx\in X^{**}$ 且 $\|Jx\|\leq\|x\|$。第\ref{chap:fa-hahn-banach}章将证明等号成立；这意味着全部连续线性观测能够恢复向量的范数。

\begin{example}[点值与积分]
在 $C([0,1])$ 的一致范数下，$\delta_t(f)=f(t)$ 与 $I(f)=\int_0^1f(s)\,\mathrm ds$ 均为范数一的泛函。
\end{example}
\begin{solve}
两者绝对值都不超过 $\|f\|_\infty$，故范数至多一；取常值函数 $f=1$ 达到一。若改用积分范数，点值泛函不连续：在 $t$ 附近取高度一、支撑长度趋零的连续尖峰函数即可。
\end{solve}

\section{数列空间的对偶表示}
\begin{theorem}\label{thm:fa-sequence-dual}
设 $1\leq p<\infty$、$1/p+1/q=1$，其中 $p=1$ 时 $q=\infty$。每个 $f\in(\ell^p)^*$ 唯一表示为
\[
 f(x)=\sum_{j=1}^{\infty}a_jx_j,\qquad a\in\ell^q,
 \qquad\|f\|=\|a\|_q.
\]
此外，$c_0^*$ 通过同一配对等距同构于 $\ell^1$。复数情形这里采用不带共轭的双线性配对。
\end{theorem}
\begin{proof}
Hölder 不等式先说明每个 $a\in\ell^q$ 都给出范数至多 $\|a\|_q$ 的泛函。反之，令 $a_j=f(e_j)$，其中 $e_j$ 是第 $j$ 个标准单位向量。有限支撑向量上已有所需公式。

若 $1<p<\infty$，对前 $N$ 项令 $x_j=|a_j|^{q-1}\overline{a_j}/|a_j|$，零坐标取零。记 $A_N=\sum_{j\leq N}|a_j|^q$，则 $A_N=f(x)\leq\|f\|A_N^{1/p}$，所以 $A_N^{1/q}\leq\|f\|$。令 $N\to\infty$ 得 $a\in\ell^q$ 及范数的反向估计。$p=1$ 时由 $|a_j|\leq\|f\|$ 即可。最后由 $c_{00}$ 的稠密性延拓表示式。

对 $c_0$，取前 $N$ 个坐标为 $\overline{a_j}/|a_j|$ 的向量，得到 $\sum_{j\leq N}|a_j|\leq\|f\|$。因此 $a\in\ell^1$；再由稠密性得到表示及范数相等。各坐标值决定 $a$，故表示唯一。
\end{proof}
在 $\sigma$ 有限测度空间上，对 $1\leq p<\infty$ 同样有 $(L^p)^*\cong L^q$，配对为 $f\mapsto\int af\,\mathrm d\mu$。这项测度论表示使用 Radon--Nikodym 定理，本书作为补充结果引用，见附录\ref{app:fa-measure}。不能把 $p=\infty$ 代入而断言 $(L^\infty)^*=L^1$；类似地，$(\ell^\infty)^*$ 一般大于由 $\ell^1$ 配对给出的部分。

\section{核、值域与对偶算子}
有界线性算子的核是闭子空间，因为它是闭集 $\{0\}$ 的逆像。值域却未必闭。
\begin{example}\label{ex:fa-diagonal-range}
在 $\ell^2$ 上令 $T(x_j)=(x_j/j)$。则 $\|T\|=1$，$T$ 单射，其值域稠密但不闭。
\end{example}
\begin{solve}
范数估计由 $1/j\leq1$ 得到，取 $e_1$ 达到等号。值域包含所有有限支撑数列，故稠密；但 $y=(1/j)$ 属于 $\ell^2$，其唯一可能原像为 $(1,1,\ldots)$，不在 $\ell^2$，所以值域不是全空间，因而不闭。截断的 $y$ 给出值域中的序列收敛到值域外的点。
\end{solve}

对 $T\in\mathcal B(X,Y)$，其\keyterm{对偶算子（dual operator）}定义为 $T':Y^*\to X^*$，$T'f=f\circ T$。直接估计给出 $\|T'\|\leq\|T\|$。对 $M\subset X$，定义零化子 $M^\perp=\{f\in X^*:f|_M=0\}$，则
\[
 \ker T'=(\operatorname{Ran}T)^\perp.
\]
这是一条代数与连续性共同作用的关系：值域上的所有观测为零，恰好表示泛函经 $T$ 拉回后为零。第\ref{chap:fa-hahn-banach}章将据此刻画值域稠密性。

核、值域与解集有直接联系。方程 $Tx=y$ 有解当且仅当 $y\in\operatorname{Ran}T$；若已有一个解 $x_0$，全部解恰为 $x_0+\ker T$。所以单射保证“至多一个解”，满射保证“每份数据都有解”，两者都没有直接给出误差估计。在例\ref{ex:fa-diagonal-range}中，数据扰动 $y^{(n)}=e_n/n$ 趋于零，其唯一原像却为 $e_n$，范数始终为一。这表明逆问题会放大小量，即使正向算子十分连续，反向求解仍可能不稳定。

\section{例题与习题}
\begin{example}[求导与定义域范数]
求导映射 $D:C^1([0,1])\to C([0,1])$ 在定义域使用一致范数时无界；使用 $\|f\|_{C^1}=\|f\|_\infty+\|f'\|_\infty$ 时有界。
\end{example}
\begin{solve}
取 $f_n(t)=\sin(nt)$，则 $\|f_n\|_\infty\leq1$ 而 $\|Df_n\|_\infty=n$，不存在统一估计。换用 $C^1$ 范数后，$\|Df\|_\infty\leq\|f\|_{C^1}$。由此可见，算子的连续性需要同时指定定义域与值域的范数。
\end{solve}
\begin{enumerate}
 \item 对 $\ell^p$ 上左右移位 $L(x_1,x_2,\ldots)=(x_2,x_3,\ldots)$、$R(x)=(0,x_1,x_2,\ldots)$，计算范数、核与值域。
 \item 证明有界线性算子在某个稠密子集上的值决定整个算子。
 \item 证明 $(ST)'=T'S'$，并注明各映射的定义域和值域。
 \item 对有界数列 $(a_j)$，求 $T(x_j)=(a_jx_j)$ 在 $\ell^p$ 上的算子范数，并判断哪些条件保证算子范数能在一个单位向量上达到；分别讨论 $p<\infty$ 与 $p=\infty$。
\end{enumerate}

\chapter{Hilbert 空间与正交展开}
\label{chap:fa-hilbert}

内积把长度、角度与正交性放在同一结构中。完备内积空间保留了有限维最小二乘的几何，同时允许用无穷多个正交坐标描述一个向量。

\section{内积及其诱导范数}
内积不仅衡量长度，还能测量一个向量沿另一个方向的分量。在实欧氏空间中，$\langle x,e\rangle$ 是 $x$ 在单位方向 $e$ 上的有向投影长度；在函数空间中，$\int f\overline g$ 扮演同样的角色。复数情形需要共轭，才能保证 $\langle x,x\rangle$ 为非负实数。

\begin{definition}
线性空间 $H$ 上的\keyterm{内积（inner product）}$\langle x,y\rangle$ 对第一个变量线性，满足 $\langle x,y\rangle=\overline{\langle y,x\rangle}$，且 $\langle x,x\rangle\geq0$，等号仅当 $x=0$。由内积定义的范数为 $\|x\|=\sqrt{\langle x,x\rangle}$；此范数完备时称 $H$ 为\keyterm{Hilbert 空间（Hilbert space）}。
\end{definition}
\begin{proposition}[Cauchy--Schwarz 不等式]
$|\langle x,y\rangle|\leq\|x\|\|y\|$，且等号成立当且仅当 $x,y$ 线性相关。
\end{proposition}
\begin{proof}
当 $y\ne0$ 时，展开 $\|x-\langle x,y\rangle y/\|y\|^2\|^2\geq0$，即得 $\|x\|^2-|\langle x,y\rangle|^2/\|y\|^2\geq0$。等号等价于该差向量为零；$y=0$ 的情形直接成立。
\end{proof}
这一不等式给出范数的三角不等式。内积还满足平行四边形恒等式
\[
 \|x+y\|^2+\|x-y\|^2=2\|x\|^2+2\|y\|^2.
\]
在复空间中可由范数恢复内积：
\[
 \langle x,y\rangle=\tfrac14\bigl(\|x+y\|^2-\|x-y\|^2
 +\mathrm i\|x+\mathrm iy\|^2-\mathrm i\|x-\mathrm iy\|^2\bigr).
\]
这是极化恒等式（polarization identity）。并非任意范数都来自内积；例如 $\mathbb R^2$ 的最大值范数不满足平行四边形恒等式。$\ell^2$ 与 $L^2(\mu)$ 的内积分别为 $\sum_jx_j\overline{y_j}$ 与 $\int f\overline g\,\mathrm d\mu$，其完备性已在上一部分证明。

上述反例可用两个坐标向量核查。取 $x=(1,0)$、$y=(0,1)$，最大值范数下 $\|x+y\|^2+\|x-y\|^2=2$，而 $2\|x\|^2+2\|y\|^2=4$。因此不能仅凭一个空间有范数，就在其中使用勾股定理或正交投影。相反，内积诱导范数的三角不等式来自
\[
 \|x+y\|^2=\|x\|^2+2\operatorname{Re}\langle x,y\rangle+\|y\|^2
 \leq(\|x\|+\|y\|)^2.
\]
这一推导也说明 Cauchy--Schwarz 是几何结构与范数估计之间的桥梁。

\section{最佳逼近与正交投影}
集合 $C$ 凸，是指 $u,v\in C$ 及 $0\leq t\leq1$ 时 $(1-t)u+tv\in C$。记 $x\perp y$ 表示 $\langle x,y\rangle=0$；对 $M\subset H$，其正交补为 $M^\perp=\{x\in H:\langle x,m\rangle=0\text{ 对所有 }m\in M\}$。这里的正交补由 $H$ 中的向量组成，与一般对偶空间中的零化子通过 Riesz 表示相联系。
\begin{theorem}[投影定理]\label{thm:fa-projection}
设 $C$ 是 Hilbert 空间 $H$ 的非空闭凸子集。对每个 $x\in H$，存在唯一 $p\in C$ 使 $\|x-p\|=\inf_{z\in C}\|x-z\|$。该点由
\[
 \operatorname{Re}\langle x-p,z-p\rangle\leq0\quad(z\in C)
\]
刻画。若 $C=M$ 是闭线性子空间，则 $x-p\perp M$，从而 $H=M\oplus M^\perp$。
\end{theorem}
\begin{proof}
记距离下确界为 $d$，取 $z_n\in C$ 使 $\|x-z_n\|\to d$。由凸性，中点在 $C$，平行四边形恒等式给出
\[
 \|z_n-z_m\|^2\leq2\|x-z_n\|^2+2\|x-z_m\|^2-4d^2\to0.
\]
完备性与闭性给出 $z_n\to p\in C$，且 $\|x-p\|=d$。同一不等式用于两个极小点即得唯一性。将 $p+t(z-p)$ 代入最小性，展开平方、除以 $t>0$ 并令 $t\downarrow0$，得到所述不等式；反之展开 $\|x-z\|^2$ 即知它保证最小性。

若 $C=M$，可令 $z=p+tv$，其中 $t$ 取正负实数，复空间中还可取纯虚数，故 $\langle x-p,v\rangle=0$。分解唯一，因为 $M\cap M^\perp=\{0\}$。
\end{proof}
记最近点为 $P_Mx$。由分解唯一性，$P_M$ 线性；由勾股恒等式 $\|x\|^2=\|P_Mx\|^2+\|x-P_Mx\|^2$，$\|P_M\|\leq1$，当 $M\ne\{0\}$ 时等号成立。闭性不能删除：若 $M$ 稠密而不闭，则 $x\notin M$ 到 $M$ 的距离为零，却无最近点。

投影定理的证明没有假定有界序列有范数收敛子列。它利用凸性允许取中点，再用平行四边形恒等式直接证明极小化序列是 Cauchy 序列。这是无限维证明中常见的替代办法。凸性也不能省略：在实直线上，闭集 $\{-1,1\}$ 到 $0$ 有两个最近点。

对一般闭凸集，记最近点映射为 $P_C$。它未必线性，但满足 $\|P_Cx-P_Cy\|\leq\|x-y\|$。事实上，令 $p=P_Cx,q=P_Cy$，分别在变分不等式中取候选点 $q,p$，相加整理得
\[
 \|p-q\|^2\leq\operatorname{Re}\langle x-y,p-q\rangle
 \leq\|x-y\|\|p-q\|.
\]
当 $p\ne q$ 时除以 $\|p-q\|$ 即得结论，$p=q$ 时直接成立。例如 $\mathbb R$ 到 $[0,\infty)$ 的投影为 $x\mapsto\max(x,0)$，它连续且不放大距离，却不是线性映射。

\section{Riesz 表示定理}
\begin{theorem}\label{thm:fa-riesz-representation}
对每个 $f\in H^*$，存在唯一 $y\in H$ 使 $f(x)=\langle x,y\rangle$ 对所有 $x\in H$ 成立，且 $\|f\|=\|y\|$。
\end{theorem}
\begin{proof}
$f=0$ 时取 $y=0$。否则 $M=\ker f$ 为真闭子空间。选 $z\notin M$，由投影定理，$z-P_Mz$ 是 $M^\perp$ 中的非零向量；将其归一化得到单位向量 $e$。若 $f(e)=0$，则 $e\in M\cap M^\perp=\{0\}$，矛盾，所以 $f(e)\ne0$。而 $x-f(x)e/f(e)\in M$，故 $\langle x,e\rangle=f(x)/f(e)$。取 $y=\overline{f(e)}e$ 即得表示。若两个向量给出同一泛函，将它们的差代入即知差为零。Cauchy--Schwarz 给出范数上界，取 $x=y/\|y\|$ 得到等号。
\end{proof}
复数情形中，$y\mapsto\langle\,\cdot\,,y\rangle$ 是共轭线性等距对应；它不是复线性映射。内积的变量约定将影响伴随及 Fourier 系数的写法。

这个定理把任意连续线性测量变成与一个固定向量的内积。比如在 $L^2(0,1)$ 中，$f(u)=\int_0^1t u(t)\,\mathrm dt$ 的表示向量是 $y(t)=t$，所以 $\|f\|=\|t\|_2=1/\sqrt3$，而不是粗略估计得到的一。复数情形中，若写成 $f(u)=\int a(t)u(t)\,\mathrm dt$，表示向量应为 $\overline a$；这与前章不带共轭的对偶配对约定一致。

\section{正交系与 Fourier 展开}
标准正交系 $(e_j)$ 满足 $\langle e_j,e_k\rangle=\delta_{jk}$，其中 $\delta_{jk}$ 在 $j=k$ 时为一，否则为零。对有限部分，正交投影为 $P_Nx=\sum_{j=1}^N\langle x,e_j\rangle e_j$，勾股恒等式给出 Bessel 不等式
\[
 \sum_{j=1}^{\infty}|\langle x,e_j\rangle|^2\leq\|x\|^2.
\]
因此 $P_Nx$ 为 Cauchy 序列，其极限是 $x$ 在 $\overline{\operatorname{span}\{e_j\}}$ 上的正交投影。若此闭包为 $H$，称 $(e_j)$ 为标准正交基（orthonormal basis），并有
\begin{equation}
 x=\sum_{j=1}^{\infty}\langle x,e_j\rangle e_j,
 \qquad \|x\|^2=\sum_{j=1}^{\infty}|\langle x,e_j\rangle|^2.
 \label{eq:fa-parseval}
\end{equation}
后一等式称为 Parseval 恒等式。可分 Hilbert 空间取可数稠密集，逐次进行 Gram--Schmidt 正交化并略去相关向量，即得到有限或可数正交基。Hamel 基只允许有限线性组合，不能与这里允许范数收敛级数的正交基混同。

\begin{proposition}[平方可和系数与向量的对应]\label{prop:fa-riesz-fischer}
若 $(e_j)_{j\geq1}$ 是 Hilbert 空间中的标准正交系，$(a_j)\in\ell^2$，则 $\sum_ja_je_j$ 在 $H$ 中收敛到某个 $x$，且 $\langle x,e_j\rangle=a_j$、$\|x\|^2=\sum_j|a_j|^2$。当该正交系为基时，这给出 $H$ 与 $\ell^2$ 的线性等距对应。
\end{proposition}
\begin{proof}
若 $s_N=\sum_{j\leq N}a_je_j$，则对 $N>M$，$\|s_N-s_M\|^2=\sum_{M<j\leq N}|a_j|^2\to0$。完备性给出 $s_N\to x$。内积连续，故逐个系数取极限得到 $\langle x,e_j\rangle=a_j$；范数连续给出范数等式。正交系为基时，每个 $x$ 都由自己的系数展开，因而对应同时单射、满射。
\end{proof}

标准正交系是否完备，也可以用“是否遗漏了一个正交方向”判断：它是正交基当且仅当 $\langle x,e_j\rangle=0$ 对所有 $j$ 成立只能推出 $x=0$。这是因为其闭线性张成空间的正交补恰由这些 $x$ 组成，再应用投影定理即可。Bessel 不等式对任何正交系成立，Parseval 等式则要求没有遗漏的方向。

在 $L^2(0,2\pi)$ 中，$e_k(t)=(2\pi)^{-1/2}e^{\mathrm ikt}$（$k\in\mathbb Z$）构成标准正交基；完备性是 Fourier 分析中的三角多项式稠密定理，这里引用该结果，参见\cite{teschl-ra}的 Fourier 级数部分。式\eqref{eq:fa-parseval}保证 $L^2$ 展开，不自动保证逐点或一致收敛。

\section{伴随与特殊算子}
对 $T\in\mathcal B(H)$ 和固定 $y$，泛函 $x\mapsto\langle Tx,y\rangle$ 连续。Riesz 表示给出唯一 $T^*y$ 满足
\[
 \langle Tx,y\rangle=\langle x,T^*y\rangle.
\]
由唯一性可知 $T^*$ 线性，且 $\|T^*y\|\leq\|T\|\|y\|$。逐一在内积中检验可得
\[
 (ST)^*=T^*S^*,\quad (T^*)^*=T,\quad
 (\alpha T)^*=\overline\alpha T^*,\quad\|T^*\|=\|T\|.
\]
此外，$\|Tx\|^2=\langle x,T^*Tx\rangle\leq\|T^*T\|\|x\|^2$，结合次乘性得到 $\|T^*T\|=\|T\|^2$。由伴随定义，$T^*y=0$ 当且仅当 $\langle Tx,y\rangle=0$ 对所有 $x$ 成立，故 $\ker T^*=(\operatorname{Ran}T)^\perp$；再取正交补，得到 $\overline{\operatorname{Ran}T}=(\ker T^*)^\perp$。

$T=T^*$ 时称自伴（self-adjoint）；$T^*T=TT^*$ 时称正规（normal）；$T^*T=TT^*=I$ 时称酉（unitary）；自伴且 $\langle Tx,x\rangle\geq0$ 时称正算子（positive operator）。正交投影恰是自伴幂等算子：正交分解直接给出 $P_M^*=P_M$；反之 $P^2=P=P^*$ 保证 $Px\perp x-Px$，且 $\operatorname{Ran}P=\ker(I-P)$ 闭。

\section{例题与习题}
\begin{example}[最佳一次逼近]
在 $L^2(0,1)$ 中，$t^2$ 的最佳一次多项式逼近是 $t-1/6$。
\end{example}
\begin{solve}
设投影为 $a+bt$。误差与 $1,t$ 正交，因此
\[
 a+\tfrac12b=\tfrac13,\qquad \tfrac12a+\tfrac13b=\tfrac14.
\]
解得 $a=-1/6$、$b=1$。这两个方程就是最小二乘的正规方程；投影定理同时保证所得解为唯一全局最佳逼近。
\end{solve}

更一般地，若 $\phi_1,\ldots,\phi_m$ 线性无关，令 $M=\operatorname{span}\{\phi_1,\ldots,\phi_m\}$，写 $P_Mx=\sum_jc_j\phi_j$，则正交条件给出
\[
 \sum_{j=1}^m\langle\phi_j,\phi_i\rangle c_j=\langle x,\phi_i\rangle
 \quad(1\leq i\leq m).
\]
矩阵 $G_{ij}=\langle\phi_j,\phi_i\rangle$ 称为 Gram 矩阵。对非零系数向量 $c$，$c^*Gc=\|\sum_jc_j\phi_j\|^2>0$，故矩阵正定、方程组唯一可解。若基标准正交，$G=I$，系数立即变成 Fourier 系数。这样，函数逼近中的投影问题被转化为可以实际求解的有限维线性方程组。
\begin{enumerate}
 \item\label{exer:fa-shift-adjoint} 证明 $\ell^2$ 中左右移位互为伴随，右移保持范数但不是酉算子。
 \item 证明无限标准正交序列的任意两项距离为 $\sqrt2$。
 \item\label{exer:fa-double-perp} 证明 $M^{\perp\perp}=\overline M$，其中 $M$ 是线性子空间。
 \item 证明 $\ker T^*=(\operatorname{Ran}T)^\perp$，并推得 $\overline{\operatorname{Ran}T}=(\ker T^*)^\perp$。
\end{enumerate}

\chapter{Hahn--Banach 定理与凸集分离}
\label{chap:fa-hahn-banach}

连续线性泛函可以看作对向量的标量测量。Hahn--Banach 定理说明，在子空间上已知的线性测量能够保持原有控制而延拓，从而保证对偶空间足够丰富。

\section{受次线性泛函控制的延拓}
实线性空间上的函数 $p:X\to\mathbb R$ 称为次线性泛函（sublinear functional），若 $p(x+y)\leq p(x)+p(y)$ 且 $p(tx)=tp(x)$ 对 $t\geq0$ 成立。

范数是次线性泛函，但次线性泛函可以不对称，也可以取负值。例如 $\mathbb R$ 上的 $p(t)=\max(t,2t)$ 满足上述条件；$p(1)=2$、$p(-1)=-1$，所以不能把 $p$ 当作绝对值使用。定理中的 $f\leq p$ 是逐点的单侧控制，只有在 $p(x)=C\|x\|$ 这样的对称情形下，才自然转化为 $|f(x)|\leq C\|x\|$。
\begin{theorem}[实 Hahn--Banach 定理]\label{thm:fa-hahn-banach-real}
设 $M\subset X$ 为实线性子空间，$p$ 次线性，$f:M\to\mathbb R$ 线性且 $f(x)\leq p(x)$。则存在全空间线性泛函 $F$，使 $F|_M=f$ 且 $F\leq p$。
\end{theorem}
\begin{proof}
先增加一个方向 $z\notin M$，令 $g(x+tz)=f(x)+tc$。当 $t>0$ 或 $t<0$ 时，要求 $g\leq p$ 分别等价于
\[
 \sup_{x\in M}\{f(x)-p(x-z)\}\leq c
 \leq\inf_{y\in M}\{p(y+z)-f(y)\}.
\]
左侧不超过右侧，因为
\[
 f(x)+f(y)=f(x+y)\leq p(x+y)\leq p(x-z)+p(y+z).
\]
两端分别具有有限上界和有限下界，故可选 $c$，完成一步延拓。

这里唯一未知的是新方向上的数值 $c=g(z)$。例如 $t>0$ 时，将 $f(x)+tc\leq p(x+tz)$ 除以 $t$，得到 $c\leq p(x/t+z)-f(x/t)$；$x/t$ 遍历 $M$，就得到右端所有上界。$t=-s<0$ 时，同样整理得到 $c\geq f(x/s)-p(x/s-z)$。只要每个下界都不超过每个上界，就能在中间选择一个实数；这正是次线性条件在证明中完成的工作。

将所有受 $p$ 控制、延拓 $f$ 的线性泛函按“定义域包含且数值一致”排序。任一链上的泛函之并仍是受控线性泛函，因而是链的上界。Zorn 引理（附录\ref{app:fa-topology}）给出极大延拓。若其定义域不是 $X$，一步延拓又能扩大定义域，与极大性矛盾。
\end{proof}

\section{保范延拓与复数情形}
\begin{corollary}\label{cor:fa-norm-preserving-extension}
若 $M$ 是实或复赋范空间 $X$ 的线性子空间，每个 $f\in M^*$ 都有延拓 $F\in X^*$，满足 $\|F\|=\|f\|$。$M$ 无须闭，$X$ 无须完备。
\end{corollary}
\begin{proof}
实情形取 $p(x)=\|f\|\|x\|$。所得 $F\leq p$ 同时应用于 $-x$，得 $|F(x)|\leq\|f\|\|x\|$；限制到 $M$ 又给出范数反向估计。

复情形先将 $\operatorname{Re}f$ 作为实线性泛函延拓为 $u$，再令 $F(x)=u(x)-\mathrm i u(\mathrm ix)$。由实线性可检验 $F(\mathrm ix)=\mathrm iF(x)$，故 $F$ 复线性，且 $F|_M=f$。选择 $|\alpha|=1$ 使 $\alpha F(x)=|F(x)|$，便有 $|F(x)|=u(\alpha x)\leq\|f\|\|x\|$。
\end{proof}
这一结论针对标量值泛函；它不保证任意向量值有界算子都能保范延拓。

从线性代数看，把一个子空间上的泛函延拓出去并不难：补足一组基，再任意指定新基向量的值即可。Hahn--Banach 定理的内容在于这些值不能任意选择，必须同时服从全空间上的同一个范数界。另一方面，若子空间稠密，连续延拓由极限唯一决定；一般不稠密的子空间没有这种唯一性，保范延拓可能有多个。

\section{支撑泛函与范数的对偶刻画}
\begin{proposition}\label{prop:fa-norming-functional}
对每个非零 $x\in X$，存在 $f\in X^*$ 使 $\|f\|=1$ 且 $f(x)=\|x\|$。因此
\begin{equation}
 \|x\|=\sup_{\|f\|\leq1}|f(x)|.
 \label{eq:fa-dual-norm}
\end{equation}
\end{proposition}
\begin{proof}
在一维子空间上定义 $f_0(\alpha x)=\alpha\|x\|$，其范数为一。由推论\ref{cor:fa-norm-preserving-extension}延拓即可。上确界不超过 $\|x\|$ 来自对偶范数定义，上述泛函给出相反不等式。
\end{proof}
由式\eqref{eq:fa-dual-norm}，典范映射 $J:X\to X^{**}$ 为线性等距嵌入。此外，对任意 $T\in\mathcal B(X,Y)$，
\[
 \|T\|=\sup_{\|x\|\leq1}\sup_{\|g\|\leq1}|g(Tx)|=\|T'\|.
\]
注意，此处为给定向量选择一个达到其范数的泛函，并不声称每个给定泛函都在单位球上达到自己的范数。

二者的量词次序不同。式\eqref{eq:fa-dual-norm}先固定 $x$，再允许选择依赖于 $x$ 的 $f$；“每个泛函达到范数”则先固定 $f$，再寻找某个单位向量。在 $\ell^1$ 上取 $f(x)=\sum_j a_jx_j$、$a_j=1-1/(j+1)$，则 $\|f\|=\sup_ja_j=1$。然而对任意非零 $x$，有 $|f(x)|\leq\sum_ja_j|x_j|<\|x\|_1$，所以它在闭单位球上不能达到模为一的值。后面将看到，自反空间具有更强的范数达到性质。

\section{凸集的分离}
\begin{theorem}[点与闭凸集的严格分离]\label{thm:fa-separation}
设 $C$ 是实赋范空间 $X$ 的非空闭凸子集，$x_0\notin C$。则存在 $f\in X^*$ 及实数 $a$，使
\[
 \sup_{c\in C}f(c)<a<f(x_0).
\]
在复空间中用 $\operatorname{Re}f$ 代替 $f$，结论仍成立。
\end{theorem}
\begin{proof}
先考虑含零的开凸集 $U$。其 Minkowski 泛函定义为 $p_U(x)=\inf\{t>0:x\in tU\}$。由于 $U$ 含一个零点邻域，它吸收每个向量，故 $p_U$ 有限；由凸性验证次线性，且 $U=\{x:p_U(x)<1\}$。若 $z\notin U$，则 $p_U(z)\geq1$。在 $\mathbb Rz$ 上令 $h(tz)=tp_U(z)$；当 $t<0$ 时，由 $p_U(z)+p_U(-z)\geq0$ 仍有 $h(tz)\leq p_U(tz)$。Hahn--Banach 定理给出 $h$ 的延拓 $f\leq p_U$。若 $B(0,r)\subset U$，则 $|f(x)|\leq\|x\|/r$，所以 $f$ 连续，且 $f(u)<1\leq f(z)$。

回到原问题，选 $r>0$ 使 $x_0\notin C+B(0,r)$，再将此开凸集平移为含零的开凸集。上一步给出非零连续泛函，使 $f(c+v)<f(x_0)$ 对 $c\in C$、$\|v\|<r$ 成立。对 $v$ 取上确界，得 $f(c)+r\|f\|\leq f(x_0)$，从而具有严格正间隙。复空间先在底层实空间分离，再以 $F(x)=f(x)-\mathrm if(\mathrm ix)$ 恢复复线性泛函。
\end{proof}
两个不交非空凸集若其中一个开，同样可通过对其差集分离零点，得到非零连续实线性泛函的非严格分离。没有开性或正距离等条件时，不能任意要求统一的严格间隙。

\begin{example}[分离定理的二维图景]
在欧氏平面中，令 $C=\{(u,v):(u-2)^2+v^2\leq1\}$，$x_0=(0,0)$。取 $f(u,v)=-u$，则 $\sup_C f=-1< -1/2<f(x_0)=0$。
\end{example}
\begin{solve}
圆盘中横坐标 $u\geq1$，所以 $f\leq-1$；直线 $f=-1/2$，也就是 $u=1/2$，位于圆盘和原点之间。一般分离定理把这条直线推广为超平面 $\{x:f(x)=a\}$，并保证无限维中也能找到一个连续线性测量，把整个闭凸集与给定的外部点区分开。定理并不要求闭凸集有界。
\end{solve}

分离证明中使用的 Minkowski 泛函也可以视为相对于凸集的“尺度”。为看清其性质，设 $U$ 开、凸且包含零。若 $a>p_U(x)$、$b>p_U(y)$，则 $x/a,y/b\in U$，由凸性
\[
 \frac{x+y}{a+b}=\frac a{a+b}\frac xa+\frac b{a+b}\frac yb\in U,
\]
故 $p_U(x+y)\leq a+b$，让 $a,b$ 分别下降到下确界即可得次可加性。正齐次性由定义中的缩放得到。若 $p_U(x)<1$，可选 $t<1$ 使 $x\in tU\subset U$；若 $x\in U$，由开放性，可把 $x$ 沿射线稍微放大后仍留在 $U$，故 $p_U(x)<1$。这解释了 $U=\{p_U<1\}$，也解释了为何证明中要先把凸集扩大为一个开集，再使用代数形式的 Hahn--Banach 定理。

\section{距离、闭包与零化子}
\begin{proposition}\label{prop:fa-distance-dual}
若 $M$ 是闭线性子空间，则
\[
 \operatorname{dist}(x,M)=
 \sup\{|f(x)|:f\in X^*,\ \|f\|\leq1,\ f|_M=0\}.
\]
\end{proposition}
\begin{proof}
对右侧每个 $f$，$|f(x)|=|f(x-m)|\leq\|x-m\|$，故上确界不超过距离。反向在商空间 $X/M$ 的向量 $x+M$ 上应用命题\ref{prop:fa-norming-functional}，再与商映射复合，即得到所需泛函；$x\in M$ 时两侧均为零。
\end{proof}
于是对任意线性子空间 $M$，向量 $x$ 属于 $\overline M$ 当且仅当所有在 $M$ 上消失的连续泛函也在 $x$ 上消失。特别地，
\[
 \overline{\operatorname{Ran}T}=Y\quad\Longleftrightarrow\quad\ker T'=\{0\}.
\]

\section{例题与习题}
\begin{example}[支撑泛函可以不唯一]
在 $\mathbb R^2$ 的 $\ell^1$ 范数下，$x=(1,0)$ 的范数为一，而 $f_a(u,v)=u+av$ 对每个 $|a|\leq1$ 都满足 $\|f_a\|=1$、$f_a(x)=1$。
\end{example}
\begin{solve}
$|u+av|\leq|u|+|v|$ 给出范数上界；在 $(1,0)$ 处达到一。这说明延拓存在性不包含唯一性，支撑泛函的几何性质还取决于范数。
\end{solve}
\begin{enumerate}
 \item 证明有限维真子空间总可以被一个非零连续泛函零化。
 \item 在欧氏平面中分离原点与闭球 $\{x:\|x-(2,0)\|_2\leq1\}$。
 \item 证明赋范空间中任意两个不同向量可被连续线性泛函区分。
 \item 对 $a\in\ell^\infty$、$a\ne0$，证明 $f(x)=\sum_ja_jx_j$ 在 $\ell^1$ 的闭单位球上达到范数，当且仅当某个坐标满足 $|a_j|=\|a\|_\infty$。
\end{enumerate}

\chapter{Baire 纲定理与算子基本定理}
\label{chap:fa-baire}

完备性不仅保证给定 Cauchy 序列收敛，还能把逐点信息转化为统一估计。本章的三个基本工具分别回答：一族算子是否统一有界，满射是否把开集映成开集，以及闭图像是否足以保证连续。

\section{Baire 纲定理}
集合 $A$ 称为无处稠密（nowhere dense），若 $\overline A$ 的内部为空；可数个无处稠密集之并称为第一纲集（meagre set）。
\begin{theorem}[Baire 纲定理]\label{thm:fa-baire}
非空完备度量空间中，可数个稠密开集的交仍然稠密。因此，该空间不能表示为可数个无处稠密集之并。
\end{theorem}
\begin{proof}
设 $U_n$ 稠密且开，$V$ 为任意非空开集。在 $V\cap U_1$ 中选择正半径闭球 $\overline B(x_1,r_1)$；再在 $B(x_1,r_1)\cap U_2$ 中选闭球，以此递归，使
\[
 \overline B(x_{n+1},r_{n+1})\subset B(x_n,r_n)\cap U_{n+1},
 \qquad0<r_n\leq2^{-n}.
\]
球心构成 Cauchy 序列，完备性给出极限 $x$。每个闭球均包含足够后面的球心，故 $x$ 属于所有闭球，从而 $x\in V\cap\bigcap_nU_n$。由于 $V$ 任意，交集稠密。对无处稠密集闭包的补集应用这一结论，即得后一断言。
\end{proof}
“第一纲”是拓扑概念，“零测”是测度概念，二者不能替代。

后面最常用的是它的等价形式：若非空完备度量空间被可数个闭集覆盖，则至少一个闭集有非空内部。这里一定要有“闭”与“可数”。例如实数可以写成不可数个单点闭集的并，每个单点内部为空；这不违反 Baire 定理。在不完备空间 $\mathbb Q$ 中，所有有理数的单点组成可数个无处稠密闭集，却覆盖了整个空间。Baire 定理把完备性转化为“某个估计必须在一整块邻域内同时成立”的工具。

\section{一致有界原理}
\begin{theorem}[Banach--Steinhaus]\label{thm:fa-uniform-boundedness}
设 $X$ 为 Banach 空间，$Y$ 为赋范空间，$\mathcal T\subset\mathcal B(X,Y)$。若对每个 $x\in X$，$\sup_{T\in\mathcal T}\|Tx\|<\infty$，则 $\sup_{T\in\mathcal T}\|T\|<\infty$。
\end{theorem}
\begin{proof}
令 $E_n=\{x:\sup_{T\in\mathcal T}\|Tx\|\leq n\}$，每个 $E_n$ 闭且 $X=\bigcup_nE_n$。Baire 定理保证某个 $E_N$ 含球 $B(x_0,r)$。当 $\|h\|<r$ 时，$x_0,x_0+h\in E_N$，所以 $\|Th\|\leq2N$ 对所有 $T$ 成立。取 $h=(r/2)x$、$\|x\|\leq1$，即得 $\|T\|\leq4N/r$。
\end{proof}
这里逐点有界的上界原本可随 $x$ 改变，结论却产生一个控制所有单位向量的常数。值域 $Y$ 不需要完备。

更明确地，假设是对每个 $x$ 存在 $C_x$，使 $\|Tx\|\leq C_x$ 对所有 $T\in\mathcal T$ 成立；结论是存在一个 $C$，使 $\|Tx\|\leq C\|x\|$ 对所有 $T,x$ 同时成立。不能把“对每个 $T$ 存在 $C_T$”误当作假设，因为那只是各算子分别有界。例如 $T_n=nI$ 每个都有界，但对任何非零输入都不逐点有界。

\begin{example}[缺少完备性时]
在 $c_{00}$ 的上确界范数下，令 $T_nx=nx_n$。对每个 $x\in c_{00}$，$\sup_n|T_nx|<\infty$，但 $\|T_n\|=n$。
\end{example}
每个向量仅有有限个非零坐标，因此逐点有界；取 $x=e_n$ 得到算子范数。定义域不完备正是定理不能应用的原因。

\section{开映射与有界逆}
\begin{theorem}[开映射定理]\label{thm:fa-open-mapping}
Banach 空间间的有界线性满射 $T:X\to Y$ 将开集映成开集。因而有界线性双射的逆也有界。
\end{theorem}
\begin{proof}
用 $B_X,B_Y$ 表示闭单位球。满射性给出 $Y=\bigcup_{n\geq1}T(nB_X)$，因而这些像的闭包覆盖 $Y$。由 $Y$ 的完备性和 Baire 定理，某个 $\overline{T(nB_X)}$ 包含开球 $B(y_0,r)$。若 $\|y\|<r$，则 $y_0+y/2$ 与 $y_0-y/2$ 都在该球中。分别用 $T(nB_X)$ 中的点逼近它们，作差得到 $y\in\overline{T(2nB_X)}$。缩放后即存在 $\delta>0$ 满足
\[
 \delta B_Y\subset\overline{T(B_X)}.
\]
必要时缩小 $\delta$，可把开球中的包含改写为这一闭球包含。

现在消去闭包。对 $\|y\|\leq\delta/2$，由缩放后的包含选择 $\|x_1\|\leq1/2$ 使 $\|y-Tx_1\|\leq\delta/4$。递归选 $\|x_k\|\leq2^{-k}$，使第 $k$ 次残差不超过 $\delta2^{-k-1}$。级数 $x=\sum_kx_k$ 在 $X$ 中收敛，$\|x\|\leq1$，连续性给出 $Tx=y$。于是 $(\delta/2)B_Y\subset T(B_X)$。再缩放一次可用定义域中的开球容纳所需原像，故 $T$ 在零点开；平移给出处处开。双射的逆因此连续，由定理\ref{thm:fa-bounded-continuous}有界。
\end{proof}
满射性控制解的存在；单射性控制唯一性；有界逆定理则在 Banach 框架下保证解对数据连续依赖。这是三个不同的判断。

上面的证明也指出两个完备性假设用在不同位置：$Y$ 完备用于 Baire 定理，$X$ 完备用于把逐次修正 $x_1+x_2+\cdots$ 求和。第一次得到的闭包包含只表示每个小数据都有近似原像；第二步不断消除残差，才把近似可解提升为真正可解。不能把 $\overline{T(B_X)}$ 直接换成 $T(B_X)$，因为一般连续映射不保持闭集。

\begin{proposition}[闭值域与稳定下界]\label{prop:fa-closed-range-bound}
设 $X,Y$ 为 Banach 空间，$T\in\mathcal B(X,Y)$ 单射。则 $\operatorname{Ran}T$ 闭当且仅当存在 $c>0$ 使 $\|Tx\|\geq c\|x\|$ 对所有 $x\in X$ 成立。
\end{proposition}
\begin{proof}
若值域闭，它本身为 Banach 空间；把 $T$ 看作 $X\to\operatorname{Ran}T$ 的有界双射，有界逆定理给出 $\|x\|\leq C\|Tx\|$。反之，若 $Tx_n\to y$，下界使 $x_n$ 为 Cauchy 序列；由 $X$ 完备得 $x_n\to x$，连续性给出 $y=Tx$，故值域闭。
\end{proof}
这个下界给出实际的稳定性：若 $Tx=y$、$T\widetilde x=\widetilde y$，则 $\|x-\widetilde x\|\leq c^{-1}\|y-\widetilde y\|$。它只控制值域内的数据，还不能单独证明满射；满射需要另外验证。

\section{闭图像定理}
\begin{theorem}\label{thm:fa-closed-graph}
设 $X,Y$ 为 Banach 空间，线性映射 $T:X\to Y$ 定义在整个 $X$ 上。则 $T$ 有界当且仅当图像 $G(T)=\{(x,Tx):x\in X\}$ 在 $X\times Y$ 中闭。
\end{theorem}
\begin{proof}
有界性保证 $x_n\to x$、$Tx_n\to y$ 时 $y=Tx$，故图像闭。反之，$X\times Y$ 在范数 $\|(x,y)\|=\|x\|+\|y\|$ 下完备，闭子空间 $G(T)$ 也完备。投影 $P:G(T)\to X$，$P(x,Tx)=x$，为有界线性双射。其逆有界，故 $\|x\|+\|Tx\|\leq C\|x\|$，得到 $T$ 有界。
\end{proof}
实际应用时，不必事先证明 $Tx_n$ 收敛；只需证明“若 $x_n\to x$ 且 $Tx_n\to y$，则 $y=Tx$”。不过，定义域必须是整个 Banach 空间；真子空间上定义的闭微分算子并不因此有界。

\begin{example}[利用闭图像验证连续性]
设 $X$ 为 Banach 空间，线性映射 $T:X\to\ell^2$ 处处有定义，并且每个坐标泛函 $x\mapsto(Tx)_j$ 都连续。则 $T$ 是有界算子。
\end{example}
\begin{solve}
设 $x_n\to x$ 且 $Tx_n\to y$ 在 $\ell^2$ 范数下成立。对每个固定 $j$，坐标泛函的连续性给出 $(Tx_n)_j\to(Tx)_j$，而 $\ell^2$ 范数收敛又给出 $(Tx_n)_j\to y_j$。因此 $y_j=(Tx)_j$ 对所有 $j$ 成立，也就是 $y=Tx$。图像闭，由闭图像定理得到统一估计 $\|Tx\|_2\leq C\|x\|$。假设中“$T$ 处处取值于 $\ell^2$”不可忽略；只给出无穷多个连续坐标公式，未必保证这些输出组成平方可和数列。
\end{solve}

\section{如何选择定理}
对一族算子从逐点估计求统一范数界，用一致有界原理；对连续线性满射求开性，用开映射定理；对双射估计逆，用有界逆定理；对难以直接估计的处处定义线性算子，可验证闭图像。除一致有界原理的值域外，这些定理中出现的空间均须满足相应完备性假设。

\begin{proposition}
若 $X,Y$ 为 Banach 空间，$T_n\in\mathcal B(X,Y)$，且每个 $x$ 上 $T_nx$ 收敛，则逐点极限 $Tx=\lim_nT_nx$ 为有界线性算子。
\end{proposition}
\begin{proof}
逐点收敛给出逐点有界，一致有界原理给出 $M=\sup_n\|T_n\|<\infty$。线性运算与极限相容，且 $\|Tx\|\leq M\|x\|$。
\end{proof}

\section{例题与习题}
\begin{example}[两个完备范数的等价性]
若 $X$ 上的两个范数均完备，且 $\|x\|_a\leq C\|x\|_b$，则二者等价。
\end{example}
\begin{solve}
恒等映射 $I:(X,\|\cdot\|_b)\to(X,\|\cdot\|_a)$ 为有界线性双射，故其逆有界，提供另一侧范数估计。若缺少完备性，这个结论会失败，例如连续函数空间上的积分范数与一致范数。
\end{solve}
\begin{enumerate}
 \item\label{exer:fa-hamel} 证明无限维 Banach 空间不可能有可数 Hamel 基。
 \item\label{exer:fa-dense-convergence} 若 $T_n$ 在某稠密线性子空间上逐点收敛，且 $\sup_n\|T_n\|<\infty$，证明在 Banach 值域中它在整个定义域上逐点收敛。
 \item 对例\ref{ex:fa-diagonal-range}的对角算子解释其逆在值域的诱导范数下为何无界，这与有界逆定理是否矛盾？
 \item 证明一个可数完备度量空间必有孤立点。
\end{enumerate}

\chapter{弱拓扑、弱星拓扑与收敛}
\label{chap:fa-weak}

无限维闭单位球通常不紧。一个自然的补救是只要求每个连续线性观测收敛，而不要求误差在所有观测上同时趋于零。这种收敛较弱，却与线性方程和凸优化相容。

\section{弱拓扑与弱星拓扑}
\begin{definition}
$X$ 的\keyterm{弱拓扑（weak topology）}$\sigma(X,X^*)$ 是使所有 $f\in X^*$ 连续的最弱拓扑。点 $x$ 的基本邻域为
\[
 U(x;f_1,\ldots,f_m;\varepsilon)=
 \{y\in X:|f_j(y-x)|<\varepsilon,\ 1\leq j\leq m\},
\]
其中 $m$ 有限且 $\varepsilon>0$。$X^*$ 上使所有取值映射 $f\mapsto f(x)$ 连续的最弱拓扑称为\keyterm{弱星拓扑（weak-star topology）}$\sigma(X^*,X)$。
\end{definition}
弱邻域只限制有限个泛函。范数小则这些观测都小，故弱拓扑弱于范数拓扑。命题\ref{prop:fa-norming-functional}保证不同点能被泛函区分，所以弱拓扑是 Hausdorff 的。

“较弱”指开集更少，于是收敛要求也更少。范数收敛要求 $\sup_{\|f\|\leq1}|f(x_n-x)|\to0$，弱收敛只要求对每个固定的 $f$，$|f(x_n-x)|\to0$。后者允许达到给定误差所需的指标 $N$ 随测试泛函改变，不能交换这里的极限与上确界。可把一个泛函理解为一个固定测量装置：每台固定装置最后都看不到误差，并不保证存在一个时刻，使所有装置同时看不到误差。

在 $\ell^2$ 中，只限制前 $m$ 个坐标的邻域 $\{x:|x_j|<\varepsilon,1\leq j\leq m\}$ 包含所有 $te_{m+1}$，$t\in\mathbb R$，因而范数无界。一般无限维空间的有限个泛函也有非零公共核：若公共核为零，映射 $x\mapsto(f_1(x),\ldots,f_m(x))$ 将是到有限维空间的线性单射，迫使原空间有限维。沿公共核中的方向可以任意放大而不改变这些测试值，所以弱邻域通常不能被理解为范数意义下的小球。

\section{序列收敛及其例子}
由基本邻域定义，
\begin{align}
 x_n\rightharpoonup x
 &\iff f(x_n)\to f(x)\quad\text{对每个 }f\in X^*,
 \label{eq:fa-weak-convergence}\\
 f_n\overset{*}{\rightharpoonup}f
 &\iff f_n(x)\to f(x)\quad\text{对每个 }x\in X.
 \label{eq:fa-weak-star-convergence}
\end{align}
$X^*$ 自身的弱拓扑是 $\sigma(X^*,X^{**})$，比弱星拓扑使用更多测试对象。一般弱拓扑不必可度量化，序列不能完整刻画其中的闭包和紧性；一般表述需要附录\ref{app:fa-topology}中的网。

\begin{example}\label{ex:fa-weak-unit-vectors}
$\ell^2$ 中 $e_n\rightharpoonup0$，但 $\|e_n\|_2=1$。把 $\ell^1$ 视为 $c_0^*$ 时，同样的 $e_n$ 弱星趋于零，却不在 $\ell^1$ 中弱趋于零。
\end{example}
\begin{solve}
$\ell^2$ 的任意泛函由一个平方可和数列表示，其第 $n$ 个坐标趋零，故 $f(e_n)\to0$。对 $c_0$ 的任意元素 $x$，$e_n(x)=x_n\to0$，得到弱星收敛。但 $\ell^1$ 上的连续泛函 $f(x)=\sum_jx_j$ 对所有 $e_n$ 取值一，排除弱收敛到零。
\end{solve}

这个例子中，$e_n$ 的质量没有缩小，只是向后面的坐标移动；这解释了范数恒为一与弱极限为零为何可以同时成立。对 $\ell^1$ 而言，同一个求和泛函会把全部坐标一起计入，因此能一直测到数值一。弱收敛不仅依赖向量的公式，还依赖所在空间的连续对偶。

\section{有界性与下半连续性}
\begin{proposition}\label{prop:fa-weak-bounds}
弱收敛序列范数有界，且 $x_n\rightharpoonup x$ 时
\[
 \|x\|\leq\liminf_n\|x_n\|.
\]
若 $X$ 为 Banach 空间，则 $X^*$ 中的弱星收敛序列也范数有界。
\end{proposition}
\begin{proof}
将 $Jx_n$ 视为完备空间 $X^*$ 上的连续泛函。弱收敛保证它们逐点有界，由一致有界原理及 $\|Jx_n\|=\|x_n\|$ 得到范数有界；这一步不需要 $X$ 本身完备。对任意 $\|f\|\leq1$，有 $|f(x)|=\lim_n|f(x_n)|\leq\liminf_n\|x_n\|$，再对 $f$ 取上确界。弱星情形直接把 $f_n$ 作为 $X$ 上的泛函使用一致有界原理，此时需要 $X$ 完备。
\end{proof}
有界线性算子弱到弱连续，因为对 $g\in Y^*$，$g(Tx_n)=(T'g)(x_n)\to(T'g)(x)$。因此线性约束能够在弱极限中保持。

\begin{proposition}[在 $\ell^2$ 中怎样验证弱收敛]\label{prop:fa-coordinate-weak}
设 $x^{(n)},x\in\ell^2$。则 $x^{(n)}\rightharpoonup x$ 当且仅当 $\sup_n\|x^{(n)}\|_2<\infty$，且对每个固定 $j$ 有 $x_j^{(n)}\to x_j$。
\end{proposition}
\begin{proof}
必要性由弱收敛的有界性及坐标泛函的连续性得到。反之，对任意 $y\in\ell^2$，把内积拆为前 $N$ 项与尾部。前 $N$ 项因逐坐标收敛趋于零；尾部满足
\[
 \left|\sum_{j>N}(x_j^{(n)}-x_j)\overline{y_j}\right|
 \leq(M+\|x\|_2)\left(\sum_{j>N}|y_j|^2\right)^{1/2},
 \qquad M=\sup_n\|x^{(n)}\|_2.
\]
先选 $N$ 使尾部对所有 $n$ 同时很小，再选 $n$ 控制有限部分，便得 $\langle x^{(n)}-x,y\rangle\to0$。Riesz 表示说明这已测试了全部连续泛函。
\end{proof}

不能删掉统一范数界。例如 $x^{(n)}=ne_n$ 的每个固定坐标最终为零，但取 $y=(1/j)\in\ell^2$，便有 $\langle x^{(n)},y\rangle=1$，所以它不弱趋于零。这也展示了证明中尾部估计的必要性。

在 Hilbert 空间中，弱收敛与范数收敛可以组合为强收敛：若 $x_n\rightharpoonup x$ 且 $\|x_n\|\to\|x\|$，则
\[
 \|x_n-x\|^2=\|x_n\|^2+\|x\|^2-2\operatorname{Re}\langle x_n,x\rangle\to0.
\]

\section{凸集的闭包与 Mazur 引理}
\begin{theorem}\label{thm:fa-convex-weak-closure}
赋范空间中，凸集的范数闭包与弱闭包相同。特别地，范数闭凸集弱闭。
\end{theorem}
\begin{proof}
由于弱拓扑较弱，范数闭包包含于弱闭包。若 $x$ 不在凸集 $C$ 的范数闭包中，定理\ref{thm:fa-separation}给出 $f\in X^*$ 和 $a$，使 $\sup_C\operatorname{Re}f<a<\operatorname{Re}f(x)$。集合 $\{y:\operatorname{Re}f(y)>a\}$ 是与 $C$ 不交的弱开邻域，所以 $x$ 也不在弱闭包中。
\end{proof}
集合 $A$ 的凸包 $\operatorname{conv}A$ 由所有有限凸组合 $\sum_{j=1}^ma_jx_j$ 组成，其中 $x_j\in A$、$a_j\geq0$、$\sum_ja_j=1$。
\begin{corollary}[Mazur 引理]\label{cor:fa-mazur}
若 $x_n\rightharpoonup x$，则存在每个尾部的有限凸组合 $y_n\in\operatorname{conv}\{x_k:k\geq n\}$，使 $\|y_n-x\|\to0$。
\end{corollary}
\begin{proof}
$x$ 属于每个尾部的弱闭包，因此属于该尾部凸包的弱闭包，也就属于范数闭包。选择 $y_n$ 使 $\|y_n-x\|<1/n$ 即可。
\end{proof}
这里使用凸组合，并不保证原序列存在强收敛子列。例\ref{ex:fa-weak-unit-vectors}的 $e_n$ 没有强收敛子列，但 $n^{-1}(e_1+\cdots+e_n)$ 的范数为 $n^{-1/2}$。

凸性是闭包定理的重要条件。在无限维 Hilbert 空间中，单位球面 $S=\{x:\|x\|=1\}$ 范数闭，却不是弱闭，因为一个标准正交序列位于 $S$ 中并弱趋于 $0\notin S$。与之相对，闭单位球既闭又凸，因此弱闭。对求极小值问题而言，约束 $\|x\|\leq1$ 可以在弱极限中保持，约束 $\|x\|=1$ 则可能丢失，这一差别会直接影响存在性论证。

\section{算子的不同收敛方式}
对 $T_n,T\in\mathcal B(X,Y)$，算子范数收敛要求 $\|T_n-T\|\to0$；\keyterm{强算子收敛（strong operator convergence）}要求对每个 $x$ 有 $\|T_nx-Tx\|\to0$；\keyterm{弱算子收敛（weak operator convergence）}只要求 $g(T_nx)\to g(Tx)$ 对每个 $x\in X,g\in Y^*$ 成立。三者依次蕴含。

在 $\ell^2$ 上令 $P_n$ 保留前 $n$ 个坐标，则 $P_nx\to x$ 对每个 $x$ 成立，而 $\|P_n-I\|=1$，所以强算子收敛不蕴含算子范数收敛。右移算子 $R$ 的幂满足 $R^nx\rightharpoonup0$：先对有限支撑测试向量验证，再用稠密性与 $\|R^n\|=1$ 延拓。因此 $R^n$ 弱算子趋零，却不强算子趋零。

\section{例题与习题}
\begin{enumerate}
 \item 若 $(x_n)$ 范数有界，且在 $X^*$ 的某个范数稠密子集上 $f(x_n)\to f(x)$，证明 $x_n\rightharpoonup x$。
 \item 证明有限维空间中的弱收敛与范数收敛等价。
 \item 对 $\ell^2$ 的单位向量，构造由第 $n$ 项开始的有限凸组合，使其范数趋零。
 \item 证明无限维赋范空间中任意基本弱零邻域都是范数无界的。提示：有限多个泛函的公共核非零。
\end{enumerate}

\chapter{弱紧性与自反空间}
\label{chap:fa-reflexive}

弱收敛的优势在于：某些不具有范数紧性的有界集合，在弱拓扑下恢复紧性。对偶空间中的一般结论是弱星紧性；要把它转化为原空间中的弱紧性，需要自反性。

\section{Banach--Alaoglu 定理}
先看一个不需要一般拓扑的模型：$\ell^2$ 中每个有界序列都有弱收敛子列。设 $\|x^{(n)}\|_2\leq M$，每个坐标序列有界。先选子列使第一坐标收敛，再从中选子列使第二坐标也收敛，依次进行；取这些嵌套子列的对角子列，使所有坐标同时收敛，记极限为 $x_j$。对每个有限 $N$，
\[
 \sum_{j=1}^N|x_j|^2=\lim_k\sum_{j=1}^N|x_j^{(n_k)}|^2\leq M^2.
\]
令 $N\to\infty$，可知 $x\in\ell^2$、$\|x\|_2\leq M$；再由命题\ref{prop:fa-coordinate-weak}得到弱收敛。一般弱紧性理论要把“逐坐标提取极限”的思想推广到由全部连续泛函提供的坐标，其中坐标集可能不可数。

\begin{theorem}\label{thm:fa-alaoglu}
任意赋范空间 $X$ 的对偶闭单位球 $B_{X^*}$ 在 $\sigma(X^*,X)$ 下紧。
\end{theorem}
\begin{proof}
将 $f\in B_{X^*}$ 视为坐标族 $(f(x))_{x\in X}$。每个坐标属于紧集 $D_x=\{z\in\mathbb K:|z|\leq\|x\|\}$，而乘积 $\prod_{x\in X}D_x$ 由 Tychonoff 定理紧。这个乘积中满足
\[
 a_{\alpha x+\beta y}=\alpha a_x+\beta a_y
 \quad(x,y\in X,\ \alpha,\beta\in\mathbb K)
\]
的坐标族构成闭集，因为每个等式都是有限坐标上的闭条件。它们恰好对应范数至多一的线性泛函。乘积拓扑在此闭集上的诱导拓扑正是弱星拓扑，所以 $B_{X^*}$ 紧。所用乘积紧性是附录\ref{app:fa-topology}中明确列出的拓扑输入。
\end{proof}
这一定理没有声称原空间 $X$ 的闭单位球弱紧，也没有在一般情形下声称每个有界泛函序列都有弱星收敛子列。

\section{可分空间中的弱星子列}
\begin{theorem}\label{thm:fa-weak-star-subsequence}
若 $X$ 可分，则 $X^*$ 中每个范数有界序列都有弱星收敛子列。
\end{theorem}
\begin{proof}
取可数稠密集 $\{x_j\}$。若 $\|f_n\|\leq M$，则对每个 $j$，标量列 $(f_n(x_j))$ 有界。逐次选子列并取对角子列，得到 $f_{n_k}(x_j)$ 对所有 $j$ 收敛。对任意 $x$，估计
\[
 |f_{n_k}(x)-f_{n_l}(x)|\leq2M\|x-x_j\|
       +|f_{n_k}(x_j)-f_{n_l}(x_j)|
\]
说明 $(f_{n_k}(x))$ 为 Cauchy 序列。令 $f(x)$ 为其极限，便得线性泛函及 $|f(x)|\leq M\|x\|$，所以 $f\in X^*$，子列弱星趋于 $f$。
\end{proof}
在 $B_{X^*}$ 上，还可用
\[
 d(f,g)=\sum_{j\geq1}2^{-j}\frac{|f(x_j)-g(x_j)|}{1+|f(x_j)-g(x_j)|}
\]
给弱星拓扑赋予度量。稠密性与统一范数界保证该度量收敛等价于逐点收敛。这里的有界集限制不可随意省略。

\section{自反空间与典范嵌入}
\begin{definition}
Banach 空间 $X$ 称为\keyterm{自反的（reflexive）}，若典范等距映射 $J:X\to X^{**}$，$(Jx)(f)=f(x)$，为满射。
\end{definition}
Hilbert 空间自反。实情形连续两次使用 Riesz 表示即可；复情形设 $R(y)=\langle\,\cdot\,,y\rangle$，对 $F\in H^{**}$，映射 $y\mapsto\overline{F(Ry)}$ 为连续线性泛函，故等于 $\langle y,z\rangle$。于是 $F(Ry)=\langle z,y\rangle=(Ry)(z)$，从而 $F=Jz$。

由定理\ref{thm:fa-sequence-dual}，$1<p<\infty$ 时 $\ell^p$ 自反；在本书采用的 $\sigma$ 有限测度空间上，$L^p$ 也由其对偶表示得出同一结论。另一方面，$c_0^{**}\cong\ell^\infty$，典范映射就是自然包含，因此 $c_0$ 不自反。自反性要求的是典范映射满射，而不只是抽象存在某个同构。

可以逐个追踪这个典范映射。一个 $x\in c_0$ 对 $a\in\ell^1=c_0^*$ 的作用是 $a\mapsto\sum_ja_jx_j$；而双对偶中还允许用任意有界数列 $b$ 定义 $a\mapsto\sum_ja_jb_j$。取 $b=\boldsymbol1$ 就得到双对偶中的元素，却没有 $c_0$ 中的原像。对 $1<p<\infty$，连续取两次共轭指数会回到 $p$，而上述配对把 $Jx$ 对应到原来的同一个数列 $x$，这才证明典范映射满射。只写两个空间“大小相同”并不足以证明自反性。

\section{自反性的弱紧性刻画}
\begin{lemma}[Goldstine 定理]\label{lem:fa-goldstine}
$J(B_X)$ 在 $B_{X^{**}}$ 中弱星稠密。
\end{lemma}
\begin{proof}
给定 $F\in B_{X^{**}}$ 及有限个 $f_1,\ldots,f_m\in X^*$，只需证明 $(F(f_j))_{j=1}^m$ 属于 $\{(f_j(x))_{j=1}^m:x\in B_X\}$ 的闭包。若不属于，在有限维实空间（复情形取实部）中分离点与闭凸集，得到 $f=\sum_jc_jf_j$ 满足
\[
 \operatorname{Re}F(f)>\sup_{x\in B_X}\operatorname{Re}f(x)=\|f\|.
\]
但 $\|F\|\leq1$ 保证 $\operatorname{Re}F(f)\leq\|f\|$，矛盾。有限坐标邻域恰好定义弱星拓扑，故引理成立。
\end{proof}
\begin{theorem}\label{thm:fa-reflexive-compact}
Banach 空间 $X$ 自反当且仅当 $B_X$ 弱紧。
\end{theorem}
\begin{proof}
自反时，$J$ 把 $B_X$ 双射到 $B_{X^{**}}$，并把弱拓扑对应到 $\sigma(X^{**},X^*)$；由 Banach--Alaoglu 定理得紧性。反之，若 $B_X$ 弱紧，则 $J(B_X)$ 在 Hausdorff 弱星拓扑中紧，因而闭；Goldstine 定理又保证它在 $B_{X^{**}}$ 中稠密，所以两者相等，缩放后得 $J(X)=X^{**}$。
\end{proof}
自反空间的闭线性子空间 $M$ 也自反：$B_M=M\cap B_X$ 为弱闭的弱紧集，而 Hahn--Banach 延拓保证 $M$ 的自身弱拓扑与诱导弱拓扑一致。闭子空间的商 $X/M$ 也自反：商映射弱连续，闭单位球的像弱紧；它包含商空间的开单位球，而后者在商闭单位球中弱稠密，故像包含整个商闭单位球，反向包含由商范数定义成立，再用定理\ref{thm:fa-reflexive-compact}。

\section{有界序列的弱收敛子列}
\begin{theorem}\label{thm:fa-reflexive-subsequence}
自反 Banach 空间中的每个有界序列都有弱收敛子列。
\end{theorem}
\begin{proof}
设 $(x_n)$ 有界，令 $M=\overline{\operatorname{span}\{x_n:n\geq1\}}$。$M$ 可分且自反。先说明 $M^*$ 可分：由上一节的度量和 Banach--Alaoglu 定理，$B_{M^*}$ 弱星紧且可度量化，故有可数弱星稠密子集。因 $M$ 自反，$M^*$ 上的弱星拓扑就是弱拓扑；该可数集的凸包弱稠密，定理\ref{thm:fa-convex-weak-closure}使其范数稠密。取有理系数的有限凸组合，便得到范数可数稠密集。

现在把 $Jx_n$ 看作 $(M^*)^*$ 中的有界序列，对可分空间 $M^*$ 使用定理\ref{thm:fa-weak-star-subsequence}，得到弱星收敛子列。极限由自反性等于某个 $Jx$，$x\in M$，故 $x_n$ 的相应子列在 $M$ 中弱收敛；将 $X^*$ 的泛函限制到 $M$，即知它也在 $X$ 中弱收敛。
\end{proof}
更一般的 Eberlein--Šmulian 定理说明：Banach 空间中，一个子集弱紧当且仅当其中每个序列都有弱收敛到该子集某点的子列。本书不证明这一更深结果；弱紧性的进一步理论可参阅\cite{teschl-fa}。上面的自反情形已足够支持后续直接法。

\section{例题与习题}
\begin{example}[最近点存在]
自反 Banach 空间中，任意非空闭凸集 $C$ 到给定点 $x$ 的距离都能达到。
\end{example}
\begin{solve}
取 $y_n\in C$ 使 $\|x-y_n\|\to d=\operatorname{dist}(x,C)$。此序列有界，故有弱收敛子列，记极限为 $y$。闭凸集弱闭，所以 $y\in C$；范数下半连续性给出 $\|x-y\|\leq d$，因而等号成立。唯一性通常需要严格凸性等额外条件。
\end{solve}

所谓严格凸，是指不同单位向量 $u,v$ 的中点满足 $\|(u+v)/2\|<1$。若两个不同点都是距离为 $d>0$ 的最近点，对它们相对于 $x$ 的误差向量归一化，严格凸性会使中点距离小于 $d$，矛盾；$d=0$ 时唯一性直接成立。Hilbert 范数严格凸，这与投影定理的唯一性相符。一般自反空间未必严格凸，例如 $\mathbb R^2$ 配最大值范数，到 $x$ 轴上点集的最近点可能形成一段线段。

\begin{example}[缺少自反性时的闭凸极小化问题]
在 $c_0$ 上令 $f(x)=\sum_{j\geq1}2^{-j}x_j$，并取闭仿射超平面 $C=\{x:f(x)=1\}$。则 $\inf_{x\in C}\|x\|_\infty=1$，但下确界不能达到。
\end{example}
\begin{solve}
因为 $\|f\|=\sum_j2^{-j}=1$，每个 $x\in C$ 都满足 $\|x\|_\infty\geq1$。令前 $N$ 个坐标均为 $(1-2^{-N})^{-1}$，其余为零，就得到 $x^{(N)}\in C$ 且范数趋于一。若某个 $x\in C$ 的范数为一，由 $1=\sum_j2^{-j}\operatorname{Re}x_j$ 及每个 $\operatorname{Re}x_j\leq1$，所有实部必须为一，结合 $|x_j|\leq1$ 得所有 $x_j=1$，不属于 $c_0$。因此闭凸约束和有界极小化序列本身还不够，还需要能保住极限的紧性条件。
\end{solve}
\begin{enumerate}
 \item 证明 $\ell^1$ 不自反。提示：若 $e_n$ 有弱收敛子列，坐标泛函迫使其极限为零，而求和泛函给出矛盾。
 \item\label{exer:fa-norm-attainment} 证明自反空间上的每个连续线性泛函都在闭单位球上达到其模的最大值。
 \item 证明严格凸 Banach 空间中，闭凸集的最近点一旦存在便唯一。
 \item 解释为何 $\ell^\infty$ 不自反：它含有闭子空间 $c_0$。
\end{enumerate}

\chapter{有界算子的谱理论}
\label{chap:fa-spectrum}

本章中 $X$ 为非零复 Banach 空间，$T\in\mathcal B(X)$。有限维矩阵不可逆等价于存在非零核；无限维中还可能因值域不满、逆映射无界等原因失败。谱把这些现象统一起来。

\section{可逆算子与 Neumann 级数}
\begin{proposition}\label{prop:fa-neumann}
若 $S\in\mathcal B(X)$ 且 $\|S\|<1$，则 $I-S$ 可逆，且
\begin{equation}
 (I-S)^{-1}=\sum_{n=0}^{\infty}S^n,
 \qquad\|(I-S)^{-1}\|\leq\frac1{1-\|S\|}.
 \label{eq:fa-neumann-series}
\end{equation}
\end{proposition}
\begin{proof}
$\mathcal B(X)$ 完备，且 $\sum_n\|S^n\|\leq\sum_n\|S\|^n<\infty$，故级数按算子范数收敛。有限几何和左右乘 $I-S$ 均等于 $I-S^{N+1}$，令 $N\to\infty$ 即得双侧逆。
\end{proof}
若 $A$ 可逆且 $\|A^{-1}E\|<1$，则 $A+E=A(I+A^{-1}E)$ 可逆。因此可逆算子构成开集，小扰动不会立即破坏可逆性。

这里的条件是充分条件，不是必要条件。例如 $S=2I$ 时 $\|S\|=2$，$I-S=-I$ 仍可逆，只是几何级数 $\sum S^n$ 不收敛。后面的定理\ref{thm:fa-spectral-radius}还会给出比 $\|S\|<1$ 更一般的级数收敛条件。证明方法的适用范围与所研究结论本身的适用范围应当区分。

扰动估计还可以定量化。令 $\eta=\|A^{-1}E\|<1$，则
\[
 \|(A+E)^{-1}\|\leq\frac{\|A^{-1}\|}{1-\eta},\qquad
 \|(A+E)^{-1}-A^{-1}\|
 \leq\frac{\eta\|A^{-1}\|}{1-\eta}.
\]
第一式由 Neumann 估计得到；第二式由 $(I+A^{-1}E)^{-1}-I=-(I+A^{-1}E)^{-1}A^{-1}E$ 得到。这说明稳定性不仅取决于扰动 $E$ 小不小，还取决于原来逆算子的范数有多大。

\section{预解集、谱与谱的分类}
\begin{definition}
\keyterm{预解集（resolvent set）}$\rho(T)$ 由所有使 $\lambda I-T$ 具有有界双侧逆的 $\lambda\in\mathbb C$ 组成。\keyterm{谱（spectrum）}为 $\sigma(T)=\mathbb C\setminus\rho(T)$。对 $\lambda\in\rho(T)$ 记 $R(\lambda,T)=(\lambda I-T)^{-1}$。
\end{definition}
由有界逆定理，$\lambda I-T$ 双射已足以保证逆有界。采用互不相交的分类：非单射时属于点谱（point spectrum）；单射、值域稠密但不满时属于连续谱（continuous spectrum）；单射且值域不稠密时属于剩余谱（residual spectrum）。点谱中的元素就是特征值。

因此检查一个数 $\lambda$ 时，可以直接研究方程 $(\lambda I-T)x=y$：若齐次方程有非零解，就是点谱；若齐次解唯一但某些数据无法求解，再检查可解数据能否逼近任意数据，分别得到连续谱或剩余谱。有限维单射等价于满射，后两种情况不会出现；这正是无限维谱比矩阵特征值集合更大的原因。

\begin{example}[对角算子的谱逐项可算]\label{ex:fa-diagonal-spectrum}
设 $(a_j)$ 是有界复数列，$Dx=(a_jx_j)$ 作用于 $\ell^2$。则
\[
 \sigma(D)=\overline{\{a_j:j\geq1\}},\qquad
 \sigma_{\mathrm p}(D)=\{a_j:j\geq1\},
\]
其中 $\sigma_{\mathrm p}$ 表示点谱。
\end{example}
\begin{solve}
若 $\lambda$ 不在对角元集合的闭包中，则 $\delta=\inf_j|\lambda-a_j|>0$，逐坐标除法 $y_j\mapsto y_j/(\lambda-a_j)$ 给出范数至多 $1/\delta$ 的逆。若 $\lambda=a_j$，则 $e_j$ 为特征向量。剩下的情形是 $\lambda$ 为极限点但不等于任何 $a_j$：取 $j_k$ 使 $a_{j_k}\to\lambda$，则 $\|(\lambda I-D)e_{j_k}\|\to0$，排除有界逆；每个坐标系数非零又保证单射，且值域包含所有有限支撑数列，因而稠密。由有界逆定理，它不能满射，所以这种点属于连续谱。特别地，$a_j=1/j$ 时，非零谱点全为特征值，零是连续谱点。
\end{solve}

\begin{example}[右移算子的剩余谱]\label{ex:fa-shift-spectrum}
在复 $\ell^2$ 上令 $R(x_1,x_2,\ldots)=(0,x_1,x_2,\ldots)$。则 $\sigma(R)$ 为闭单位圆盘，点谱为空，开单位圆盘为剩余谱，单位圆周为连续谱。
\end{example}
\begin{solve}
先检查核。若 $Rx=\lambda x$，当 $\lambda\ne0$ 时，第一坐标给出 $x_1=0$，依次推出所有坐标为零；当 $\lambda=0$ 时，$R$ 本身单射，仍只有零解。因此点谱为空。$|\lambda|>1$ 时，由 $\|R\|=1$ 与 Neumann 级数得到可逆性。

若 $|\lambda|<1$，向量 $y=(1,\overline\lambda,\overline\lambda^2,\ldots)$ 属于 $\ell^2$，而左移 $L=R^*$ 满足 $Ly=\overline\lambda y$。所以 $(\lambda I-R)^*y=0$；值域与非零向量 $y$ 正交，不稠密，属于剩余谱。特别地，$R$ 的值域是第一坐标为零的闭超平面，这已经解释了零点的情形。

若 $|\lambda|=1$，取单位向量
\[
 x^{(N)}=\frac1{\sqrt N}(1,\lambda^{-1},\ldots,\lambda^{-(N-1)},0,\ldots).
\]
$ (\lambda I-R)x^{(N)}$ 只有第一坐标与第 $N+1$ 个坐标可能非零，各自模为 $1/\sqrt N$，所以其范数为 $\sqrt{2/N}\to0$，不可能有有界逆。同时 $Ly=\overline\lambda y$ 迫使 $y_j=\overline\lambda^{j-1}y_1$，平方可和性又迫使 $y_1=0$，故伴随核为零，值域稠密。结合单射与不可逆性，单位圆周上的点属于连续谱。
\end{solve}
这个算子没有任何特征值，谱却有整整一个圆盘。因而求出全部特征向量，并不等于已经了解一个无限维算子的全部谱。另一方面，伴随算子能把值域的稠密性问题转化为核的问题，这是判断剩余谱的直接工具。

\section{谱的非空性与预解恒等式}
直接将逆算子相减得到
\[
 R(\lambda,T)-R(\mu,T)=(\mu-\lambda)R(\lambda,T)R(\mu,T).
\]
若 $|h|\|R(\lambda,T)\|<1$，Neumann 展开还给出
\[
 R(\lambda+h,T)=\sum_{n=0}^{\infty}(-h)^nR(\lambda,T)^{n+1},
\]
所以预解集开，预解算子在其中解析。

\begin{theorem}\label{thm:fa-spectrum-compact}
$\sigma(T)$ 是非空紧集，且包含于 $\{\lambda:|\lambda|\leq\|T\|\}$。
\end{theorem}
\begin{proof}
若 $|\lambda|>\|T\|$，则
\[
 R(\lambda,T)=\lambda^{-1}\sum_{n=0}^{\infty}(T/\lambda)^n,
 \qquad\|R(\lambda,T)\|\leq\frac1{|\lambda|-\|T\|}.
\]
因此谱有界；预解集开保证谱闭。若谱为空，对每个 $x\in X,f\in X^*$，函数 $\lambda\mapsto f(R(\lambda,T)x)$ 为整函数，在大圆外有界且趋零，在圆内由连续性有界。Liouville 定理使其恒为零。由对偶分离点知 $R(\lambda,T)=0$，这与它是逆算子且 $X\ne\{0\}$ 矛盾。
\end{proof}
复标量域在非空性证明中至关重要；实矩阵的实特征值集合可能为空。实算子的谱通常通过附录\ref{app:fa-complex}中的复化定义。

\section{谱半径与多项式谱映射}
\begin{theorem}\label{thm:fa-spectral-radius}
定义\keyterm{谱半径（spectral radius）}$r(T)=\max_{\lambda\in\sigma(T)}|\lambda|$，则
\begin{equation}
 r(T)=\lim_{n\to\infty}\|T^n\|^{1/n}.
 \label{eq:fa-spectral-radius}
\end{equation}
对任意复系数多项式 $p$，有 $\sigma(p(T))=p(\sigma(T))$。
\end{theorem}
\begin{proof}
次乘性 $\|T^{m+n}\|\leq\|T^m\|\|T^n\|$ 保证根式极限 $L$ 存在并等于其下确界：固定 $m$，写 $n=km+j$，用余数 $0\leq j<m$ 的有限个范数控制，再令 $n\to\infty$。若某次幂为零，则 $L=0$ 单独成立。当 $|\lambda|>L$ 时，根值判别保证 $\sum_nT^n/\lambda^{n+1}$ 收敛并给出逆，故 $r(T)\leq L$。

反向取 $R>r(T)$。由大圆上的 Neumann 级数逐项积分，再在预解区域内变形路径，得到
\[
 T^n=\frac1{2\pi\mathrm i}\int_{|\lambda|=R}\lambda^nR(\lambda,T)\,\mathrm d\lambda.
\]
算子值积分的使用可逐个向量、泛函标量化，见附录\ref{app:fa-complex}。故 $\|T^n\|\leq R^{n+1}\max_{|\lambda|=R}\|R(\lambda,T)\|$，取 $n$ 次方根得 $L\leq R$，再令 $R\downarrow r(T)$。

对非常数 $p$，将 $p(z)-\mu$ 分解为一次因子的乘积。相应算子因子两两交换；有限个交换因子的乘积可逆当且仅当每个因子可逆。因此 $\mu\notin\sigma(p(T))$ 当且仅当 $p(z)=\mu$ 的所有根均不在 $\sigma(T)$。常数多项式情形直接验证。
\end{proof}
一般有 $r(T)\leq\|T\|$，但不必相等；非零幂零矩阵的谱半径为零。

例如欧氏平面上的矩阵 $T=\begin{pmatrix}0&M\\0&0\end{pmatrix}$ 满足 $T^2=0$、$\sigma(T)=\{0\}$，而 $\|T\|=|M|$ 可以很大。范数描述单次作用最大的放大程度，谱半径通过 $\|T^n\|^{1/n}$ 描述反复作用时的渐近尺度。本例第一次可能放大误差，第二次却使所有向量归零，因此两个量没有矛盾。

\section{自伴与正算子的谱}
\begin{proposition}\label{prop:fa-selfadjoint-spectrum}
有界自伴算子 $T$ 的谱为实集，且 $r(T)=\|T\|$。正算子的谱包含于 $[0,\|T\|]$。
\end{proposition}
\begin{proof}
对非实 $\lambda$，内积的虚部给出 $\|(\lambda I-T)x\|\geq|\operatorname{Im}\lambda|\|x\|$。这保证单射及闭值域：像序列收敛时，原像由下界为 Cauchy 序列。同时 $(\lambda I-T)^*=\overline\lambda I-T$ 的核为零，故值域稠密，因而满射。于是非实数均在预解集中。

由 $\|T^*T\|=\|T\|^2$ 及自伴性，归纳得 $\|T^{2^k}\|=\|T\|^{2^k}$；谱半径公式即给出 $r(T)=\|T\|$。若 $T\geq0$、$a>0$，则 $\langle(T+aI)x,x\rangle\geq a\|x\|^2$，同样用下界与伴随核证明 $T+aI$ 可逆，排除负谱点。
\end{proof}

\section{例题与习题}
\begin{example}[谱点不必是特征值]\label{ex:fa-multiplication-spectrum}
在 $L^2(0,1)$ 上，$(Tf)(t)=tf(t)$ 的谱为 $[0,1]$，但点谱为空。
\end{example}
\begin{solve}
若 $\lambda\notin[0,1]$，乘以 $(\lambda-t)^{-1}$ 为有界逆。若 $\lambda\in[0,1]$，取支撑在 $|t-\lambda|<1/n$ 内的归一化示性函数 $f_n$，则 $\|f_n\|_2=1$ 而 $\|(\lambda I-T)f_n\|_2\leq1/n$，排除有界逆。特征方程 $(t-\lambda)f(t)=0$ 迫使 $f$ 支撑在零测单点，故 $f=0$。因算子自伴且核为零，其值域稠密，所以这些谱点均为连续谱点。
\end{solve}
\begin{enumerate}
 \item 对 $a_j=(-1)^j+1/j$ 的 $\ell^2$ 对角算子，求谱、点谱及连续谱，并讨论谱点的孤立性。
 \item 求 $\ell^2$ 上左移 $L$ 的点谱，并证明其谱为闭单位圆盘、剩余谱为空，比较它与例\ref{ex:fa-shift-spectrum}中右移的谱分类。
 \item 证明当 $|\lambda|>\|T\|$ 时，预解算子的上述上界成立，讨论 $T=cI$ 时何时达到等号。
 \item 构造两个具有相同谱但范数不同的有限维算子。
\end{enumerate}

\chapter{紧算子与 Fredholm 理论}
\label{chap:fa-compact}

紧算子把有界输入转化为具有收敛子列的输出。这一性质弱于有限秩，却足以使非零谱点及方程不可解的障碍具有有限维结构。

\section{紧算子的定义与基本性质}
\begin{definition}
线性算子 $K:X\to Y$ 称为\keyterm{紧算子（compact operator）}，若每个有界集的像相对紧，即其闭包紧。对赋范空间，这等价于每个有界序列 $(x_n)$ 都有子列，使 $(Kx_{n_j})$ 在 $Y$ 中收敛。
\end{definition}
紧算子有界，因为单位球像的紧闭包有界。有限秩有界算子紧，因为其有界像位于有限维空间。紧算子与有界算子在任意一侧复合仍紧：一侧保持输入有界，另一侧把输出收敛序列连续地映为收敛序列。有限个紧算子之和也紧，可依次选子列证明。

\begin{proposition}\label{prop:fa-compact-closed}
若 $Y$ 完备，$K_n:X\to Y$ 紧且 $\|K_n-K\|\to0$，则 $K$ 紧。
\end{proposition}
\begin{proof}
给定 $\varepsilon>0$，取 $n$ 使 $\|K_n-K\|<\varepsilon/3$。$K_n(B_X)$ 有有限个半径 $\varepsilon/3$ 的球覆盖，因此 $K(B_X)$ 有有限个半径 $\varepsilon$ 的球覆盖。其闭包全有界，且作为完备空间的闭集完备，由定理\ref{thm:fa-metric-compactness}可知紧。
\end{proof}
在无限维空间中，恒等算子不紧，这是推论\ref{cor:fa-unit-ball}的另一种表述。

\begin{example}[尾部衰减怎样产生紧性]\label{ex:fa-diagonal-compactness}
$\ell^2$ 上的有界对角算子 $Dx=(a_jx_j)$ 紧，当且仅当 $a_j\to0$。
\end{example}
\begin{solve}
若 $a_j\to0$，只保留前 $N$ 个输出坐标，得到有限秩算子 $D_N$，且 $\|D-D_N\|=\sup_{j>N}|a_j|\to0$，所以 $D$ 紧。反之，若对角元不趋零，可取子列使 $|a_{j_k}|\geq\varepsilon>0$。输入 $e_{j_k}$ 全在单位球中，输出却两两相距至少 $\sqrt2\varepsilon$，不可能有收敛子列，矛盾。紧性在这里表示：无论单位输入怎样分布，高编号坐标对输出的影响最终都统一变小。
\end{solve}

这个例子也说明紧性与把距离缩小不是一回事。无限维空间上的 $I/2$ 把任意距离缩小一半，但不紧；有限秩算子 $100P$ 可以显著放大长度，却仍然紧。压缩性控制任意两个输入之间的距离，紧性控制整族有界输出的子列行为。

\section{有限秩逼近与积分算子}
若 $K:H_1\to H_2$ 是 Hilbert 空间间的紧算子，对 $K(B_{H_1})$ 取有限 $\varepsilon$ 网 $y_1,\ldots,y_m$。设 $P$ 为到其线性张成空间的正交投影，则最佳逼近性质给出 $\|(I-P)K\|\leq\varepsilon$，而 $PK$ 为有限秩。因此 Hilbert 空间间的紧算子恰为有限秩算子的范数极限；这一逆向逼近结论不能无条件推广到所有 Banach 空间。

\begin{theorem}[Arzelà--Ascoli]\label{thm:fa-ascoli}
设 $Q$ 为紧度量空间。$\mathcal F\subset C(Q)$ 在一致范数下相对紧，当且仅当它一致有界且等度连续。等度连续指对每个 $\varepsilon>0$ 存在 $\delta>0$，使 $d(s,t)<\delta$ 时 $|f(s)-f(t)|<\varepsilon$ 对所有 $f\in\mathcal F$ 成立。
\end{theorem}
\begin{proof}
相对紧集可用有限个一致范数小球覆盖。有限个中心函数均有界且一致连续，利用三角不等式把这两种性质传给整个函数族。

反之，由等度连续性选 $\delta$，再以有限个 $\delta$ 球覆盖 $Q$，记球心为 $t_1,\ldots,t_m$。一致有界性保证向量 $(f(t_j))_{j=1}^m$ 处于有限维有界集，故可取有限细网，并从每个非空网格选择一个函数。两个函数若各球心处数值足够接近，利用等度连续性即可在整个 $Q$ 上一致接近。因此 $\mathcal F$ 全有界；$C(Q)$ 完备，故闭包紧。
\end{proof}
若 $k\in C([a,b]^2)$，则
\[
 (Kf)(t)=\int_a^b k(t,s)f(s)\,\mathrm ds
\]
在 $C([a,b])$ 上紧：单位球的像由 $(b-a)\|k\|_\infty$ 控制，而 $|Kf(t)-Kf(t')|\leq\int_a^b|k(t,s)-k(t',s)|\,\mathrm ds$ 给出等度连续性。在 $L^2$ 框架中，若核平方可积，则 Cauchy--Schwarz 给出 $\|K\|\leq\|k\|_{L^2}$；用有限个乘积函数逼近核，可得有限秩逼近，故算子也紧。

\section{紧扰动的核与闭值域}
在 Arzelà--Ascoli 定理中，等度连续要求先选择同一个 $\delta$，再对函数族中所有函数使用它。各个函数分别连续还不够：$f_n(t)=t^n$ 在 $[0,1]$ 上一致有界，但在 $1$ 附近越来越陡，不能选择统一的 $\delta$。任何子列的逐点极限在 $[0,1)$ 上为零、在 $1$ 处为一，因此不可能一致收敛到连续函数。另一方面，常值函数族 $f_n(t)=n$ 等度连续，却不一致有界，也没有一致收敛子列。

\begin{proposition}[紧算子把弱收敛变为范数收敛]\label{prop:fa-compact-weak-strong}
若 $K:X\to Y$ 是赋范空间间的紧线性算子，$x_n\rightharpoonup x$，则 $\|Kx_n-Kx\|\to0$。
\end{proposition}
\begin{proof}
弱收敛序列有界。若结论不成立，可取 $\varepsilon>0$ 与子列，使 $\|Kx_{n_j}-Kx\|\geq\varepsilon$。紧性使其有进一步子列在范数下趋于某个 $y\in Y$。有界线性算子保持弱收敛，故该子列也弱趋于 $Kx$；范数收敛又蕴含弱趋于 $y$。连续对偶分离点，迫使 $y=Kx$，与距离下界矛盾。
\end{proof}
证明中的“任取一个反常子列，再提取更好的子列”是常用方法。只证明存在一个子列收敛到 $Kx$，不足以证明整个输出序列收敛。

\section{紧扰动的核与闭值域}
本节设 $X$ 为 Banach 空间，$K\in\mathcal B(X)$ 紧，$A=I-K$。
\begin{lemma}\label{lem:fa-compact-range}
$\ker A$ 有限维，$\operatorname{Ran}A$ 闭，并存在 $C>0$ 使
\[
 \operatorname{dist}(x,\ker A)\leq C\|Ax\|\quad(x\in X).
\]
\end{lemma}
\begin{proof}
在 $\ker A$ 上 $K=I$，故其单位球紧，核有限维。若距离估计失败，可缩放并减去核内向量，得到 $\operatorname{dist}(x_n,\ker A)=1$、$\|x_n\|\leq2$ 且 $Ax_n\to0$。紧性给出 $Kx_n$ 的收敛子列，于是 $x_n=Ax_n+Kx_n$ 的相应子列收敛到 $x\in\ker A$，与距离为一矛盾。

若 $Ax_n\to y$，由距离估计可在不改变 $Ax_n$ 的条件下选取有界代表元。再由紧性选 $Kx_n$ 收敛子列，便有 $x_n\to x$ 且 $Ax=y$，所以值域闭。
\end{proof}

\section{Fredholm 二择一与指数}
\begin{theorem}[Fredholm 二择一]\label{thm:fa-fredholm-alternative}
设 $X$ 为 Banach 空间，$K\in\mathcal B(X)$ 为紧算子，$A=I-K$。则 $A$ 单射当且仅当满射，此时逆有界。一般地，$A$ 的核有限维、值域闭，且
\[
 \dim\ker A=\dim(X/\operatorname{Ran}A).
\]
方程 $Ax=y$ 可解当且仅当 $f(y)=0$ 对所有 $f\in\ker A'$ 成立。
\end{theorem}
\begin{proof}
考虑 $N_j=\ker A^j$ 与 $R_j=\operatorname{Ran}A^j$。因 $A^j=I-K_j$，其中 $K_j$ 是含因子 $K$ 的算子多项式，故 $K_j$ 紧；由引理\ref{lem:fa-compact-range}，$N_j$ 有限维且 $R_j$ 闭。

两个链都必在有限步稳定。若 $N_j$ 严格递增，用 Riesz 引理选单位向量 $x_j\in N_j$，到 $N_{j-1}$ 的距离大于 $1/2$。当 $j>i$ 时，$Ax_j,Kx_i\in N_{j-1}$，故 $\|Kx_j-Kx_i\|>1/2$，违反紧性。若 $R_j$ 严格递减，选单位向量 $x_j\in R_{j-1}$，到 $R_j$ 的距离大于 $1/2$；当 $j>i$ 时，$Kx_j,Ax_i\in R_i$，再次得到像之间距离大于 $1/2$。链一旦相邻相等，以取原像或取像可知后面也都相等。

取足够大 $m$ 使两链稳定。则
\[
 X=N_m\oplus R_m.
\]
确实，若 $x=A^my\in N_m$，则 $y\in N_{2m}=N_m$，所以 $x=0$。又因 $R_m=R_{2m}$，任意 $x$ 都可选 $y$ 使 $A^mx=A^{2m}y$，从而 $x-A^my\in N_m$。在 $R_m$ 上，$A$ 单射且满射；在有限维 $N_m$ 上，核维数等于值域余维。分块考察便得全空间的维数等式，所以单射与满射等价，逆的有界性由有界逆定理得到。

最后，$\ker A'=(\operatorname{Ran}A)^\perp$。由值域闭及命题\ref{prop:fa-distance-dual}，一个 $y$ 属于值域当且仅当所有这些泛函都在 $y$ 上消失。
\end{proof}
有界线性算子 $A:X\to Y$ 若核有限维、值域闭且值域余维有限，称为\keyterm{Fredholm 算子（Fredholm operator）}。其指数定义为
\begin{equation}
 \operatorname{ind}A=\dim\ker A-\dim(Y/\operatorname{Ran}A).
 \label{eq:fa-fredholm-index}
\end{equation}
上面的定理说明 $I-K$ 的指数为零。在 Hilbert 空间中，相容条件写成 $y\perp\ker(I-K^*)$。

\begin{example}[把 Fredholm 相容条件算出来]\label{ex:fa-rank-one-fredholm}
在实 $C([0,1])$ 中求解
\[
 u(t)-\alpha t\int_0^1u(s)\,\mathrm ds=g(t),
 \qquad g\in C([0,1]),\quad\alpha\in\mathbb R.
\]
\end{example}
\begin{solve}
令 $c=\int_0^1u$、$b=\int_0^1g$。原方程给出 $u(t)=g(t)+\alpha tc$，两侧积分得到 $(1-\alpha/2)c=b$。当 $\alpha\ne2$ 时，
\[
 u(t)=g(t)+\frac{\alpha b}{1-\alpha/2}t
\]
是唯一解。当 $\alpha=2$ 时，必须有 $b=0$；若此条件成立，每个 $u(t)=g(t)+2ct$ 都是解，$c$ 任意。算子 $Ku(t)=\alpha t\int u$ 的值域至多一维，所以紧。在临界参数 $\alpha=2$ 时，齐次解空间为 $\operatorname{span}\{t\}$，值域为 $\{g:\int g=0\}$，核维数与值域余维同为一；泛函 $g\mapsto\int g$ 正是检测可解性的观测。分母还表明，接近临界参数时，唯一解虽然存在，却可能对数据十分敏感。
\end{solve}

\section{紧算子的非零谱}
\begin{theorem}\label{thm:fa-compact-spectrum}
复 Banach 空间上紧算子 $K$ 的每个非零谱点都是特征值，其特征空间有限维；对每个 $\varepsilon>0$，模至少为 $\varepsilon$ 的谱点只有有限个。因此非零谱至多可数，唯一可能的聚点是零。
\end{theorem}
\begin{proof}
对 $\lambda\ne0$，将 $\lambda I-K$ 写为 $\lambda(I-K/\lambda)$，Fredholm 二择一及有限维核立即给出前两项。

若有互异特征值 $\lambda_n$ 满足 $|\lambda_n|\geq\varepsilon$，取相应特征向量，令其前 $n$ 个的张成空间为 $E_n$。不同特征值的特征向量线性无关，且 $(K-\lambda_n I)E_n\subset E_{n-1}$。由 Riesz 引理取 $x_n\in E_n$，$\|x_n\|=1$，到 $E_{n-1}$ 距离大于 $1/2$。若 $n>m$，则 $Kx_n-Kx_m$ 与 $\lambda_nx_n$ 相差 $E_{n-1}$ 中的向量，故其范数大于 $\varepsilon/2$，与紧性矛盾。
\end{proof}
无限维空间上零必在紧算子的谱中，否则 $I=K^{-1}K$ 紧，违反闭单位球不紧。但零未必是特征值，如 $\ell^2$ 上对角元为 $1/n$ 的算子。

\section{例题与习题}
\begin{example}[Volterra 算子]
$C([0,1])$ 上 $(Vf)(t)=\int_0^t f(s)\,\mathrm ds$ 紧，且 $\sigma(V)=\{0\}$。
\end{example}
\begin{solve}
单位球像一致有界且满足 $|Vf(t)-Vf(s)|\leq|t-s|$，由 Arzelà--Ascoli 得紧性。逐次积分或归纳得 $\|V^n\|\leq1/n!$，故谱半径为零。谱非空，所以只有零；$V$ 单射，表明零可以是紧算子的非特征谱点。
\end{solve}

这里可以把幂算子写得更具体：对 $n\geq1$，反复积分得到
\[
 V^nf(t)=\frac1{(n-1)!}\int_0^t(t-s)^{n-1}f(s)\,\mathrm ds,
 \qquad \|V^n\|=\frac1{n!}.
\]
上界由积分核的绝对值积分得到，下界用 $f=1$、$t=1$ 达到。于是对任意标量 $\alpha$，$\sum_{n\geq0}\alpha^nV^n$ 都在算子范数下绝对收敛，并给出 $(I-\alpha V)^{-1}$。所以某个简单的压缩估计要求 $|\alpha|<1$，并不意味着参数超出这个范围后方程就没有唯一解；更精确的幂估计能扩大求解范围。
\begin{enumerate}
 \item 在 $\ell^2$ 上令 $Kx=(x_j/j^a)$，其中 $a>0$。求前 $N$ 项有限秩截断的精确算子范数误差，并给出误差不超过 $\varepsilon>0$ 所需的 $N$。
 \item 将例\ref{ex:fa-rank-one-fredholm}推广为 $Ku=\phi\,\ell(u)$，其中 $\phi\in X$、$\ell\in X^*$。分别在 $\ell(\phi)\ne1$ 与 $\ell(\phi)=1$ 时讨论 $I-K$。
 \item 设 $K$ 为 Hilbert 空间上的紧算子，$(e_n)$ 为标准正交序列。证明 $\|Ke_n\|\to0$，并用它再次说明无限维恒等算子不紧。
 \item 说明紧算子的有限秩近似为何不意味着它具有有限维值域。
\end{enumerate}

\chapter{紧自伴算子的谱分解}
\label{chap:fa-compact-spectral}

本章把有限维对称矩阵的正交对角化推广到紧自伴算子。紧性保证能够从近似特征向量中提取极限，自伴性保证不同特征方向正交，两者结合得到可用于计算的展开。以下几何证明同时适用于实、复 Hilbert 空间；因此尚未学习复分析的读者，也可以先读本章的离散谱分解，再回头学习一般有界算子的谱理论。

\section{特征值与正交性}
若 $Tx=\lambda x$、$x\ne0$，且 $T=T^*$，则
\[
 \lambda\|x\|^2=\langle Tx,x\rangle
 =\langle x,Tx\rangle=\overline\lambda\|x\|^2,
\]
所以 $\lambda$ 为实数。若 $Ty=\mu y$，$\lambda\ne\mu$，则 $(\lambda-\mu)\langle x,y\rangle=0$，故 $x\perp y$。某个特征空间的正交补在 $T$ 下保持不变，因为对该特征空间中的 $e$，$\langle Tx,e\rangle=\langle x,Te\rangle=0$。

\section{极值与非零特征值}
\begin{proposition}\label{prop:fa-compact-extreme}
非零紧自伴算子 $T$ 有特征值 $\lambda$ 满足 $|\lambda|=\|T\|$。因此
\[
 \max_{\|x\|=1}|\langle Tx,x\rangle|=\|T\|.
\]
\end{proposition}
\begin{proof}
先证明对任意有界自伴算子，都有
\[
 m:=\sup_{\|x\|=1}|\langle Tx,x\rangle|=\|T\|.
\]
由 Cauchy--Schwarz 知 $m\leq\|T\|$。反向取单位向量 $x,y$，必要时给 $y$ 乘一个模为一的标量，使 $\langle Tx,y\rangle$ 为非负实数。记 $q(z)=\langle Tz,z\rangle$，利用自伴性和平行四边形恒等式，得到
\[
 |\langle Tx,y\rangle|
 =\frac14\bigl(q(x+y)-q(x-y)\bigr)
 \leq\frac m4\bigl(\|x+y\|^2+\|x-y\|^2\bigr)=m.
\]
对 $y$ 取上确界，再对 $x$ 取上确界，便得 $\|T\|\leq m$。实情形只需必要时把 $y$ 换为 $-y$。

因 $T\ne0$，$m>0$。选单位向量 $x_n$ 使 $|q(x_n)|\to m$，再取子列，使 $q(x_n)\to\lambda$，其中 $\lambda=m$ 或 $-m$。由于 $\|Tx_n\|\leq m$，
\[
 \|Tx_n-\lambda x_n\|^2
 =\|Tx_n\|^2-2\lambda q(x_n)+\lambda^2
 \leq 2m^2-2\lambda q(x_n)\longrightarrow0.
\]
紧性给出 $Tx_n$ 的收敛子列；上式与 $\lambda\ne0$ 又使相应的 $x_n$ 收敛到某个单位向量 $x$。取极限得 $Tx=\lambda x$，于是极值达到，且 $|\lambda|=\|T\|$。
\end{proof}
这段证明只用了 Hilbert 空间的几何与紧性，没有依赖复分析，因此也直接适用于实 Hilbert 空间。关键是先把接近二次型极值的向量变成近似特征向量，再借助紧算子的输出收敛推出输入收敛；并没有把整个单位球当作范数紧集。

\section{紧自伴谱定理}
\begin{theorem}\label{thm:fa-compact-spectral}
设 $T$ 为可分 Hilbert 空间 $H$ 上的紧自伴算子。其非零特征值可按重数列为有限列或可数列 $(\lambda_n)$，相应标准正交特征向量为 $(e_n)$，并且
\[
 H=\ker T\oplus\overline{\operatorname{span}\{e_n\}}.
\]
若列无穷，则 $\lambda_n\to0$；对任意 $x\in H$，
\begin{equation}
 Tx=\sum_n\lambda_n\langle x,e_n\rangle e_n,
 \label{eq:fa-compact-spectral-expansion}
\end{equation}
级数按 Hilbert 范数收敛。
\end{theorem}
\begin{proof}
若 $T\ne0$，由命题\ref{prop:fa-compact-extreme}选择绝对值最大的特征值及一个单位特征向量，然后把 $T$ 限制到其正交补，继续这一过程。限制仍紧且自伴。若有限步后限制为零，分解已经得到；否则得到正交列，$\|Te_n-Te_m\|^2=\lambda_n^2+\lambda_m^2$。紧性迫使 $\lambda_n\to0$。

记 $M=\overline{\operatorname{span}\{e_n\}}$。任意 $x\in M^\perp$ 都属于每一步的剩余空间，其上算子范数等于下一步选出的 $|\lambda_n|$；令 $n\to\infty$ 得 $Tx=0$。反之，$x\in\ker T$ 必与所有非零特征向量正交，所以 $M^\perp=\ker T$。在 $M$ 中使用正交展开，在核中算子为零，即得式\eqref{eq:fa-compact-spectral-expansion}。
\end{proof}
$x$ 本身的展开应为 $x=P_{\ker T}x+\sum_n\langle x,e_n\rangle e_n$，不可遗漏核中的分量。按 $|\lambda_n|$ 递减排列时，有限秩截断 $T_Nx=\sum_{n\leq N}\lambda_n\langle x,e_n\rangle e_n$ 满足
\[
 \|T-T_N\|=\sup_{n>N}|\lambda_n|.
\]
上界由正交性得到，下界在相应特征向量上取到或逼近；有限列情形中，空尾部的上确界在此约定为零，即保留全部非零特征项后误差为零。

\section{正算子的极小极大原理}
设 $T\geq0$ 紧，非零特征值按 $\lambda_1\geq\lambda_2\geq\cdots>0$ 排列。对存在的第 $n$ 个特征值，有
\[
 \lambda_n=\min_{\dim L=n-1}\ \sup_{\substack{x\perp L\\\|x\|=1}}
 \langle Tx,x\rangle.
\]
证明只需两个方向。取 $L=\operatorname{span}\{e_1,\ldots,e_{n-1}\}$，谱展开给出上界 $\lambda_n$；任意 $n-1$ 维的 $L$ 都使 $\operatorname{span}\{e_1,\ldots,e_n\}\cap L^\perp$ 含单位向量，该向量的二次型至少为 $\lambda_n$，得到下界。这是极小极大原理（min--max principle）的一个常用形式。

当 $n=1$ 时，它就是在全部单位向量中最大化 $\langle Tx,x\rangle$。当 $n=2$ 时，先删去一个方向，再在其正交补上寻找最大的二次型；最有效的删法是删去第一特征方向，此时剩余最大值为 $\lambda_2$。一般地，删去前 $n-1$ 个最强方向，剩余最大值就是第 $n$ 个特征值。这个解释把公式中的“先上确界、再极小化”与逐个寻找特征方向联系起来。

\section{奇异值与 Green 算子}
若 $T:H_1\to H_2$ 是可分 Hilbert 空间间的紧算子，则 $T^*T$ 为正紧自伴算子。设其正特征值为 $s_n^2$，对应单位特征向量 $v_n$，定义 $u_n=Tv_n/s_n$。计算内积可知 $(u_n)$ 标准正交，并且
\[
 Tx=\sum_n s_n\langle x,v_n\rangle u_n.
\]
数 $s_n$ 称为\keyterm{奇异值（singular value）}。核 $\ker(T^*T)=\ker T$ 保证展开没有遗漏非零输出。按降序截断后的误差为 $s_{N+1}$。若至少有 $N+1$ 个正奇异值，任何秩至多 $N$ 的算子 $F$ 在 $\operatorname{span}\{v_1,\ldots,v_{N+1}\}$ 中都有单位核向量，故 $\|T-F\|\geq s_{N+1}$；若正奇异值只有至多 $N$ 个，截断已经等于 $T$，并约定 $s_{N+1}=0$。因此奇异值截断是算子范数下的最佳秩 $N$ 逼近。

\begin{example}[区间上的 Green 算子]\label{ex:fa-green}
在 $L^2(0,1)$ 上令
\[
 (Gf)(t)=\int_0^1(\min(t,s)-ts)f(s)\,\mathrm ds.
\]
则 $G$ 紧、自伴且单射，特征值为 $1/(n\pi)^2$，标准正交特征函数为 $e_n(t)=\sqrt2\sin(n\pi t)$。
\end{example}
\begin{solve}
核连续且对称，故紧性与自伴性成立。将积分在 $s=t$ 处分开并求导，得到 $u=Gf$ 满足 $-u''=f$（弱导数意义）及 $u(0)=u(1)=0$，从而 $G$ 单射。若 $Gf=\lambda f$，$\lambda\ne0$，便有 $-f''=\lambda^{-1}f$ 及零边界条件。常微分方程求解给出 $\lambda=1/(n\pi)^2$ 和正弦函数；负值及零频率仅有零解。由紧自伴谱定理及核为零，这些特征函数的闭线性张成空间就是 $L^2(0,1)$。
\end{solve}

积分与微分之间的计算为
\[
 u(t)=(1-t)\int_0^t sf(s)\,\mathrm ds
       +t\int_t^1(1-s)f(s)\,\mathrm ds,
\]
\[
 u'(t)=-\int_0^t sf(s)\,\mathrm ds
       +\int_t^1(1-s)f(s)\,\mathrm ds,
 \qquad u''(t)=-f(t)\quad\text{几乎处处}.
\]
因为区间有限，$L^2$ 函数属于 $L^1$，上述变上限积分绝对连续，微分公式在几乎处处意义下成立。若 $f=\sum_nf_ne_n$，则 $u=Gf=\sum_n f_n e_n/(n\pi)^2$。积分求解把高频坐标缩小得更多，这既解释了 $G$ 的紧性，也说明逆向求导为什么会放大高频误差。

\section{谱展开怎样解方程}
设 $T$ 为紧自伴算子，非零特征数据如式\eqref{eq:fa-compact-spectral-expansion}。考虑实参数 $\alpha$ 下的方程 $(I-\alpha T)u=f$。写 $f=f_0+\sum_nf_ne_n$，其中 $f_0\in\ker T$，则逐个正交方向给出
\[
 P_{\ker T}u=f_0,\qquad (1-\alpha\lambda_n)\langle u,e_n\rangle=f_n.
\]
若所有分母均非零，则
\[
 u=f_0+\sum_n\frac{f_n}{1-\alpha\lambda_n}e_n.
\]
这个形式解确实属于 $H$：无穷列情形中 $\lambda_n\to0$，尾部分母的绝对值至少为 $1/2$，剩下有限个分母有正的最小绝对值，因此系数仍平方可和。若某个分母为零，必须要求对应的 $f_n=0$；满足这些条件时，相应的 $u$ 系数任意，其余系数按上式确定。这就是正交坐标中的 Fredholm 二择一。

直接解 $Tu=f$ 则有另一种困难。除了 $f\perp\ker T$，还必须有 $\sum_n|f_n/\lambda_n|^2<\infty$，才能把形式原像放回 $H$。仅有正交条件只能保证 $f$ 属于值域的闭包，不能保证属于值域本身。对于 $\lambda_n\to0$ 的无限秩情形，除法会不断放大某些坐标；这与 $I-\alpha T$ 的分母趋于一形成清楚的区别。

\section{习题}
\begin{enumerate}
 \item 验证实对角紧算子的谱分解，并计算截断误差。
 \item 若紧自伴算子单射，证明其值域稠密；它的值域是否一定闭？
 \item 利用例\ref{ex:fa-green}，求 $-u''=\sin(2\pi t)$、$u(0)=u(1)=0$ 的解。
 \item 用乘法算子例\ref{ex:fa-multiplication-spectrum}说明一般有界自伴算子不能都按离散特征向量展开。
\end{enumerate}

\chapter{变分方法与方程的弱解}
\label{chap:fa-variational}

许多方程可以从能量极小化或对所有测试函数成立的积分恒等式得到。以下以实空间为主，展示抽象定理怎样给出具体问题的存在性、唯一性和稳定性。

\section{能量最小化的直接法}
函数 $E:C\to\mathbb R\cup\{+\infty\}$ 称为弱序列下半连续，若 $u_n\rightharpoonup u$ 且 $u_n,u\in C$ 时，$E(u)\leq\liminf_nE(u_n)$；称为\keyterm{强制的（coercive）}，若沿 $C$ 中任意满足 $\|u_n\|\to\infty$ 的序列都有 $E(u_n)\to+\infty$。

允许取 $+\infty$ 可以把不允许的状态并入目标函数；有效定义域是 $\operatorname{dom}E=\{u\in C:E(u)<+\infty\}$。所谓严格凸，是要求有效定义域为凸集，且不同的 $u,v$ 与 $0<t<1$ 满足 $E((1-t)u+tv)<(1-t)E(u)+tE(v)$。弱下半连续表示取弱极限时能量不会向上跳；这恰好适合证明极小点，因为极小化序列的极限只要能量不增，就必须达到下确界。

\begin{theorem}[直接法]\label{thm:fa-direct-method}
设 $X$ 为实自反 Banach 空间，$C\subset X$ 非空且弱序列闭。若 $E:C\to\mathbb R\cup\{+\infty\}$ 弱序列下半连续、强制，且 $-\infty<\inf_CE<+\infty$，则 $E$ 在 $C$ 上取得最小值。若 $C$ 凸且 $E$ 在有效定义域上严格凸，则极小点唯一。
\end{theorem}
\begin{proof}
选 $u_n\in C$ 使 $E(u_n)\to m=\inf_CE$。若序列不有界，便可取范数趋于无穷的子列，强制性将迫使其能量趋于无穷，与趋于有限 $m$ 矛盾。因此存在弱收敛子列 $u_{n_k}\rightharpoonup u$。约束的弱序列闭性给出 $u\in C$，下半连续性给出 $E(u)\leq m$，所以等号成立。若存在两个不同极小点，严格凸性使其中点能量严格小于 $m$，矛盾。
\end{proof}
范数连续的凸函数弱下半连续：其各个下水平集是范数闭凸集，因而弱闭。这个判别常比直接估计弱极限更方便。

\begin{example}[直接法的完整使用]
在实 Hilbert 空间 $H$ 中，设 $C\subset H$ 非空、闭且凸，$g\in H$。则 $E(u)=\frac12\|u\|^2-\langle u,g\rangle$ 在 $C$ 上有唯一极小点。
\end{example}
\begin{solve}
Cauchy--Schwarz 给出 $E(u)\geq\frac12\|u\|^2-\|g\|\|u\|$，所以 $E$ 有下界且强制；在任一固定 $u_0\in C$ 处有限，故下确界也是有限数。平方范数弱下半连续，线性项弱连续，所以 $E$ 弱下半连续。Hilbert 空间自反，闭凸集弱闭，直接法给出存在性。平方范数严格凸，减去线性项仍严格凸，所以极小点唯一。配方 $E(u)=\frac12\|u-g\|^2-\frac12\|g\|^2$ 又说明它恰为 $P_Cg$，把直接法与投影定理联系起来。
\end{solve}

\section{Lax--Milgram 定理}
\begin{theorem}\label{thm:fa-lax-milgram}
设 $H$ 为实 Hilbert 空间，双线性形式 $a:H\times H\to\mathbb R$ 满足
\[
 |a(u,v)|\leq M\|u\|\|v\|,\qquad
 a(v,v)\geq\alpha\|v\|^2,
 \qquad M<\infty,\quad\alpha>0.
\]
则对每个 $f\in H^*$，存在唯一 $u\in H$ 满足
\begin{equation}
 a(u,v)=f(v)\quad(v\in H),\qquad
 \|u\|\leq\alpha^{-1}\|f\|.
 \label{eq:fa-lax-milgram}
\end{equation}
\end{theorem}
\begin{proof}
由 Riesz 表示，存在有界线性算子 $A$ 使 $a(u,v)=\langle Au,v\rangle$，并存在 $g$ 使 $f(v)=\langle g,v\rangle$。强制性给出 $\|Au\|\geq\alpha\|u\|$，所以 $A$ 单射且值域闭。若 $z\perp\operatorname{Ran}A$，则 $a(v,z)=0$ 对所有 $v$ 成立，取 $v=z$ 得 $z=0$。故值域又稠密，从而等于 $H$。唯一解为 $A^{-1}g$，同一下界给出所述估计。
\end{proof}
若 $a$ 对称，解还等价于最小化 $E(v)=\frac12a(v,v)-f(v)$，因为
\[
 E(u+h)-E(u)=\tfrac12a(h,h)\geq\tfrac\alpha2\|h\|^2.
\]
反过来，在任意方向对能量的一元二次式取导，得到弱方程。非对称情形不能用这一能量直接替代原方程，因为能量只反映双线性形式的对称部分。

\section{弱导数与 Sobolev 空间}
Lax--Milgram 定理的下界 $a(v,v)\geq\alpha\|v\|^2$ 也称强椭圆性或强制性估计；它比仅有 $a(v,v)>0$（$v\ne0$）更强。在 $\ell^2$ 上，$a(u,v)=\sum_ju_jv_j/j$ 对非零 $v$ 确有 $a(v,v)>0$，但 $a(e_j,e_j)=1/j\to0$，不存在统一的 $\alpha>0$。对应算子正是 $u_j\mapsto u_j/j$，其值域不满，故并非每个连续右端都可解。这说明无限维中的正定性不能自动提供稳定下界。

若右端从 $f$ 变为 $\widetilde f$，相减并再次使用定理，可得 $\|u-\widetilde u\|\leq\alpha^{-1}\|f-\widetilde f\|$。因此强制常数同时控制存在性证明中的逆算子和数据扰动的放大程度。接下来选择 Sobolev 空间，正是为了使微分方程转成满足这类估计的双线性形式。

\section{弱导数与 Sobolev 空间}
设 $\Omega\subset\mathbb R^d$ 开，$C_c^\infty(\Omega)$ 表示光滑且支撑为 $\Omega$ 内紧集的测试函数。对局部可积 $u$，若局部可积 $v_j$ 满足
\[
 \int_\Omega u\,\partial_j\varphi\,\mathrm dx
 =-\int_\Omega v_j\varphi\,\mathrm dx
 \quad(\varphi\in C_c^\infty(\Omega)),
\]
称 $v_j$ 为 $u$ 的第 $j$ 个\keyterm{弱导数（weak derivative）}，记为 $\partial_ju$。它在几乎处处意义下唯一；所用事实是局部可积函数若对所有测试函数积分为零，则几乎处处为零，可由光滑逼近恒等元证明。

这个定义把导数转移到光滑测试函数上，使 $u$ 本身无需处处可微。测试函数支撑远离边界，所以没有边界项。若 $u$ 经典可微且导数局部可积，分部积分立即表明经典导数也是弱导数；新定义因此延续了原有微积分，而非另行指定一个不相关的量。

\begin{example}[尖角允许，跳跃需要不同的导数概念]
在 $(-1,1)$ 上，$u(t)=|t|$ 的弱导数为 $v(t)=-1$（$t<0$）、$v(t)=1$（$t>0$），在零点任意赋值均可；而阶跃函数 $w(t)=\mathbf1_{(0,1)}(t)$ 没有局部可积的弱导数。
\end{example}
\begin{solve}
把 $\int |t|\varphi'(t)$ 在零点分开，两侧分部积分，零点处的边界项因 $|0|=0$ 消失，得到 $-\int v\varphi$。所以 $|t|\in H^1(-1,1)$，尽管它在零点没有经典导数。阶跃函数则满足 $\int w\varphi'=-\varphi(0)$。若存在局部可积 $v$，应有 $\int v\varphi=\varphi(0)$。选择满足 $0\leq\varphi_n\leq1$、$\varphi_n(0)=1$ 且支撑缩到零点的光滑函数，由局部可积性及支配收敛，左侧趋零，右侧却恒为一，矛盾。在分布理论中这个导数称为 Dirac 点质量，但它不是 $L^1_{\mathrm{loc}}$ 函数，因此阶跃函数不属于 $H^1(-1,1)$。
\end{solve}

\keyterm{Sobolev 空间（Sobolev space）}$H^1(\Omega)$ 定义为 $u\in L^2(\Omega)$ 且所有一阶弱导数属于 $L^2(\Omega)$ 的集合，赋予范数
\[
 \|u\|_{H^1}^2=\|u\|_2^2+\sum_{j=1}^d\|\partial_ju\|_2^2.
\]
它是 Hilbert 空间：若 $u_n$ 及各弱导数分别在 $L^2$ 中收敛到 $u,v_j$，则在测试恒等式中利用 Cauchy--Schwarz 取极限，得 $v_j=\partial_ju$。定义 $H_0^1(\Omega)$ 为 $C_c^\infty(\Omega)$ 在 $H^1$ 中的闭包，所以它也完备；这是齐次边界条件的一种不依赖逐点边界值的定义。

若 $\Omega$ 有界，则存在 $C_\Omega$ 使
\begin{equation}
 \|u\|_2\leq C_\Omega\|\nabla u\|_2\quad(u\in H_0^1(\Omega)).
 \label{eq:fa-poincare}
\end{equation}
这是 Poincaré 不等式。对测试函数，将其零延拓到包含 $\Omega$ 的有界长方体，沿一个坐标方向使用微积分基本定理及 Cauchy--Schwarz，再积分即可得估计；一般 $u$ 由定义中的稠密性取极限。因此 $\|\nabla u\|_2$ 是 $H_0^1$ 上与原范数等价的 Hilbert 范数。

在区间 $(0,1)$ 上，上述计算尤其清楚：对 $u\in C_c^\infty(0,1)$，
\[
 |u(t)|^2=\left|\int_0^tu'(s)\,\mathrm ds\right|^2
 \leq t\int_0^t|u'(s)|^2\,\mathrm ds\leq\|u'\|_2^2.
\]
再对 $t$ 积分，得到 $\|u\|_2\leq\|u'\|_2$，一般 $H_0^1$ 函数由稠密性得到同一结论。若换成 $H^1(0,1)$，常值函数 $u=1$ 的导数为零而函数范数非零，估计立刻失败。因此边界条件承担了排除常数核的作用。

\section{Poisson 方程的弱形式}
对有界开集 $\Omega$，用梯度范数赋予 $H_0^1(\Omega)$ Hilbert 结构，并定义 $H^{-1}(\Omega)=(H_0^1(\Omega))^*$。给定 $f\in H^{-1}(\Omega)$，求 $u\in H_0^1(\Omega)$ 满足
\[
 \int_\Omega\nabla u\cdot\nabla v\,\mathrm dx=f(v)
 \quad(v\in H_0^1(\Omega)).
\]
这里 $a(u,v)=\int\nabla u\cdot\nabla v$ 的有界常数和强制常数都为一，故 Lax--Milgram 定理给出唯一解及 $\|\nabla u\|_2\leq\|f\|_{H^{-1}}$。若 $f\in L^2$，将其视为 $v\mapsto\int fv$，Poincaré 不等式给出连续性。

若 $u$ 足够光滑，分部积分将该恒等式化为 $-\Delta u=f$ 及零边界条件；反之，弱解自动具有经典导数或逐点边界值并非形式推理的结果，还需要额外正则性条件。

\section{Galerkin 逼近与 Céa 估计}
取有限维子空间 $H_n\subset H$，求 $u_n\in H_n$ 使 $a(u_n,v_n)=f(v_n)$ 对所有 $v_n\in H_n$ 成立。限制后的形式仍有界且强制，所以离散解唯一。令 $e=u-u_n$，相减得到 Galerkin 正交性 $a(e,v_n)=0$，于是对任意 $w_n\in H_n$，
\[
 \alpha\|e\|^2\leq a(e,e)=a(e,u-w_n)
 \leq M\|e\|\|u-w_n\|.
\]
因此
\begin{equation}
 \|u-u_n\|\leq\frac M\alpha\inf_{w_n\in H_n}\|u-w_n\|.
 \label{eq:fa-cea}
\end{equation}
这称为 Céa 估计。若右侧最佳逼近误差趋零，则离散解收敛；单凭“子空间有限维”不能推出该逼近性质。

取 $H_n$ 的基 $\phi_1,\ldots,\phi_m$，令 $u_n=\sum_jc_j\phi_j$，则待解的有限维系统为 $Ac=b$，其中 $A_{ij}=a(\phi_j,\phi_i)$、$b_i=f(\phi_i)$。如果 $a$ 对称，$A$ 对称正定；若 $a$ 不对称，不能称矩阵对称，但强制性仍保证齐次系统只有零解，因此有限维系统可逆。Céa 估计把离散求解的误差归结为所选子空间的逼近能力；它本身不指定逼近误差关于网格尺度的收敛阶，那还取决于解的光滑性和基函数的选择。

\section{例题与习题}
\begin{example}
方程 $-u''=1$、$u(0)=u(1)=0$ 的弱解为 $u(t)=t(1-t)/2$。
\end{example}
\begin{solve}
此函数属于 $H_0^1(0,1)$。先对光滑紧支撑 $v$ 分部积分得 $\int_0^1u'v'=\int_0^1v$，再由稠密性扩展到所有 $v\in H_0^1$，故它是弱解；唯一性由 Lax--Milgram 定理保证。若用分片线性基 $\phi_j$ 逼近，系数满足 $\sum_j(\int\phi_j'\phi_i')c_j=\int\phi_i$，刚度矩阵的正定性来自梯度能量。
\end{solve}

例如将区间分成三段，节点为 $0,1/3,2/3,1$，取在两个内部节点分别为一、其余节点为零的连续分片线性帽函数 $\phi_1,\phi_2$。每段斜率为 $3$、$-3$ 或零，直接积分得到
\[
 \begin{pmatrix}6&-3\\-3&6\end{pmatrix}
 \begin{pmatrix}c_1\\c_2\end{pmatrix}
 =\begin{pmatrix}1/3\\1/3\end{pmatrix},
 \qquad c_1=c_2=1/9.
\]
因此近似解是依次经过 $(0,0),(1/3,1/9),(2/3,1/9),(1,0)$ 的折线。本例的节点值恰等于精确解的节点值，但两个节点之间的抛物线与折线仍有差别；节点上的相等不代表两函数处处相等。梯度范数下的误差由式\eqref{eq:fa-cea}控制，这才是相应函数空间中的收敛含义。
\begin{enumerate}
 \item 在 $H_0^1(0,1)$ 上证明 $\|u\|_2\leq\|u'\|_2$，不要求求最佳常数。
 \item 用能量差证明对称强制形式的极小点唯一。
 \item 说明若改在 $H^1(0,1)$ 上只使用 $a(u,v)=\int u'v'$，常数函数如何导致不唯一。
 \item\label{exer:fa-galerkin} 对嵌套且并集稠密的 $H_n$，由式\eqref{eq:fa-cea}证明 $u_n\to u$。
\end{enumerate}

\chapter{不动点方法与非线性方程}
\label{chap:fa-fixed-point}

线性方程可以研究核、值域与谱；非线性方程常改写为 $u=F(u)$。不动点方法的第一步是选定空间与不变集合，随后才判断压缩性、连续性或紧性。

\section{压缩映射与非线性积分方程}
考虑 $C([0,1])$ 上的方程
\[
 u(t)=g(t)+\lambda\int_0^1k(t,s)\phi(u(s))\,\mathrm ds,
\]
其中 $g,k$ 连续，$\phi:\mathbb R\to\mathbb R$ 为 Lipschitz 函数，常数为 $L$。右侧定义连续函数，且对应映射 $F$ 满足
\[
 \|Fu-Fv\|_\infty\leq q\|u-v\|_\infty,
 \qquad q=|\lambda|L\max_t\int_0^1|k(t,s)|\,\mathrm ds.
\]
若 $q<1$，定理\ref{thm:fa-contraction}给出唯一解。算法为任取 $u_0\in C([0,1])$，计算 $u_{n+1}=F(u_n)$；可用后验误差界 $q\|u_n-u_{n-1}\|_\infty/(1-q)$ 作为停止依据。

这个误差界针对精确计算 $F(u_n)$ 得到的迭代。若每步只算到误差 $\delta_n$，即 $\|\widetilde u_{n+1}-F(\widetilde u_n)\|_\infty\leq\delta_n$，则相对于真解 $u_*$，
\[
 \|\widetilde u_{n+1}-u_*\|_\infty
 \leq q\|\widetilde u_n-u_*\|_\infty+\delta_n.
\]
若始终有 $\delta_n\leq\delta$，递推给出 $\|\widetilde u_n-u_*\|_\infty\leq q^n\|\widetilde u_0-u_*\|_\infty+\delta/(1-q)$。所以实际积分计算的误差会留下一个误差上界；要保证逼近真解的精度任意提高，需要进一步控制这些计算误差。

\begin{example}[一个既能用压缩性又能用紧性的方程]\label{ex:fa-nonlinear-integral}
考虑 $u(t)=t+\lambda\int_0^1ts\sin u(s)\,\mathrm ds$，其中 $\lambda\in\mathbb R$。当 $|\lambda|<2$ 时有唯一连续解；对任意实数 $\lambda$ 至少有一个连续解。
\end{example}
\begin{solve}
因为 $|\sin a-\sin b|\leq|a-b|$，且 $\max_{0\leq t\leq1}\int_0^1ts\,\mathrm ds=1/2$，压缩常数可取 $q=|\lambda|/2$，给出第一项结论。为考察任意参数，注意每个解都必须形如 $u(t)=ct$，其中
\[
 c=1+\lambda\int_0^1s\sin(cs)\,\mathrm ds.
\]
右端作为 $c$ 的连续函数，始终落在区间 $[1-|\lambda|/2,1+|\lambda|/2]$ 中。在这个闭区间的左端，右端减 $c$ 非负；在右端非正，由介值定理得到一个不动点，因而得到连续解。后面的 Schauder 定理将把这种“连续映射保持一个闭有界凸集，且输出足够紧”的思想推广到不能化成一个标量的问题。
\end{solve}

\section{常微分方程的局部存在性}
考虑 $u'(t)=b(t,u(t))$、$u(t_0)=u_0$，其中 $b$ 在矩形 $|t-t_0|\leq a$、$|u-u_0|\leq r$ 上连续，并对第二变量统一 Lipschitz，常数为 $L$。设该矩形上 $|b|\leq M$，选 $h\leq a$ 使 $hM\leq r$ 且 $hL<1$。在完备闭集
\[
 C=\{u\in C([t_0-h,t_0+h]):\|u-u_0\|_\infty\leq r\}
\]
上，$Fu(t)=u_0+\int_{t_0}^t b(s,u(s))\,\mathrm ds$ 保持 $C$ 不变且为压缩。其不动点连续，积分表达又给出可微性，因此是微分方程的局部唯一解。延长到更大时间区间需要进一步控制解的范围，不能仅凭局部估计断言全局存在。

这里 $hM\leq r$ 用于确保近似曲线不离开已知 $b$ 的区域，$hL<1$ 用于确保两条近似曲线的距离收缩，二者解决不同问题。例如 $u'=u$、$u(0)=1$ 的 Picard 迭代从常值函数 $u_0(t)=1$ 开始，依次得到 $u_n(t)=\sum_{k=0}^nt^k/k!$，极限为 $e^t$。短区间上的压缩估计证明局部收敛，而级数的进一步估计还能证明任意有限区间上的收敛。相反，$u'=u^2$、$u(0)=1$ 的局部解 $u(t)=1/(1-t)$ 在 $t\uparrow1$ 时发散，说明局部 Lipschitz 条件不能单独保证全局解。

\section{Schauder 不动点定理}
本节引用有限维拓扑中的 Brouwer 不动点定理：有限维实线性空间中，非空紧凸集的每个连续自映射都有不动点。其拓扑证明超出本书范围，可参阅\cite{teschl-fa}的 Brouwer 映射度部分；下面给出从它推出无限维紧映射结论的完整论证。

\begin{theorem}[Schauder]\label{thm:fa-schauder}
设 $C$ 是 Banach 空间中的非空闭有界凸集，$F:C\to C$ 连续且 $F(C)$ 相对紧，则 $F$ 有不动点。
\end{theorem}
\begin{proof}
记 $K=\overline{F(C)}\subset C$。对 $\varepsilon>0$ 选 $y_1,\ldots,y_m\in K$ 构成 $\varepsilon/2$ 网。在 $K$ 上定义
\[
 w_j(y)=\max(\varepsilon-\|y-y_j\|,0),\qquad
 Q_\varepsilon(y)=\frac{\sum_jw_j(y)y_j}{\sum_jw_j(y)}.
\]
分母严格为正，映射连续，且 $\|Q_\varepsilon(y)-y\|<\varepsilon$。凸包 $D=\operatorname{conv}\{y_1,\ldots,y_m\}$ 为 $C$ 内有限维紧凸集，$Q_\varepsilon F:D\to D$ 连续，所以有不动点 $x_\varepsilon$。于是 $\|x_\varepsilon-Fx_\varepsilon\|<\varepsilon$。

令 $\varepsilon\downarrow0$。由于 $Fx_\varepsilon\in K$，可选择一列使 $Fx_\varepsilon$ 收敛，$x_\varepsilon$ 也趋于同一极限 $x\in C$。由连续性 $Fx=x$。
\end{proof}
紧性提供的是可收敛的近似解，不提供压缩映射所具有的误差收缩率。

证明中的有限维凸包是关键的过渡：$Q_\varepsilon$ 用有限多个代表点的加权平均近似紧集中的任意输出，把无限维映射变成有限维映射，再调用 Brouwer 定理。凸性保证这些平均点仍在允许集合内；相对紧性保证可以用有限个代表点控制全部输出，并从近似不动点提取极限；连续性保证极限真正满足方程。这些条件各自承担了明确的一步。

\section{先验估计与 Schaefer 定理}
非线性映射 $F:X\to X$ 称为紧映射，若它连续并把有界集映成相对紧集。
\begin{theorem}[Schaefer 不动点定理]
若 Banach 空间上的紧映射 $F$ 满足集合
\[
 \{x\in X:x=\lambda F(x)\text{ 对某个 }0\leq\lambda\leq1\}
\]
有界，则 $F$ 有不动点。
\end{theorem}
\begin{proof}
取 $R$ 大于上述集合中所有元素的范数。在闭球 $RB_X$ 上定义 $G(x)=F(x)$（当 $\|F(x)\|\leq R$），否则定义 $G(x)=RF(x)/\|F(x)\|$。它连续、像相对紧，并将该闭球映入自身。由 Schauder 定理有 $x=G(x)$。若截断发生，则 $\|x\|=R$ 且 $x=\lambda F(x)$，$0<\lambda<1$，与 $R$ 的选取矛盾。因此没有截断，$x=F(x)$。
\end{proof}
这解释了先验估计的用途：不必事先找到所有解，只需给整个参数族的候选解一个统一界。

\section{存在性、唯一性与算法收敛}
$F(x)=1-x$ 在 $[0,1]$ 上连续，唯一不动点为 $1/2$，但除初值 $1/2$ 外，简单迭代在两个数之间交替。$F(x)=x$ 的每个点都是不动点。因此，Schauder 结论不包含唯一性，也不保证迭代收敛。研究算法时，必须另外给出收缩、单调性、能量下降或其他收敛依据。

\section{例题与习题}
\begin{example}[仅连续非线性的存在性]
设 $g\in C([0,1])$、$k\in C([0,1]^2)$，$\phi:\mathbb R\to\mathbb R$ 连续且 $|\phi(s)|\leq A$。则方程 $u(t)=g(t)+\int_0^1k(t,s)\phi(u(s))\,\mathrm ds$ 至少有一个连续解。
\end{example}
\begin{solve}
取 $R\geq\|g\|_\infty+A\max_t\int|k(t,s)|\,\mathrm ds$，右端映射保持闭球 $\|u\|_\infty\leq R$。$\phi$ 在 $[-R,R]$ 上一致连续，保证该映射连续；核的一致连续性给出像的等度连续性，一致有界性已由 $R$ 的选取保证。Arzelà--Ascoli 与 Schauder 定理给出解。这里不要求 $\phi$ 为压缩型非线性。
\end{solve}
\begin{enumerate}
 \item 将例\ref{ex:fa-nonlinear-integral}的 $\sin u$ 改为 $\tanh u$，证明相同的唯一性充分条件仍成立，并证明任意实参数下至少有解。
 \item 在常微分方程例中，明确 $hM\leq r$ 与 $hL<1$ 分别用于哪一步。
 \item 给出连续自映射有多个不动点、且某些初值的迭代不收敛的区间例子。
 \item 检查 Schauder 证明中闭性、凸性、连续性和紧性分别出现的位置。
\end{enumerate}

\chapter{无界算子与定义域}
\label{chap:fa-unbounded}

本章为选读内容。求导在 $L^2$ 范数下不能由函数本身的大小控制，因此不适合放进处处定义的有界算子框架。正确做法是把定义域一并写入算子的定义。

\section{稠密定义、闭性与可闭性}
记线性算子为 $A:D(A)\subset H\to H$，其中 $D(A)$ 为线性子空间。若 $\overline{D(A)}=H$，称其\keyterm{稠密定义（densely defined）}；若图像 $G(A)$ 在 $H\times H$ 中闭，称为\keyterm{闭算子（closed operator）}。等价的序列条件是：$x_n\in D(A)$、$x_n\to x$、$Ax_n\to y$ 时，必有 $x\in D(A)$ 且 $Ax=y$。

若图像闭包仍为某个算子的图像，则称 $A$ 可闭（closable），该算子记为 $\overline A$。一个线性图像表示单值算子的充要条件是它与竖直子空间 $\{0\}\times H$ 仅相交于零，所以可闭性等价于
\[
 x_n\to0,\quad Ax_n\to y\quad\Longrightarrow\quad y=0.
\]

\section{图像范数}
“闭算子”不是说定义域本身在 $H$ 中闭，也不是说算子把闭集映成闭集。它要求输入和输出一起收敛时，极限仍在同一张图像上。只知道 $x_n\to x$ 并不能要求 $Ax_n$ 收敛；但若后者也收敛，就不允许出现一个与算子关系不符的极限。这正是微分方程中同时控制函数与导数时需要的性质。

\section{图像范数}
在 $D(A)$ 上定义 $\|x\|_A=(\|x\|^2+\|Ax\|^2)^{1/2}$，称为图像范数（graph norm）。映射 $x\mapsto(x,Ax)$ 是到 $G(A)$ 的等距同构，因此 $A$ 闭当且仅当 $D(A)$ 在图像范数下完备。

\begin{example}\label{ex:fa-unbounded-diagonal}
在 $\ell^2$ 上，令 $(Ax)_n=nx_n$，$D(A)=\{x\in\ell^2:(nx_n)\in\ell^2\}$。它稠密定义、闭而无界。
\end{example}
\begin{solve}
$D(A)$ 包含 $c_{00}$，所以稠密。若 $x^{(k)}\to x$、$Ax^{(k)}\to y$，逐坐标取极限得 $y_n=nx_n$，于是 $x\in D(A)$ 且 $Ax=y$，证明闭性。又 $\|Ae_n\|=n$，故它在原 $\ell^2$ 范数下无界。图像范数则为 $(\sum_n(1+n^2)|x_n|^2)^{1/2}$。
\end{solve}

取 $x=(1/n)\in\ell^2$，它的有限截断都在 $D(A)$ 中，却按 $\ell^2$ 范数趋于不属于 $D(A)$ 的 $x$，所以定义域确实不闭。相应输出的前 $N$ 个坐标均为一，其范数为 $\sqrt N$，没有收敛；因此这不违反算子的闭性。图像范数同时惩罚输入与输出的大小，在这个更强的范数下，刚才的截断序列不是 Cauchy 序列。

作为 $D(A)$ 配图像范数到 $H$ 的映射，任何这样的 $A$ 都满足 $\|Ax\|\leq\|x\|_A$。所以“无界”始终指相对于原来的环境空间范数而言，并不排斥通过改变定义域范数获得一个有界映射；改变范数同时也改变了对输入数据精度的要求。

\section{无界算子的伴随}
若 $A$ 稠密定义，规定 $y\in D(A^*)$ 当且仅当存在 $C_y<\infty$ 使
\[
 |\langle Ax,y\rangle|\leq C_y\|x\|\quad(x\in D(A)).
\]
该泛函先唯一延拓到 $H$，再由 Riesz 表示得到唯一 $A^*y$，使 $\langle Ax,y\rangle=\langle x,A^*y\rangle$。

这里估计右端必须是 $\|x\|$，不能擅自改成图像范数。对任意固定 $y$，Cauchy--Schwarz 总会给出 $|\langle Ax,y\rangle|\leq\|y\|\|Ax\|\leq\|y\|\|x\|_A$，但这样的估计并不能保证存在 $H$ 中的向量代表该泛函。伴随定义域正是挑出那些可以仅由原输入范数控制的测试向量。

\begin{proposition}
伴随 $A^*$ 总是闭算子；$A$ 可闭当且仅当 $D(A^*)$ 稠密，此时 $\overline A=A^{**}$。
\end{proposition}
\begin{proof}
若 $y_n\to y$、$A^*y_n\to z$，则对任意 $x\in D(A)$ 取极限，得到 $\langle Ax,y\rangle=\langle x,z\rangle$，故 $y\in D(A^*)$ 且 $A^*y=z$。

在乘积 Hilbert 空间中直接由定义得到
\[
 G(A)^\perp=\{(-A^*y,y):y\in D(A^*)\}.
\]
所以 $(0,z)\in\overline{G(A)}=G(A)^{\perp\perp}$ 当且仅当 $z\perp D(A^*)$。图像闭包不含非零竖直向量，恰好等价于 $D(A^*)$ 稠密。此时再次使用同一正交关系，得到 $G(A^{**})=\overline{G(A)}$。
\end{proof}

\section{对称、自伴与本质自伴}
稠密定义的 $A$ 称为对称（symmetric），若 $\langle Ax,y\rangle=\langle x,Ay\rangle$ 对 $x,y\in D(A)$ 成立，即 $A\subset A^*$；称为自伴，若 $D(A)=D(A^*)$ 且两算子在其上相等。若 $A$ 可闭且 $\overline A$ 自伴，称为本质自伴（essentially self-adjoint）。对称性只检验共同定义域上的恒等式，自伴性还要求定义域已经足够大。

\begin{theorem}[Hellinger--Toeplitz]
在整个 Hilbert 空间上定义的对称线性算子必有界。
\end{theorem}
\begin{proof}
设 $x_n\to x$、$Ax_n\to y$。对任意 $z\in H$，由对称性得 $\langle y,z\rangle=\lim_n\langle x_n,Az\rangle=\langle Ax,z\rangle$，因此 $y=Ax$。图像闭，且定义域为整个 $H$，由闭图像定理得到有界性。
\end{proof}
所以真正无界的自伴算子必有真定义域。例\ref{ex:fa-unbounded-diagonal}还是自伴的：用 $x=e_n$ 测试伴随关系，迫使 $(A^*y)_n=ny_n$，其属于 $\ell^2$ 恰好给出同一个定义域。

\section{边界条件决定算子}
在一维，$H^1(0,1)$ 的元素有唯一的连续绝对连续代表，满足 $u(t)=u(0)+\int_0^tu'(s)\,\mathrm ds$。这是因为减去弱导数的积分后，剩余函数弱导数为零，故为常数。因而端点值有意义，且可进行分部积分。

端点取值在 $H^1$ 范数下连续。例如对任意 $t\in[0,1]$，$|u(0)|\leq|u(t)|+\int_0^1|u'|$，对 $t$ 积分并用 Cauchy--Schwarz 得 $|u(0)|\leq\|u\|_2+\|u'\|_2$，另一个端点同理。由此，$H_0^1$ 中的函数端点为零；反向结论在区间上由零延拓、光滑逼近及端点截断得到。因此本节可把 $H_0^1(0,1)$ 识别为端点均为零的 $H^1$ 函数。这个端点识别是一维结论；一般区域的边界值需要迹理论，不能对任意 $L^2$ 函数直接使用。

在复 $L^2(0,1)$ 中考虑表达式 $-\mathrm i u'$，有边界恒等式
\[
 \langle-\mathrm iu',v\rangle-\langle u,-\mathrm iv'\rangle
 =-\mathrm i[u\overline v]_0^1.
\]
若 $D(A_0)=H_0^1(0,1)$，两端点均为零，故 $A_0$ 对称。然而测试函数检验伴随只要求 $v\in H^1(0,1)$，边界项不再限制 $v$，因此 $D(A_0^*)=H^1(0,1)$，$A_0$ 不自伴。

若改为周期定义域 $D(A_{\mathrm{per}})=\{u\in H^1(0,1):u(0)=u(1)\}$，伴随关系则迫使 $v(0)=v(1)$，故 $A_{\mathrm{per}}$ 自伴。特征方程给出 $u(t)=ce^{\mathrm i\lambda t}$，周期条件给出 $\lambda=2\pi n$，$n\in\mathbb Z$。同一微分表达式配上不同定义域，确实产生不同算子。

\section{习题}
\begin{enumerate}
 \item 证明无界对角算子的定义域在 $\ell^2$ 中不闭，却在图像范数下完备。
 \item 对 $A(x_n)=(a_nx_n)$，其中 $a_n$ 为任意实数列，写出最大自然定义域并证明自伴性。
 \item 核对周期微分算子的伴随定义域，说明为何只验证分部积分恒等式还不够。
 \item 构造稠密定义但不可闭的线性算子。提示：在 $c_{00}\subset\ell^2$ 上考虑 $Ax=(\sum_nnx_n)e_1$。
\end{enumerate}

\chapter{一般谱定理与函数演算}
\label{chap:fa-general-spectral}

本章为选读内容。乘法算子可以没有任何特征向量，因而一般自伴算子不能只靠特征值求和表示。谱积分用可测集合上的正交投影代替单个特征空间。本章证明连续函数演算的基本构造，对一般谱定理给出准确陈述和证明思路；完整论证可参阅文献\cite{teschl}的谱定理部分。

\section{投影值测度}
先看矩阵 $T=\operatorname{diag}(1,2,2)$。令 $P_1$ 投影到第一坐标轴，$P_2$ 投影到后两个坐标组成的平面，就有 $T=P_1+2P_2$。对任意数集 $B$，令 $E(B)$ 保留那些特征值落在 $B$ 中的坐标。例如 $E(\{2\})=P_2$，$E(\{1,2\})=I$，$E(\{3\})=0$。于是 $T=\int\lambda\,\mathrm dE(\lambda)$ 在这里就是一个有限和。无限维谱积分沿用同样的思想，只是不能保证全部信息都集中在有限或可数个单点上。

\begin{definition}
设 $S$ 是实轴或复平面的 Borel 子集。\keyterm{投影值测度（projection-valued measure）}是从 $S$ 的 Borel 集到 $H$ 上正交投影的映射 $E$，满足 $E(S)=I$、$E(\varnothing)=0$、$E(B\cap C)=E(B)E(C)$，且对两两不交的 $B_n$，
\[
 E\Bigl(\bigcup_nB_n\Bigr)x=\sum_nE(B_n)x\quad(x\in H),
\]
其中级数在 $H$ 中收敛。
\end{definition}
给定 $x$，$\mu_x(B)=\langle E(B)x,x\rangle=\|E(B)x\|^2$ 为有限正测度。用 $\mathbf1_B$ 表示在 $B$ 上取一、在其外取零的示性函数。对简单函数 $s=\sum_ja_j\mathbf1_{B_j}$，定义 $s(E)=\sum_ja_jE(B_j)$。若 $B_j$ 两两不交，则
\[
 \|s(E)x\|^2=\sum_j|a_j|^2\mu_x(B_j)=\int|s|^2\,\mathrm d\mu_x.
\]
有界 Borel 函数可由简单函数一致逼近，这个等式保证相应算子在范数下收敛，定义 $f(E)=\int f\,\mathrm dE$。

强可数可加性表示对每个固定向量 $x$，各个谱区间中的分量可以在范数下相加；它通常不是算子范数意义的可数可加性。例如在 $\ell^2$ 中，第 $j$ 个坐标投影记为 $P_j$，则 $\sum_{j=1}^NP_jx\to x$ 对每个 $x$ 成立，但 $\|I-\sum_{j=1}^NP_j\|=1$。因此定义中明确写出向量 $x$ 与收敛意义，是不可省略的部分。

\section{连续函数演算}
先设 $T$ 有界自伴。对多项式 $p$，$p(T)$ 正规。正规算子 $N$ 满足 $\|N^2\|=\|N\|^2$：计算
\[
 \|N^2\|^2=\|(N^*)^2N^2\|=\|(N^*N)^2\|=\|N^*N\|^2,
\]
最后一步用自伴算子的范数恒等式。重复取平方并用谱半径公式，得 $\|N\|=r(N)$。因此多项式谱映射定理给出
\[
 \|p(T)\|=\max_{\lambda\in\sigma(T)}|p(\lambda)|.
\]

对 $f\in C(\sigma(T))$，由实紧集上的 Weierstrass 逼近定理，取多项式 $p_n$ 在谱上一致趋于 $f$；定义 $f(T)=\lim_np_n(T)$。上式保证极限存在、与逼近选择无关，并且
\[
 (fg)(T)=f(T)g(T),\quad \overline f(T)=f(T)^*,\quad
 \|f(T)\|=\max_{\sigma(T)}|f|.
\]
这称为\keyterm{连续函数演算（continuous functional calculus）}。对复平面上的正规算子需要同时使用 $z,\overline z$ 的多项式以及 Stone--Weierstrass 定理，不能只用解析多项式逼近所有连续函数。

函数演算的实际含义是把数值函数作用到每个谱方向上。在有限维对角情形，$f(T)=\operatorname{diag}(f(\lambda_1),\ldots,f(\lambda_m))$；一般情形则用上述极限保持同样的运算规律。若 $0\notin\sigma(T)$，函数 $f(\lambda)=1/\lambda$ 在谱上连续，乘法性质给出 $f(T)T=I$，于是恢复逆算子。若 $T\geq0$，同样可以用 $f(\lambda)=\sqrt\lambda$ 构造平方根。

谱区间的示性函数通常在谱上不连续，因此不能只靠连续函数演算就直接代入 $\mathbf1_B$。一般谱定理提供的 Borel 函数演算，正是为了同时容纳这类谱投影；它是理论中的进一步扩展，而非把连续逼近步骤原样重复一遍。

\section{一般有界谱定理}
\begin{theorem}[谱定理；本节引用]\label{thm:fa-general-spectral}
每个复 Hilbert 空间上的有界正规算子 $T$ 都有唯一的投影值测度 $E$，定义在 $\sigma(T)$ 上，使
\begin{equation}
 T=\int_{\sigma(T)}\lambda\,\mathrm dE(\lambda).
 \label{eq:fa-general-spectral-integral}
\end{equation}
有界 Borel 函数演算由 $f(T)=\int f\,\mathrm dE$ 给出。若 $T$ 自伴，谱在实轴上。
\end{theorem}
\begin{proof}[证明思路]
先建立连续函数演算。对每个 $x$，$f\mapsto\langle f(T)x,x\rangle$ 是正线性泛函，用 Riesz--Markov 表示得到测度 $\mu_x$。对不同向量使用极化恒等式，恢复交叉项的测度；再证明这些测度给出算子 $E(B)$，且满足投影性、乘法性与强可数可加性，最后由简单函数逼近构造 Borel 演算。由标量测度恢复算子并验证投影性是关键步骤，此处仅说明路线，不作为完整证明。
\end{proof}
投影 $E(\{\lambda\})$ 的像恰为 $\lambda$ 的特征空间。紧自伴算子具有可数个非零谱原子，再加零处的核投影，其谱积分就是第\ref{chap:fa-compact-spectral}章的求和公式。若谱测度没有相应单点原子，谱点就不产生特征向量。

\section{无界自伴谱定理}
对无界算子，表达式必须与定义域同时给出。一般谱定理的无界版本断言：每个自伴算子 $A$ 有唯一实轴上的投影值测度 $E$，使
\[
 D(A)=\left\{x:\int_{\mathbb R}\lambda^2\,\mathrm d\mu_x(\lambda)<\infty\right\},
 \qquad Ax=\lim_{N\to\infty}\int_{[-N,N]}\lambda\,\mathrm dE(\lambda)x.
\]
这里的极限按向量范数取得。更一般地，有限值 Borel 函数 $f$ 给出
\[
 D(f(A))=\{x:\int|f|^2\,\mathrm d\mu_x<\infty\},\qquad
 \|f(A)x\|^2=\int|f|^2\,\mathrm d\mu_x.
\]
本书引用这一版本；它可由有界谱定理配合 Cayley 变换等方法建立，完整条件与证明见\cite{teschl}。无界函数演算中的算子乘法需要核对定义域，不能无条件照搬有界代数公式。

以 $\ell^2$ 上的 $Ax=(nx_n)$ 为例，谱投影 $E(B)$ 保留满足 $n\in B$ 的坐标，故 $\mu_x(B)=\sum_{n\in B}|x_n|^2$。积分条件立即变成 $\sum_n n^2|x_n|^2<\infty$，正是例\ref{ex:fa-unbounded-diagonal}中的定义域；而 $A^2$ 的定义域要求 $\sum_n n^4|x_n|^2<\infty$。例如 $x_n=1/n^2$ 属于 $D(A)$，却不属于 $D(A^2)$。所以即使 $Ax$ 已有定义，也不能自动再作用一次 $A$。

\section{平方根、谱投影与演化算子}
若 $T\geq0$ 有界自伴，则 $T^{1/2}=\int\sqrt\lambda\,\mathrm dE(\lambda)$ 为正算子且平方等于 $T$。其正平方根唯一：若 $B\geq0$ 且 $B^2=T$，对 $B$ 的连续函数演算应用恒等式 $\sqrt{s^2}=s$（$s\geq0$），即得 $B=(B^2)^{1/2}=T^{1/2}$。

若 $A$ 自伴，$e^{\mathrm itA}=\int e^{\mathrm it\lambda}\,\mathrm dE(\lambda)$ 为酉算子。由谱积分的乘法性质形成群；由 $|e^{\mathrm it\lambda}-1|^2\leq4$ 和支配收敛定理，对每个 $x$ 有 $e^{\mathrm itA}x\to x$。若 $A\geq0$，则 $e^{-tA}$（$t\geq0$）同样强连续，且范数至多一，是下一章的一个重要半群。

\section{例题与习题}
\begin{example}[乘法算子的谱投影]
在 $L^2(0,1)$ 上令 $Tf(t)=tf(t)$。对 Borel 集 $B\subset[0,1]$，$E(B)f=\mathbf1_Bf$，且 $g(T)f(t)=g(t)f(t)$。
\end{example}
\begin{solve}
示性函数乘法为正交投影，不交集合对应正交分量；可数可加性由 $L^2$ 收敛得到。简单函数逐项验证谱积分，再用一致逼近得到有界 Borel 函数公式。每个单点零测，所以 $E(\{\lambda\})=0$，这与算子没有特征值相符。
\end{solve}
\begin{enumerate}
 \item 对 $\ell^2$ 上的自伴对角算子写出 $E(B)$，由谱积分恢复其定义域。
 \item 验证非负乘法算子的平方根就是乘以原函数的平方根。
 \item 比较谱投影 $E(B)$ 与某个单点特征空间投影，说明二者信息量的差别。
 \item 证明 $e^{-tA}$ 在 $A\geq0$ 时保持范数不增，但不一定按算子范数趋于 $I$。
\end{enumerate}

\chapter{算子半群与演化方程}
\label{chap:fa-semigroups}

本章为选读内容。有限维方程 $u'(t)=Au(t)$ 的解由矩阵指数给出。无限维中，即使生成时间演化的算子 $A$ 无界，每个固定时刻的演化仍可能由有界算子描述。

\section{强连续半群}
\begin{definition}
Banach 空间 $X$ 上的\keyterm{强连续半群（strongly continuous semigroup）}是算子族 $S(t)\in\mathcal B(X)$，$t\geq0$，满足
\[
 S(0)=I,\qquad S(t+s)=S(t)S(s),
\]
并且对每个 $x\in X$，映射 $t\mapsto S(t)x$ 连续。若 $\|S(t)\|\leq1$，称为压缩半群。
\end{definition}
当 $A\in\mathcal B(X)$ 时，$S(t)=\sum_{n\geq0}t^nA^n/n!$ 按算子范数收敛；绝对收敛允许重排双重级数，从二项式公式得到半群性质，而且 $S'(t)=AS(t)$。无界情形不应直接写未经定义的算子幂级数。

任意强连续半群在有限时间区间上算子范数一致有界：对每个 $x$，轨道在该区间上有界，再用一致有界原理。取 $M=\max\{1,\sup_{0\leq t\leq1}\|S(t)\|\}$，由半群性质得到 $\|S(t)\|\leq M e^{\omega t}$，其中 $\omega=\log M$。

“半群”只要求非负时间，是因为演化未必能够稳定地倒退。矩阵指数 $e^{tA}$ 在有限维中总有逆 $e^{-tA}$；热方程的正向演化不断压低高频误差，倒退则要把这些坐标指数放大，通常不能得到整个 $L^2$ 空间上的有界逆。因此半群性质描述时间拼接，并不自动提供负时间。

\section{生成元与闭性}
\begin{definition}
半群的\keyterm{无穷小生成元（infinitesimal generator）}由
\begin{equation}
 Ax=\lim_{h\downarrow0}\frac{S(h)x-x}{h},\qquad
 D(A)=\left\{x\in X:\text{此极限在 }X\text{ 中存在}\right\}
 \label{eq:fa-semigroup-generator}
\end{equation}
定义。
\end{definition}
若 $x\in D(A)$，由半群性质可得 $S(t)x\in D(A)$，且 $AS(t)x=S(t)Ax$；右、左差商进一步给出 $\frac{\mathrm d}{\mathrm dt}S(t)x=S(t)Ax$（在 $t=0$ 取右导数）。因此
\begin{equation}
 S(t)x-x=\int_0^tS(s)Ax\,\mathrm ds.
 \label{eq:fa-generator-integral}
\end{equation}
连续 Banach 空间值函数的积分可定义为 Riemann 和的范数极限；一致连续性和完备性保证极限存在。

生成元的定义只使用固定初值 $x$ 的差商。在有界情形可以把极限统一为算子范数极限，得到处处定义的 $A$；在无界情形，某些初值的轨道只有连续性，没有初始导数。定义域 $D(A)$ 因而可以理解为允许在初始时刻写出微分方程的那一部分初值，而不是人为附加的技术集合。

\begin{proposition}
强连续半群的生成元稠密定义且闭。
\end{proposition}
\begin{proof}
对任意 $x\in X$，令 $x_h=h^{-1}\int_0^hS(s)x\,\mathrm ds$。利用半群性质移动积分区间可得 $x_h\in D(A)$ 且 $Ax_h=(S(h)x-x)/h$；强连续性给出 $x_h\to x$，故定义域稠密。

若 $x_n\in D(A)$、$x_n\to x$、$Ax_n\to y$，在式\eqref{eq:fa-generator-integral}中利用有限时间区间上的统一算子界取极限，得到 $S(t)x-x=\int_0^tS(s)y\,\mathrm ds$。除以 $t$ 并令 $t\downarrow0$，右侧趋于 $y$，所以 $x\in D(A)$ 且 $Ax=y$，即 $A$ 闭。
\end{proof}

\section{预解算子与 Hille--Yosida 定理}
对闭算子 $A$，$\lambda\in\rho(A)$ 指 $\lambda I-A:D(A)\to X$ 双射且其逆作为 $X\to X$ 的算子有界。若 $A$ 生成满足 $\|S(t)\|\leq Me^{\omega t}$ 的半群，则对实数 $\lambda>\omega$，
\[
 R(\lambda,A)x=\int_0^\infty e^{-\lambda t}S(t)x\,\mathrm dt,
 \qquad\|R(\lambda,A)\|\leq\frac M{\lambda-\omega}.
\]
积分由增长估计绝对收敛。其逆关系可先在 $D(A)$ 中用式\eqref{eq:fa-generator-integral}分部积分验证，再利用稠密性和生成元闭性扩展；这是半群与预解估计之间的联系。

\begin{theorem}[Hille--Yosida 的压缩半群版本；引用]
设 $A$ 是 Banach 空间上的稠密定义闭算子。它生成强连续压缩半群，当且仅当每个实数 $\lambda>0$ 属于 $\rho(A)$，且 $\|R(\lambda,A)\|\leq1/\lambda$。
\end{theorem}
必要性由上述积分公式得到。充分性通过有界近似 $A_\lambda=\lambda A R(\lambda,A)=\lambda^2R(\lambda,A)-\lambda I$ 及指数半群的极限构造证明；本书不展开这一极限论证，完整证明见\cite{engel-nagel}第二章第三节。一般增长常数 $M$ 的版本需要对预解算子的各次幂给出相应估计，不能只替换单次估计中的常数。

\section{经典解与温和解}
设 $A$ 生成 $S(t)$。若 $u_0\in D(A)$，则 $u(t)=S(t)u_0$ 是 $u'=Au$ 的经典解，即轨道可微、取值在 $D(A)$ 中并满足微分方程。唯一性可固定终点 $t$，对 $s\mapsto S(t-s)u(s)$ 求导，得到导数为零。

一般 $u_0\in X$ 未必属于定义域，仍可定义 $u(t)=S(t)u_0$，称为\keyterm{温和解（mild solution）}。对连续外力 $g:[0,T]\to X$，非齐次问题的温和解为
\[
 u(t)=S(t)u_0+\int_0^tS(t-s)g(s)\,\mathrm ds.
\]
连续外力保证这条公式有意义，但要逐项微分并取得经典解，还需初值与外力的额外正则性；不能把温和解公式本身当作经典可微性的证明。

\section{区间上的热半群}
采用例\ref{ex:fa-green}中的正交基 $e_n(x)=\sqrt2\sin(n\pi x)$。若 $u_0=\sum_nc_ne_n\in L^2(0,1)$，令
\[
 S(t)u_0=\sum_{n=1}^{\infty}e^{-(n\pi)^2t}c_ne_n.
\]
系数逐项满足乘法规律，故这是半群；Parseval 恒等式与级数的支配收敛给出强连续性，且 $\|S(t)\|=e^{-\pi^2t}$。其生成元为
\[
 Au=-\sum_n(n\pi)^2c_ne_n,\qquad
 D(A)=\left\{u:\sum_n(n\pi)^4|c_n|^2<\infty\right\}.
\]
定义域内，差商由 $|(e^{-ah}-1)/h|\leq a$ 控制，故按 $L^2$ 范数收敛；反之，若差商收敛，逐坐标极限迫使所得系数平方可和，恰好给出该定义域。

定义 $H^2(0,1)=\{u\in H^1(0,1):u''\in L^2(0,1)\}$，二阶导数仍为弱导数。由 Green 算子的逆关系可将上述定义域识别为 $H^2(0,1)\cap H_0^1(0,1)$，并有 $Au=u''$。所以 $u(t)=S(t)u_0$ 描述零边界条件下的热方程 $u_t=u_{xx}$。每个 $t>0$ 的指数衰减还保证 $u(t)\in D(A)$；在初始时刻是否可微则取决于 $u_0$ 是否属于 $D(A)$。

\begin{example}[两个频率的衰减速度]
若 $u_0=e_1+2e_3$，则
\[
 u(t)=e^{-\pi^2t}e_1+2e^{-9\pi^2t}e_3,
 \qquad\|u(t)\|_2^2=e^{-2\pi^2t}+4e^{-18\pi^2t}.
\]
\end{example}
\begin{solve}
第一式逐项应用半群的定义，第二式由正交性得到。第三频率的衰减指数是第一频率的九倍，所以较晚时刻主要剩下第一项。在无穷多个频率同时存在时，对每个固定初值，尾部系数的平方和已经很小，故 $S(t)u_0\to u_0$；但对每个 $t>0$，总能找到足够大的 $n$，使 $1-e^{-(n\pi)^2t}$ 任意接近一，所以 $\|S(t)-I\|=1$。强连续与算子范数连续在这个例子中的差别，来自是否允许测试向量随时间改变。
\end{solve}

更定量地，对固定 $t>0$，$\sup_{r\geq0}r^2e^{-2rt}<\infty$，取 $r=(n\pi)^2$ 得
\[
 \sum_n(n\pi)^4e^{-2(n\pi)^2t}|c_n|^2
 \leq\left(\sup_{r\geq0}r^2e^{-2rt}\right)\sum_n|c_n|^2<\infty.
\]
这证明任意 $L^2$ 初值在正时间后进入生成元定义域。相反，若要从时刻 $t$ 的系数倒推初值，需要乘以 $e^{(n\pi)^2t}$，这些放大倍数没有统一上界，不能得到有界的反向演化算子。

\section{习题}
\begin{enumerate}
 \item 对有界 $A$ 证明 $\|e^{tA}\|\leq e^{t\|A\|}$。
 \item 在上述热半群中证明 $\|S(t)-I\|=1$ 对每个 $t>0$ 成立，说明强连续不等于算子范数连续。
 \item 设 $u_0\perp e_1,\ldots,e_{m-1}$，证明 $\|S(t)u_0\|_2\leq e^{-(m\pi)^2t}\|u_0\|_2$，并说明常数的最优性。
 \item 证明非负自伴算子的谱演算产生强连续压缩半群 $e^{-tA}$。
\end{enumerate}

\appendix

\chapter{测度与积分的必要知识}
\label{app:fa-measure}

本附录补充 $L^p$ 理论中直接用到的概念与估计。测度构造、一般测度表示和乘积测度理论只陈述所需结果；它们可在实分析课程中系统学习，相关测度、积分、对偶表示及 Sobolev 空间内容参阅\cite{teschl-ra}。

\section{测度、可测函数与几乎处处}
集合 $\Omega$ 上的 $\sigma$ 代数 $\mathcal A$ 包含 $\Omega$，并对补集与可数并封闭。测度 $\mu:\mathcal A\to[0,\infty]$ 满足 $\mu(\varnothing)=0$，且对两两不交的 $E_n$，$\mu(\bigcup_nE_n)=\sum_n\mu(E_n)$。若 $\Omega$ 可由可数个有限测度集覆盖，称测度 $\sigma$ 有限。

实函数 $f$ 可测，是指 $\{x:f(x)>a\}\in\mathcal A$ 对每个实数 $a$ 成立；复函数的实部与虚部均可测时称可测。某性质在一个零测集之外成立，称其几乎处处（almost everywhere）成立。零测集的任意子集仍可测时，称测度完备。欧氏空间的 Borel $\sigma$ 代数由开集生成；Lebesgue 可测集还包括对 Borel 测度作完备化时加入的零测子集。

\section{积分与收敛定理}
非负简单函数写为 $s=\sum_{j=1}^ma_j\mathbf1_{E_j}$，其中 $a_j\geq0$ 且 $E_j$ 两两不交，定义 $\int s\,\mathrm d\mu=\sum_ja_j\mu(E_j)$，约定 $0\cdot\infty=0$。对非负可测 $f$，定义
\[
 \int f\,\mathrm d\mu=\sup\left\{\int s\,\mathrm d\mu:0\leq s\leq f,\ s\text{ 为简单函数}\right\}.
\]
对实函数使用正负部分之差；当 $\int|f|<\infty$ 时称可积，复函数再分别积分实部与虚部。

\begin{theorem}[三个收敛工具]\label{thm:fa-measure-convergence}
以下函数均可测，等式与不等式允许非负积分为无穷。
\begin{enumerate}
 \item 若 $0\leq f_n\uparrow f$ 几乎处处，则 $\int f_n\uparrow\int f$（单调收敛定理）。
 \item 若 $f_n\geq0$，则 $\int\liminf_nf_n\leq\liminf_n\int f_n$（Fatou 引理）。
 \item 若 $f_n\to f$ 几乎处处，且 $|f_n|\leq g$，其中 $g$ 非负可积，则 $\int|f_n-f|\to0$，从而 $\int f_n\to\int f$（支配收敛定理）。
\end{enumerate}
\end{theorem}
\begin{proof}
单调情形中，$\lim_n\int f_n\leq\int f$。反向任取简单函数 $0\leq s\leq f$ 与 $0<\alpha<1$，在 $\{s>0\}$ 上集合 $E_n=\{f_n\geq\alpha s\}$ 递增并覆盖该集合。由测度的从下连续性，$\int\alpha s\mathbf1_{E_n}\to\alpha\int s$，而每一项不超过 $\int f_n$。令 $\alpha\uparrow1$ 再对 $s$ 取上确界得到等式。

令 $h_n=\inf_{k\geq n}f_k$，则 $h_n\uparrow\liminf f_n$ 且 $\int h_n\leq\inf_{k\geq n}\int f_k$，由单调收敛得到 Fatou 引理。支配情形中 $|f|\leq g$，对非负函数 $2g-|f_n-f|$ 用 Fatou 引理，得到 $\limsup_n\int|f_n-f|\leq0$。
\end{proof}
对两个 $\sigma$ 有限测度空间，本书引用 Tonelli--Fubini 定理：非负可测函数可交换积分次序，允许结果无穷；若函数的绝对值在乘积测度下可积，则两次迭代积分存在并等于乘积积分。交换极限、积分或求和前，应先说明使用的是哪一项条件。

\begin{example}[为什么极限与积分不能随意交换]
在 $(0,1)$ 上，$f_n(t)=n\mathbf1_{(0,1/n)}(t)$ 几乎处处趋于零，但 $\int_0^1f_n=1$ 对每个 $n$ 成立。
\end{example}
\begin{solve}
固定 $t>0$ 后，只要 $n>1/t$ 就有 $f_n(t)=0$；积分却等于高度乘宽度 $n\cdot(1/n)=1$。这里质量越来越集中在零点附近，缺少一个统一的可积控制函数。若存在 $g\in L^1$ 支配全部 $f_n$，支配收敛定理会给出积分趋零，矛盾。因此逐点极限只观察每个固定位置，不能单独排除移动或集中的质量。Fatou 引理在本例仍给出正确的不等式 $0\leq1$，但并不声称等号。
\end{solve}

\section{Hölder 与 Minkowski 不等式}
\begin{proposition}\label{prop:fa-holder-minkowski}
若 $1\leq p,q\leq\infty$、$1/p+1/q=1$，则对 $f\in L^p(\mu)$、$g\in L^q(\mu)$，
\[
 \int|fg|\,\mathrm d\mu\leq\|f\|_p\|g\|_q.
\]
对 $1\leq p\leq\infty$ 及 $f,g\in L^p(\mu)$，还成立 $\|f+g\|_p\leq\|f\|_p+\|g\|_p$。
\end{proposition}
\begin{proof}
当 $1<p<\infty$ 时，Young 不等式 $ab\leq a^p/p+b^q/q$ 可通过对 $a$ 的一元极值计算证明。把 $f,g$ 分别除以其非零范数，积分后即得 Hölder；范数为零时直接成立，端点 $p=1,q=\infty$ 用几乎处处上界。

对 $1<p<\infty$，先由 $|f+g|^p\leq2^{p-1}(|f|^p+|g|^p)$ 保证可积，再写
\[
 \|f+g\|_p^p
 \leq\int(|f|+|g|)|f+g|^{p-1}
 \leq(\|f\|_p+\|g\|_p)\|f+g\|_p^{p-1}.
\]
除去公共因子得到 Minkowski；若该因子为零则结论直接成立。端点 $p=1,\infty$ 由三角不等式得到。
\end{proof}
简单函数在 $L^p(\mu)$（$p<\infty$）中稠密：先截断函数的幅值与小值部分，使所剩非零支撑有限测度，再将值域划分为有限小区间，最后用支配收敛控制误差。在 $\mathbb R^d$ 的 Lebesgue 测度下，进一步有 $C_c^\infty$ 在 $L^p$ 中稠密；该光滑逼近结论在本书作为实分析结果引用，不对任意抽象测度空间无条件使用。

\section{对偶表示所用结果}
有限有符号或复测度 $\nu$ 的全变差记为 $|\nu|$；对可测集 $E$，$|\nu|(E)$ 是有限可测分割 $E=\bigcup_jE_j$ 上 $\sum_j|\nu(E_j)|$ 的上确界。若 $\mu(E)=0$ 蕴含 $\nu(E)=0$，记 $\nu\ll\mu$，称为绝对连续。

本书引用 Radon--Nikodym 定理的如下形式：若 $\mu$ 为 $\sigma$ 有限正测度、$\nu$ 为全变差有限且 $\nu\ll\mu$ 的有符号或复测度，则存在 $g\in L^1(\mu)$，使 $\nu(E)=\int_Eg\,\mathrm d\mu$，且 $g$ 几乎处处唯一。

它给出 $L^p$ 对偶表示的基本思路：对 $L^p$ 上的有界泛函，在有限测度区域内定义 $\nu(E)=F(\mathbf1_E)$，验证可数可加性与绝对连续性；用 Radon--Nikodym 得到密度，再用 Hölder 型测试证明密度属于 $L^q$ 且范数相等。将可数个有限测度区域拼接，就得到 $\sigma$ 有限情形。这里仅说明表示证明的路线，不代替测度论中的完整论证。

另一项引用结果是 Riesz--Markov 表示：若 $K$ 为紧度量空间，$C(K)$ 上每个连续线性泛函都唯一写为 $F(f)=\int_Kf\,\mathrm d\nu$，其中 $\nu$ 为有限正则 Borel 有符号测度或复测度，且 $\|F\|=|\nu|(K)$。正则性指其全变差可由内侧紧集和外侧开集逼近。它与 Hilbert 空间的 Riesz 表示有不同的表示对象，不可混同。

\section{习题}
\begin{enumerate}
 \item 用计数测度从积分不等式推出数列版本。
 \item 若 $\mu(\Omega)<\infty$、$1\leq p<q<\infty$，证明 $\|f\|_p\leq\mu(\Omega)^{1/p-1/q}\|f\|_q$。
 \item 取 $f_n=n\mathbf1_{(0,1/n)}$，说明几乎处处趋零为何不保证积分趋零。
 \item 用 $L^p$ 的稠密性区分“连续函数可逼近可积函数”与“可积函数本身连续”。
\end{enumerate}

\chapter{拓扑与集合论工具}
\label{app:fa-topology}

度量空间足以处理主线中的大多数例子，但一般弱紧性涉及任意坐标乘积，必须补充少量不依赖度量的语言。本附录区分定义、可直接验证的事实与作为先修知识引用的定理。进一步论证可参阅\cite{teschl-fa}的集合论及拓扑附录。

\section{拓扑、紧性与乘积}
集合 $X$ 上的拓扑 $\tau$ 是一族子集，包含 $\varnothing,X$，对任意并和有限交封闭；其成员称开集。映射连续是指开集的逆像开。任意两个不同点都有不交邻域时，空间称为 Hausdorff 空间。

紧性的定义仍是每个开覆盖有有限子覆盖。紧空间的闭子集紧，连续像紧；Hausdorff 空间中的紧子集闭。最后一个事实可如下证明：对紧集外的点 $x$，分别把它与紧集中每一点用不交邻域分开，再对紧集选有限子覆盖；相应的有限个 $x$ 邻域之交就是与紧集不交的邻域。

乘积 $\prod_{i\in I}X_i$ 的基本开集只限制有限个坐标：在有限指标集 $F\subset I$ 上要求 $x_i\in U_i$，其余坐标任意。本书引用\keyterm{Tychonoff 定理}：任意一族紧空间的乘积在乘积拓扑下紧。它是 Banach--Alaoglu 定理中唯一使用的任意乘积紧性结果。

\section{网与子网}
有向集是带有反身、传递关系的非空集合 $D$，且任意两个元素都有共同后继。\keyterm{网（net）}是以有向集为指标的族 $(x_\alpha)_{\alpha\in D}$；它趋于 $x$，是指每个 $x$ 的邻域最终包含全部网元素。

序列就是以自然数为指标的网。若 $x\in\overline A$，用 $x$ 的邻域按反向包含排序，在每个邻域 $U$ 中选择 $a_U\in U\cap A$，便得到趋于 $x$ 的网。反之，一个来自 $A$ 的网若趋于 $x$，则每个邻域与 $A$ 相交。因此网总能刻画闭包，而序列只有在额外条件下才能做到。

若 $\phi:E\to D$ 保序，且对每个 $\alpha\in D$，最终有 $\phi(\beta)\geq\alpha$，则 $(x_{\phi(\beta)})_{\beta\in E}$ 称为一个子网。一般紧性判别断言：拓扑空间紧当且仅当每个网都有收敛子网。这里引用该判别；它解释为什么一般弱星紧性不能直接改写成序列子列结论。

\section{Zorn 引理}
偏序关系除反身、传递外还满足反对称性。链是其中任意两个元素都可比较的子集；上界是大于等于该子集全部元素的元素；极大元是无法被严格扩大的元素，不要求它大于整个集合的所有元素。

本书采用 Zorn 引理：若非空偏序集的每条链都在该集合内有上界，则它存在极大元。这是与选择公理等价的集合论原则。在 Hahn--Banach 证明中，元素是定义在不同子空间上的受控延拓；链的上界由这些兼容映射之并给出。使用引理时必须验证并后的定义域是线性子空间、数值无冲突且控制不等式仍成立。

\section{半范数族与局部凸空间}
若向量空间的加法与数乘连续，称为拓扑向量空间（topological vector space）；若零点有凸邻域基，称为局部凸空间（locally convex space）。一族半范数 $(p_i)_{i\in I}$ 可用基本邻域
\[
 \{x:p_{i_1}(x)<\varepsilon,\ldots,p_{i_m}(x)<\varepsilon\}
\]
生成局部凸拓扑；若半范数族能分离非零向量，则拓扑为 Hausdorff。弱拓扑就是取 $p_f(x)=|f(x)|$ 得到的情形。

例如，对开集 $\Omega\subset\mathbb R^d$ 中的光滑函数，用 $p_{K,m}(u)=\max_{|\alpha|\leq m}\sup_{x\in K}|\partial^\alpha u(x)|$ 描述各阶导数在紧集 $K\subset\Omega$ 上一致收敛。这里 $\alpha=(\alpha_1,\ldots,\alpha_d)$ 为非负整数多重指标，$|\alpha|=\sum_j\alpha_j$。这种结构用于分布理论，不必强行由单个范数描述。

\section{习题}
\begin{enumerate}
 \item 直接证明坐标投影在乘积拓扑下连续。
 \item 用网的定义证明闭集包含来自该集合的所有收敛网的极限。
 \item 解释 Zorn 引理中的极大元与最大元的区别。
 \item 由半范数邻域验证弱拓扑下的加法和数乘连续。
\end{enumerate}

\chapter{复分析与谱论的补充工具}
\label{app:fa-complex}

本附录只列出谱论直接使用的复分析工具，并说明算子值表达式的含义。复分析中的积分定理与逼近定理作为先修结论使用。

\section{实空间的复化}
若 $X$ 为实 Banach 空间，令 $X_{\mathbb C}=\{x+\mathrm iy:x,y\in X\}$，按通常规则定义复数乘法。一种可用范数是
\[
 \|x+\mathrm iy\|_{\mathbb C}=
 \sup_{\theta\in\mathbb R}\|x\cos\theta-y\sin\theta\|.
\]
旋转参数给出复齐次性，三角不等式逐个 $\theta$ 验证；并有
\[
 \max(\|x\|,\|y\|)\leq\|x+\mathrm iy\|_{\mathbb C}\leq\|x\|+\|y\|.
\]
所以它完备，且实空间等距嵌入。实有界算子 $T$ 延拓为 $T_{\mathbb C}(x+\mathrm iy)=Tx+\mathrm iTy$，其范数仍为 $\|T\|$。实算子的复谱定义为 $\sigma(T_{\mathbb C})$；实平面九十度旋转没有实特征值，却有复特征值 $\mathrm i,-\mathrm i$。

\section{解析函数的基本结论}
复函数在开集内每一点复可微时称解析（holomorphic）。本书使用 Cauchy 积分公式及其估计：若 $f$ 在包含闭圆盘 $|z-z_0|\leq R$ 的邻域内解析，则
\[
 f^{(n)}(z_0)=\frac{n!}{2\pi\mathrm i}
 \int_{|z-z_0|=R}\frac{f(z)}{(z-z_0)^{n+1}}\,\mathrm dz,
 \qquad |f^{(n)}(z_0)|\leq\frac{n!}{R^n}\max_{|z-z_0|=R}|f(z)|.
\]
若 $f$ 是有界整函数，让 $R\to\infty$ 即得 $f'=0$，所以 $f$ 为常数；这就是 Liouville 定理。若一条闭曲线在解析区域内连续变形，积分值不变。谱半径公式的证明使用了这一性质。

连续函数演算还使用 Weierstrass 逼近定理：实紧区间上的连续函数可由多项式一致逼近；对实紧子集先作连续延拓即可。复平面紧集上的一般连续函数则需使用含 $z,\overline z$ 的多项式及 Stone--Weierstrass 定理，不能不加条件地只用 $z$ 的多项式。

\section{算子值积分与标量化}
若 $F:[a,b]\to\mathcal B(X)$ 按算子范数连续，则其 Riemann 和在完备空间 $\mathcal B(X)$ 中收敛，定义积分，并满足
\[
 \left\|\int_a^bF(t)\,\mathrm dt\right\|\leq\int_a^b\|F(t)\|\,\mathrm dt.
\]
分段光滑曲线上的积分由参数化定义。每个 $x\in X$ 和 $f\in X^*$ 都可移入积分，即 $f((\int F)x)=\int f(Fx)$。因此可以先证明标量积分恒等式，再利用连续对偶分离向量恢复算子恒等式。

预解算子的 Neumann 展开按算子范数在足够小的闭圆盘上一致收敛，所以可逐项积分。在大圆上提取 $R(\lambda,T)=\sum_{k\geq0}T^k\lambda^{-k-1}$ 的第 $n$ 项，就得到谱半径证明中的积分表达。这里的一致收敛是交换求和与积分的依据。

\section{习题}
\begin{enumerate}
 \item 验证上述复化范数的复齐次性及完备性。
 \item 从预解恒等式推出 $\frac{\mathrm d}{\mathrm d\lambda}R(\lambda,T)=-R(\lambda,T)^2$。
 \item 在有限维对角矩阵上直接核验谱半径证明中的轮廓积分公式。
 \item 说明算子值函数弱意义上的逐点信息为何不足以无条件交换极限与算子范数。
\end{enumerate}

\chapter{部分习题解答与证明方法回顾}
\label{app:fa-solutions}

本附录选取能够代表主要证明方法的习题，给出完整论证。阅读时宜先独立确定所需估计，再与解答比较；只记最后的结论，难以掌握这些方法在新问题中的用法。其余习题可沿相邻例题与定理的证明尝试。

\section{由局部估计控制整个序列}
第\ref{chap:fa-metric}章习题\ref{exer:fa-cauchy}要求证明 Cauchy 序列有界，并说明一个收敛子列足以决定整个 Cauchy 序列的极限。取 $\varepsilon=1$，存在 $N$ 使 $m,n\geq N$ 时 $d(x_m,x_n)<1$。固定 $m=N$，后面所有项都落在 $B(x_N,1)$ 中；前面只有有限项，取
\[
 R=1+\max_{1\leq j\leq N}d(x_j,x_N),
\]
即可把整个序列包含在 $B(x_N,R)$ 中。这一步使用的是“尾部统一受控，加上有限个例外”的结构。

若 $x_{n_k}\to x$，给定 $\varepsilon>0$，先由 Cauchy 条件选择 $N$，使后部两项距离小于 $\varepsilon/2$，再选 $k$ 使 $n_k\geq N$ 且 $d(x_{n_k},x)<\varepsilon/2$。对所有 $n\geq N$，三角不等式给出 $d(x_n,x)<\varepsilon$。注意这里不是给每个 $n$ 临时寻找一个不同的近点，而是用同一个子列元素作为整个尾部的比较点。这个论证经常用于完备性证明：先为一个精心选择的子列找到极限，再回到原序列。

\section{数列范数的包含关系}
第\ref{chap:fa-banach}章习题\ref{exer:fa-lp-inclusion}可以先通过归一化消除大小因素。若 $x=0$，结论直接成立；否则令 $z=x/\|x\|_p$，则 $\sum_j|z_j|^p=1$，所以每个 $|z_j|\leq1$。当 $q<\infty$ 时，$|z_j|^q\leq|z_j|^p$，故 $\sum_j|z_j|^q\leq1$，从而 $\|x\|_q\leq\|x\|_p$。当 $q=\infty$ 时，逐坐标估计直接给出同一结论。

当 $p<q<\infty$，选择 $1/q<a<1/p$，则 $(j^{-a})\in\ell^q$ 而不在 $\ell^p$，说明包含通常严格。$q=\infty$ 时，常数列已经给出严格性。这里的证明依赖计数测度的每个单点都有固定正质量，因此一个坐标不能在 $\ell^p$ 范数有界时任意变大；有限区间上的 $L^p$ 函数却可在很小的集合上取很大的值，这正是两类包含方向不同的原因。

\section{伴随与正交补的计算}
第\ref{chap:fa-hilbert}章习题\ref{exer:fa-shift-adjoint}中，Cauchy--Schwarz 保证下列级数绝对收敛，因此移换指标合法：
\[
 \langle Rx,y\rangle=\sum_{j\geq2}x_{j-1}\overline{y_j}
 =\sum_{k\geq1}x_k\overline{y_{k+1}}=\langle x,Ly\rangle.
\]
由伴随唯一性，$R^*=L$、$L^*=R$。又 $LR=I$，而 $RLx=x-x_1e_1$，所以右移保持范数，却不满射，且 $RR^*\ne I$，不是酉算子。等距算子只要求 $T^*T=I$；酉算子还要求能够从每个输出倒推输入，也就是在此基础上满射。

习题\ref{exer:fa-double-perp}则适合使用正交分解。首先 $M^\perp=(\overline M)^\perp$，因为内积连续；每个 $m\in\overline M$ 与 $M^\perp$ 正交，所以 $\overline M\subset M^{\perp\perp}$。反向对 $x\in M^{\perp\perp}$，分解 $x=m+z$，其中 $m\in\overline M$、$z\in M^\perp$。由于 $x$ 与 $z$ 正交，而 $m$ 也与 $z$ 正交，便有 $0=\langle x,z\rangle=\|z\|^2$，所以 $z=0$、$x=m\in\overline M$。两次正交补会自动补齐极限点，因此结果是闭包，未必是原来的非闭子空间。

\section{Baire 方法与稠密集上的验证}
第\ref{chap:fa-baire}章习题\ref{exer:fa-hamel}中，反设 $e_1,e_2,\ldots$ 为无限维 Banach 空间 $X$ 的可数 Hamel 基。令 $M_n=\operatorname{span}\{e_1,\ldots,e_n\}$，每个 $M_n$ 有限维，因而闭。真线性子空间内部必为空：若 $B(x_0,r)\subset M_n$，相减得到 $B(0,r)\subset M_n$，再由数乘可把每个向量缩入该球，迫使 $M_n=X$，矛盾。于是每个 $M_n$ 无处稠密。但 Hamel 展开仅有有限项，所以 $X=\bigcup_nM_n$，违反 Baire 定理。可数标准正交基则允许无穷级数，不给出这个可数并等式，因此不会产生矛盾。

习题\ref{exer:fa-dense-convergence}设 $D\subset X$ 稠密，$Y$ 完备，且 $\|T_n\|\leq M$。给定 $x\in X$ 与 $z\in D$，
\[
 \|T_nx-T_mx\|\leq2M\|x-z\|+\|T_nz-T_mz\|.
\]
若 $M=0$，所有算子都为零；否则先取 $z$ 使第一项小于 $\varepsilon/2$，再由 $D$ 上的收敛选择 $n,m$ 使第二项小于 $\varepsilon/2$。于是 $(T_nx)$ 在 $Y$ 中 Cauchy，故收敛。逐点极限 $T$ 线性且 $\|Tx\|\leq M\|x\|$。这项结论只需要值域完备；若统一算子界已经给出，定义域本身无需完备。统一估计负责把稠密集中的结论传到全空间，完备性负责接住输出的极限。

\section{弱紧性怎样保证最值}
第\ref{chap:fa-reflexive}章习题\ref{exer:fa-norm-attainment}中，若 $f=0$，任取 $x=0$ 即可。若 $f\ne0$，按上确界定义选 $x_n\in B_X$，使 $|f(x_n)|\to\|f\|$。自反性给出弱收敛子列 $x_{n_k}\rightharpoonup x$；范数弱下半连续，故 $\|x\|\leq1$。连续线性泛函按弱收敛的定义满足 $f(x_{n_k})\to f(x)$，取绝对值得 $|f(x)|=\|f\|$。由于 $|f(x)|\leq\|f\|\|x\|$，进一步得到 $\|x\|=1$。

这个论证与紧集上连续函数取得最大值的证明相同，只是把范数子列换成弱收敛子列。不能把结论推广到单位球上的任意范数连续函数，因为范数连续不保证弱连续；例如无限维 Hilbert 空间中的范数函数在 $e_n\rightharpoonup0$ 时不保持函数值极限。

\section{Galerkin 收敛中嵌套性的作用}
第\ref{chap:fa-variational}章习题\ref{exer:fa-galerkin}中，已知 $H_n\subset H_{n+1}$ 且 $\bigcup_nH_n$ 稠密。给定 $\varepsilon>0$，稠密性保证存在某个 $N$ 与 $w\in H_N$，使 $\|u-w\|<\varepsilon$；嵌套性保证同一个 $w$ 属于所有 $n\geq N$ 的 $H_n$。因此这些 $n$ 都满足
\[
 \inf_{w_n\in H_n}\|u-w_n\|\leq\|u-w\|<\varepsilon.
\]
Céa 估计随即给出 $u_n\to u$。这里先找到一个好近似，再保证它在以后所有逼近空间中仍可使用，正是嵌套性的作用。

仅有并集稠密却未必足够。在 $\ell^2$ 中，可令 $H_{2n}=\operatorname{span}\{e_1,\ldots,e_n\}$，$H_{2n-1}=\operatorname{span}\{e_1\}$。它们的并集稠密，但对于 $u=e_2$，奇数步的最佳逼近误差始终为一。若取 $a(u,v)=\langle u,v\rangle$ 的实 Hilbert 空间问题，Galerkin 解就是正交投影，因此这一序列的离散解也不收敛。非嵌套方法并非不可用，但必须直接验证对整个指标序列成立的逼近条件。

\begin{thebibliography}{9}
\bibitem{mit-syllabus}
Massachusetts Institute of Technology.
\emph{18.102 Introduction to Functional Analysis, Spring 2021: Syllabus}.
\url{https://ocw.mit.edu/courses/18-102-introduction-to-functional-analysis-spring-2021/pages/syllabus/}.

\bibitem{mit-notes}
Massachusetts Institute of Technology.
\emph{18.102 Introduction to Functional Analysis, Spring 2021: Lecture Notes and Readings}.
\url{https://ocw.mit.edu/courses/18-102-introduction-to-functional-analysis-spring-2021/pages/lecture-notes-and-readings/}.
\bibitem{teschl}
Gerald Teschl.
\emph{Mathematical Methods in Quantum Mechanics: With Applications to Schrödinger Operators}.
作者公开版本，重点参阅自伴算子与谱定理部分。
\url{https://www.mat.univie.ac.at/~gerald/ftp/book-schroe/schroe.pdf}.

\bibitem{engel-nagel}
Klaus-Jochen Engel and Rainer Nagel.
\emph{One-Parameter Semigroups for Linear Evolution Equations}.
第二章：生成元、预解算子与生成定理。
\url{https://www.math.uni-tuebingen.de/de/forschung/agfa/members/engel-nagel_one-parameter-semigroups.pdf}.
\bibitem{teschl-fa}
Gerald Teschl.
\emph{Topics in Linear and Nonlinear Functional Analysis}.
作者书目与内容目录页。
\url{https://www.mat.univie.ac.at/~gerald/ftp/book-fa/index.html}.

\bibitem{teschl-ra}
Gerald Teschl.
\emph{Topics in Real Analysis}.
作者书目与内容目录页。
\url{https://www.mat.univie.ac.at/~gerald/ftp/book-ra/index.html}.
\end{thebibliography}

\end{document}
